% =====================================================================
%  รายงานโครงงานวิศวกรรมคอมพิวเตอร์ ENG23 4080
%  CPE-691008 — หุ่นยนต์เคลื่อนที่อัตโนมัติสำหรับการสำรวจพื้นที่ฯ
%  จัดรูปแบบตาม Template 04_Project_ENG23_4080
%  ฉบับนี้: เพิ่มตำแหน่งรูปภาพ (placeholder) ด้วยคำสั่ง \imgph
%           ดูวิธีแทนที่ด้วยรูปจริงได้ที่นิยามของ \imgph ในส่วน preamble
%
%  *** ต้องคอมไพล์ด้วย XeLaTeX เท่านั้น ***
%  ไฟล์ที่ต้องอัปโหลดไว้โฟลเดอร์เดียวกับ .tex :
%     THSarabun.ttf / THSarabun Bold.ttf
%     THSarabun Italic.ttf / THSarabun BoldItalic.ttf
%     (ชื่อไฟล์ต้องตรงกับชื่อใน \setmainfont ทุกตัวอักษร รวมช่องว่างด้วย
%      ถ้าไฟล์ชื่อ THSarabunNew*.ttf ให้แก้ชื่อในบล็อกฟอนต์เป็น THSarabunNew)
%     sut_logo.png , system_architecture.jpg
% =====================================================================
\documentclass[a4paper,12pt]{report}

% ─── ขนาดหน้ากระดาษ (ตรงตาม Template: ซ้าย 1.5 นิ้ว ที่เหลือ 1 นิ้ว) ──
\usepackage{geometry}
\geometry{
    a4paper,
    left   = 3.81cm,
    right  = 2.54cm,
    top    = 2.54cm,
    bottom = 2.54cm,
    footskip = 1.25cm
}

\usepackage{fontspec}
\usepackage{titlesec}
\usepackage{titletoc}
\usepackage{enumitem}
\usepackage{tabularx}
\usepackage{array}
\usepackage{multirow}
\usepackage{graphicx}
\usepackage{adjustbox}  % กันรูป TikZ ใหญ่เกินกรอบหน้า
\graphicspath{{images/}{./}}  % รูปอยู่โฟลเดอร์เดียวกับ .tex หรือในโฟลเดอร์ images/
\usepackage{tikz}
\usetikzlibrary{shapes.geometric, shapes.multipart, arrows.meta, positioning, calc}
\usepackage{amsmath}
\usepackage{amssymb}
\usepackage{anyfontsize}  % กันคำเตือนเรื่องฟอนต์คณิตศาสตร์ตอนใช้ขนาดแปลก ๆ
\usepackage[labelsep=space]{caption}
\usepackage{tocloft}
\usepackage{indentfirst}
\usepackage{fancyhdr}
\usepackage[hidelinks]{hyperref}

% ─── ฟอนต์ TH Sarabun ────────────────────────────────────────────────
%  *** ชื่อฟอนต์ต้องตรงกับชื่อไฟล์ .ttf ที่วางไว้โฟลเดอร์เดียวกับ .tex ***
%  ชื่อที่พบบ่อยมีสามแบบ : THSarabun / THSarabunNew / THSarabunPSK
%  Template และรายงานรุ่นพี่ฝังฟอนต์ชื่อ THSarabunPSK
%  ถ้าเปลี่ยนชื่อ ต้องเปลี่ยนให้ครบทั้งสามบรรทัดด้านล่าง
%
%  *** เหตุผลที่ต้องมี \setsansfont และ \setmonofont ***
%  ถ้าประกาศแต่ \setmainfont ตัวอักษรอังกฤษทุกตัวที่อยู่ใน \texttt{...}
%  (ในไฟล์นี้มี 55 จุด เช่น ชื่อโหนด ชื่อเฟรม ชื่อไฟล์) จะถูกเรนเดอร์ด้วย
%  Latin Modern Mono ซึ่งเป็นฟอนต์พิมพ์ดีดคนละแบบกับ TH Sarabun
%  และ \textsf{...} จะถูกเรนเดอร์ด้วย Latin Modern Sans
%  นี่คือสาเหตุที่ตัวอักษรอังกฤษบางจุดหน้าตาไม่เหมือนที่อื่นในเล่ม
%
%  หมายเหตุ: ห้ามใส่ Scale ที่นี่ ขนาดตัวอักษรคุมด้วย package fontsize แทน
\defaultfontfeatures{
    Extension      = .ttf,
    UprightFont    = *,
    BoldFont       = * Bold,
    ItalicFont     = * Italic,
    BoldItalicFont = * BoldItalic
}

\setmainfont{THSarabun}[Ligatures = TeX]
\setsansfont{THSarabun}[Ligatures = TeX]
\setmonofont{THSarabun}
%  ^ ไม่ใส่ Ligatures = TeX ให้ฟอนต์พิมพ์ดีด เพื่อไม่ให้เครื่องหมาย --
%    ในชื่อไฟล์หรือโค้ดถูกรวมเป็นขีดยาวโดยอัตโนมัติ
%
%  ถ้าอยากให้โค้ดใน \texttt{...} เป็นฟอนต์พิมพ์ดีดแยกจากเนื้อความจริง ๆ
%  (รายงานรุ่นพี่ใช้ Courier New สำหรับส่วนนี้) ให้ลบบรรทัด \setmonofont ออก

% ─── ขนาดตัวอักษรพื้นฐานของทั้งเล่ม = 16pt (ตาม Template) ────────────
\usepackage[fontsize=16pt]{fontsize}

% ─── คำสั่งช่วย: สั่งขนาดตัวอักษรเป็น pt ตรง ๆ ───────────────────────
%   \ptsize{20}  → ตัวอักษร 20pt (ระยะบรรทัด 1.3 เท่าให้อัตโนมัติ)
\newcommand{\ptsize}[1]{\fontsize{#1pt}{\fpeval{#1*1.3}pt}\selectfont}

% การตัดคำภาษาไทย (XeLaTeX เท่านั้น)
\XeTeXlinebreaklocale "th"
\XeTeXlinebreakskip = 0pt plus 1pt

% ─── ย่อหน้าและระยะบรรทัด ────────────────────────────────────────────
\setlength{\parindent}{1.25cm}
\setlength{\parskip}{0pt}
\linespread{1.0}

% ─── แปลคำสั่งของ report class เป็นภาษาไทย ───────────────────────────
\renewcommand{\chaptername}{บทที่}
\renewcommand{\contentsname}{สารบัญ}
\renewcommand{\listtablename}{สารบัญตาราง}
\renewcommand{\listfigurename}{สารบัญรูป}
\renewcommand{\bibname}{เอกสารอ้างอิง}
\renewcommand{\appendixname}{ภาคผนวก}
\renewcommand{\figurename}{รูปที่}
\renewcommand{\tablename}{ตารางที่}
\captionsetup{labelsep=space, justification=centering, skip=6pt}

% ─── เลขหน้าไทย ก ข ค ──────────────────────────────────────────────
\makeatletter
\def\@thaialph#1{%
  \ifcase#1\or ก\or ข\or ค\or ง\or จ\or ฉ\or ช\or ซ\or ฌ\or ญ%
  \or ฎ\or ฏ\or ฐ\or ฑ\or ฒ\or ณ\or ด\or ต\or ถ\or ท%
  \or ธ\or น\or บ\or ป\or ผ\or ฝ\or พ\or ฟ\or ภ\or ม%
  \or ย\or ร\or ล\or ว\or ศ\or ษ\or ส\or ห\or ฬ\or อ\or ฮ%
  \else\@ctrerr\fi}
\newcommand{\thaialph}[1]{\expandafter\@thaialph\csname c@#1\endcsname}
\makeatother

% ตัวนับสำหรับ "จำเลขหน้าไทยหน้าสุดท้ายของส่วนนำ"
\newcounter{savedfront}
\newcommand{\savefrontpage}{\setcounter{savedfront}{\value{page}}}
\newcommand{\resumefrontpage}{%
  \clearpage
  \pagenumbering{thaialph}\setcounter{page}{\value{savedfront}}}

% ─── เลขหน้าท้ายกระดาษชิดขวา (ตรงตาม Template) ───────────────────────
\pagestyle{fancy}
\fancyhf{}
\fancyfoot[R]{\thepage}
\renewcommand{\headrulewidth}{0pt}
\fancypagestyle{plain}{%
  \fancyhf{}%
  \fancyfoot[R]{\thepage}%
  \renewcommand{\headrulewidth}{0pt}%
}

% ─── รูปแบบหัวข้อบท : "บทที่ N" บรรทัดบน / ชื่อบทบรรทัดล่าง กึ่งกลาง ──
\titleformat{\chapter}[display]
  {\normalfont\bfseries\centering}
  {\ptsize{18}\chaptertitlename\ \thechapter}{12pt}{\ptsize{18}}
\titlespacing*{\chapter}{0pt}{0pt}{24pt}

% ขนาดหัวข้อตาม Template: ชื่อบท 18pt ตัวหนา, หัวข้อย่อย 16pt ตัวหนา
\titleformat{name=\chapter,numberless}[display]
  {\normalfont\bfseries\centering}{}{0pt}{\ptsize{18}}
\titlespacing*{name=\chapter,numberless}{0pt}{0pt}{24pt}

\titleformat{\section}{\normalfont\normalsize\bfseries}{\thesection}{1em}{}
\titlespacing*{\section}{0pt}{12pt}{6pt}
\titleformat{\subsection}{\normalfont\normalsize\bfseries}{\thesubsection}{1em}{}
\titlespacing*{\subsection}{0pt}{10pt}{4pt}

% ─── รูปแบบสารบัญ : "บทที่ N   ชื่อบท" + จุดไข่ปลา ────────────────────
\renewcommand{\cftchapfont}{\normalfont\bfseries}
\renewcommand{\cftchappagefont}{\normalfont\bfseries}
\renewcommand{\cftchapleader}{\cftdotfill{\cftdotsep}}
\renewcommand{\cftchappresnum}{บทที่ }
\renewcommand{\cftchapaftersnum}{\hspace{1em}}
\setlength{\cftchapnumwidth}{3.2cm}
\setlength{\cftsecindent}{3.2cm}
\setlength{\cftsecnumwidth}{1.6cm}
\setlength{\cftbeforechapskip}{6pt}

\renewcommand{\cfttabpresnum}{ตารางที่ }
\renewcommand{\cfttabaftersnum}{\hspace{1em}}
\setlength{\cfttabnumwidth}{3.2cm}
\renewcommand{\cftfigpresnum}{รูปที่ }
\renewcommand{\cftfigaftersnum}{\hspace{1em}}
\setlength{\cftfignumwidth}{3.2cm}

\renewcommand{\cfttoctitlefont}{\hfill\normalfont\ptsize{18}\bfseries}
\renewcommand{\cftaftertoctitle}{\hfill}
\renewcommand{\cftlottitlefont}{\hfill\normalfont\ptsize{18}\bfseries}
\renewcommand{\cftafterlottitle}{\hfill}
\renewcommand{\cftloftitlefont}{\hfill\normalfont\ptsize{18}\bfseries}
\renewcommand{\cftafterloftitle}{\hfill}
\setlength{\cftbeforetoctitleskip}{0pt}
\setlength{\cftbeforelottitleskip}{0pt}
\setlength{\cftbeforeloftitleskip}{0pt}
\setlength{\cftaftertoctitleskip}{18pt}
\setlength{\cftafterlottitleskip}{18pt}
\setlength{\cftafterloftitleskip}{18pt}

% คำว่า "หน้า" มุมขวาเหนือรายการในสารบัญทั้งสาม
\newcommand{\tocpagehead}{\hfill\textbf{หน้า}\par\vspace{6pt}}


% ─── กล่องพื้นที่สำหรับรูปภาพ (Placeholder) ──────────────────────────
%  ใช้ชั่วคราวแทนรูปจริง เมื่อได้ไฟล์รูปแล้วให้แทนที่บรรทัด \imgph{...}
%  ทั้งบรรทัดด้วย  \includegraphics[width=0.85\textwidth]{ชื่อไฟล์}
%  รูปแบบ: \imgph[ความกว้าง]{ความสูง}{ชื่อไฟล์ที่จะใส่}{สิ่งที่ต้องถ่าย/จับภาพ}
\newcommand{\imgph}[4][\textwidth]{%
  \setlength{\fboxsep}{6pt}%
  \fbox{\begin{minipage}[c][#2][c]{\dimexpr#1-2\fboxsep-2\fboxrule\relax}
    \centering\small
    [ พื้นที่สำหรับรูปภาพ ]\\[2pt]
    \texttt{\detokenize{#3}}\\[6pt]
    \rule{0.85\linewidth}{0.4pt}\\[6pt]
    #4
  \end{minipage}}%
}

\begin{document}

% =====================================================================
%  ใบปะหน้า
% =====================================================================
\begin{titlepage}
\ptsize{20}
\centering
\includegraphics[height=2.6cm]{sut_logo.png}
\vspace{\stretch{1}}

{\bfseries หุ่นยนต์เคลื่อนที่อัตโนมัติสำหรับการสำรวจพื้นที่และระบุตำแหน่งวัตถุ\\[4pt]
โดยใช้ไลดาร์และเอดจ์เอไอ\par}
\vspace{\stretch{0.6}}
{\bfseries Autonomous Mobile Robot for Area Exploration and Object\\[4pt]
Localization Using LiDAR and Edge AI\par}
\vspace{\stretch{2}}

{\bfseries โดย}\\[0.5cm]
\begin{tabular}{@{}l@{\hspace{1em}}l@{}}
นายชทัตพล เสริมศรี      & รหัสนักศึกษา B6603908\\[3pt]
นายสุรเกียรติ สิงขรอาสน์  & รหัสนักศึกษา B6603946\\[3pt]
นายกิตตินันท์ ปัจจัยโคถา  & รหัสนักศึกษา B6618643\\[3pt]
นายทองนรินทร์ แย้มศรี    & รหัสนักศึกษา B6627713\\
\end{tabular}
\vspace{\stretch{2}}

{\bfseries อาจารย์ที่ปรึกษา}\\[0.5cm]
อาจารย์ ดร.วิชัย ศรีสุรักษ์
\vspace{\stretch{2.5}}

{\bfseries รายงานนี้เป็นส่วนหนึ่งของรายวิชา ENG23 4080 โครงงานวิศวกรรมคอมพิวเตอร์}\\[4pt]
{\bfseries หลักสูตรวิศวกรรมศาสตรบัณฑิต สาขาวิชาวิศวกรรมคอมพิวเตอร์}\\[4pt]
{\bfseries สำนักวิชาวิศวกรรมศาสตร์ มหาวิทยาลัยเทคโนโลยีสุรนารี}\\[4pt]
{\bfseries ประจำภาคการศึกษาที่ 1 ปีการศึกษา 2569}
\end{titlepage}


% =====================================================================
%  ส่วนนำ — เลขหน้า ก ข ค
% =====================================================================
\pagenumbering{thaialph}

\addtocontents{toc}{\protect\tocpagehead}
\addtocontents{lot}{\protect\tocpagehead}
\addtocontents{lof}{\protect\tocpagehead}

% ------------------- บทคัดย่อ -------------------
%  ตาม Template: ข้อมูลโครงงาน (ใบรองปก) อยู่ส่วนบนของหน้าบทคัดย่อ หน้า ก
\clearpage
\phantomsection
\addcontentsline{toc}{chapter}{บทคัดย่อ}
{\noindent
\newcommand{\lbl}[1]{\textbf{#1}}
\begin{tabular}{@{}>{\raggedright\arraybackslash}p{4cm}@{\hspace{4mm}}p{\dimexpr\textwidth-4.4cm\relax}@{}}
\lbl{ชื่อโครงงาน} &
หุ่นยนต์เคลื่อนที่อัตโนมัติสำหรับการสำรวจพื้นที่และระบุตำแหน่งวัตถุ\\[2pt]
 & โดยใช้ไลดาร์และเอดจ์เอไอ\\[2pt]
 & (Autonomous Mobile Robot for Area Exploration and\\[2pt]
 & Object Localization Using LiDAR and Edge AI)\\[16pt]

\lbl{ผู้จัดทำ} &
\begin{tabular}[t]{@{}l@{\hspace{1em}}l@{}}
นายชทัตพล เสริมศรี      & รหัสนักศึกษา B6603908\\[2pt]
นายสุรเกียรติ สิงขรอาสน์  & รหัสนักศึกษา B6603946\\[2pt]
นายกิตตินันท์ ปัจจัยโคถา  & รหัสนักศึกษา B6618643\\[2pt]
นายทองนรินทร์ แย้มศรี    & รหัสนักศึกษา B6627713
\end{tabular}\\[16pt]

\lbl{อาจารย์ที่ปรึกษา} & อาจารย์ ดร.วิชัย ศรีสุรักษ์\\[16pt]

\lbl{ภาคการศึกษา} & 1/2569\\
\end{tabular}
}

\vspace{24pt}
{\centering\ptsize{18}\bfseries บทคัดย่อ\par}
\vspace{12pt}

โครงงานนี้พัฒนาหุ่นยนต์เคลื่อนที่อัตโนมัติแบบขับเคลื่อนสองล้ออิสระ (Differential Drive)
สำหรับสำรวจพื้นที่ปิดภายในอาคาร สร้างแผนที่ และระบุตำแหน่งของวัตถุที่ตรวจพบลงบนแผนที่
โดยประมวลผลทั้งหมดภายในตัวหุ่นยนต์
ระบบทำงานบน ROS 2 Humble และหน่วยประมวลผล NVIDIA Jetson Orin NX (Advantech MIC-711-ON)
ใช้ไลดาร์สองมิติร่วมกับ slam\_toolbox สร้างแผนที่และระบุตำแหน่งของหุ่นยนต์
ใช้ Nav2 นำทางและสำรวจพื้นที่อัตโนมัติ
ตรวจจับวัตถุจากกล้องด้วยแบบจำลอง YOLO26l บน TensorRT
แล้วผสานทิศทางของวัตถุในภาพกับระยะจากไลดาร์เพื่อระบุพิกัดบนแผนที่
และแสดงผลแบบเรียลไทม์ผ่านเว็บแอปพลิเคชัน

ผลการทดสอบบนหุ่นยนต์จริงผ่านเกณฑ์ที่กำหนดทั้งห้าข้อ
ความคลาดเคลื่อนเชิงตำแหน่งของ SLAM สูงสุด 6.4 เซนติเมตร (เกณฑ์ 10 เซนติเมตร)
ความแม่นยำของการตรวจจับวัดเป็นค่า F1 เฉลี่ยได้ร้อยละ 82.2 (เกณฑ์ร้อยละ 80)
เพิ่มขึ้นจากร้อยละ 64.4 ของแบบจำลองรุ่นแรก
การตรวจจับทำงานได้ 15.16 เฟรมต่อวินาทีขณะ SLAM ทำงานอยู่ (เกณฑ์ 5 เฟรมต่อวินาที)
ตำแหน่งวัตถุบนแผนที่คลาดเคลื่อนเฉลี่ย 0.19 เมตร และสูงสุด 0.34 เมตรเมื่อวัดจากสี่ทิศทาง
(เกณฑ์ 0.50 เมตร)
และอัตราการใช้หน่วยประมวลผลกลางเฉลี่ยร้อยละ 55.8 (เกณฑ์ร้อยละ 80)
หลังแก้ไขการตั้งค่าเธรดของไลบรารี OpenBLAS ที่ทำให้เดิมสูงถึงร้อยละ 81.5
ผลดังกล่าวแสดงว่าอุปกรณ์เอดจ์เอไอเพียงชุดเดียวสามารถทำ SLAM ตรวจจับวัตถุ
และระบุตำแหน่งวัตถุบนแผนที่ไปพร้อมกันได้แบบเรียลไทม์

% ------------------- กิตติกรรมประกาศ -------------------
\chapter*{กิตติกรรมประกาศ}
\addcontentsline{toc}{chapter}{กิตติกรรมประกาศ}

โครงงานวิศวกรรมนี้ได้บรรลุตามวัตถุประสงค์ของผู้จัดทำโครงงานไปด้วยดี
คณะผู้จัดทำขอขอบพระคุณ อาจารย์ ดร.วิชัย ศรีสุรักษ์ อาจารย์ที่ปรึกษาโครงงาน
ที่ให้คำปรึกษา คำแนะนำ และข้อคิดเห็นที่เป็นประโยชน์อย่างยิ่งตลอดระยะเวลาการดำเนินโครงงาน

ขอขอบพระคุณคณาจารย์สาขาวิชาวิศวกรรมคอมพิวเตอร์ สำนักวิชาวิศวกรรมศาสตร์
มหาวิทยาลัยเทคโนโลยีสุรนารี ที่ประสิทธิ์ประสาทวิชาความรู้
และขอขอบคุณผู้ที่มีส่วนเกี่ยวข้องทุกท่านที่ให้ความช่วยเหลือจนโครงงานนี้สำเร็จลุล่วง

\vspace{1cm}
\begin{flushright}
คณะผู้จัดทำ
\end{flushright}

% ------------------- สารบัญ -------------------
\clearpage
\tableofcontents
\clearpage
\listoftables
\clearpage
\listoffigures

\clearpage
\savefrontpage

% =====================================================================
%  เนื้อหา — เลขหน้า 1 2 3
% =====================================================================
\pagenumbering{arabic}

% ---------------------------------------------------------------------
\chapter{บทนำ}

\section{ที่มาและความสำคัญของโครงงาน}

ในการบริหารจัดการพื้นที่ปิดขนาดใหญ่ เช่น คลังสินค้า โกดัง หรืออาคารสำนักงาน
การตรวจสอบและสำรวจพื้นที่เพื่อระบุตำแหน่งของวัตถุภายในอาคารยังคงต้องอาศัยการเดินสำรวจ
ด้วยแรงงานคนเป็นหลัก ซึ่งเป็นงานที่ใช้เวลามาก ทำซ้ำบ่อย และให้ข้อมูลตำแหน่งที่ขาดความแม่นยำ
โดยเฉพาะเมื่อพื้นที่มีขนาดใหญ่หรือมีวัตถุจำนวนมากที่ต้องติดตาม

สาเหตุหลักมาจากการขาดระบบที่สามารถสำรวจพื้นที่และบันทึกตำแหน่งวัตถุได้โดยอัตโนมัติ
การพึ่งพาคนเดินสำรวจทำให้ข้อมูลที่ได้ขึ้นอยู่กับการสังเกตและการจดบันทึกด้วยมือ
ซึ่งเกิดความคลาดเคลื่อนได้ง่าย
อีกทั้งยังไม่มีการเชื่อมโยงตำแหน่งของวัตถุที่พบเข้ากับแผนที่ของพื้นที่อย่างเป็นระบบ
ทำให้ยากต่อการนำข้อมูลไปใช้อ้างอิงในภายหลัง

ผลที่ตามมาคือความล่าช้าในการตรวจสอบพื้นที่ ต้นทุนด้านแรงงานที่สูงอย่างต่อเนื่อง
และความผิดพลาดในการระบุพิกัดตำแหน่งของวัตถุ
ซึ่งส่งผลโดยตรงต่อประสิทธิภาพในการบริหารจัดการพื้นที่
และทำให้ผู้ดูแลไม่สามารถติดตามสภาพภายในอาคารได้อย่างทันท่วงที

แม้ในปัจจุบันจะมีการนำหุ่นยนต์เข้ามาช่วยในงานอัตโนมัติ
แต่หุ่นยนต์จำนวนมากยังทำงานในลักษณะ ``Input-Process-Output''
ตามลำดับขั้นที่กำหนดไว้ล่วงหน้า
และต้องอาศัยการวางโครงสร้างพื้นฐานที่ซับซ้อนและมีต้นทุนสูงเพื่อให้ทำงานได้ถูกต้อง
ช่องว่างสำคัญที่พบคือ การขาดการรวมความสามารถในการระบุตำแหน่งของหุ่นยนต์ (Localization)
เข้ากับการตรวจจับและระบุตำแหน่งของวัตถุไว้ในระบบเดียว
ที่ประมวลผลได้ในตัวอุปกรณ์โดยไม่ต้องพึ่งพาการเชื่อมต่อภายนอกตลอดเวลา

โครงงานนี้จึงเสนอการพัฒนาหุ่นยนต์เคลื่อนที่อัตโนมัติที่ติดตั้งระบบ SLAM
(Simultaneous Localization and Mapping)
เพื่อให้หุ่นยนต์สร้างแผนที่ของพื้นที่และระบุตำแหน่งของตนเองได้แบบเรียลไทม์
ควบคู่กับการติดตั้งเอดจ์เอไอ (Edge AI) เพื่อทำหน้าที่ตรวจจับและจำแนกวัตถุจากกล้อง
แล้วผสานข้อมูลทั้งสองส่วนเข้าด้วยกันเพื่อระบุพิกัดของวัตถุที่ตรวจพบลงบนแผนที่
โดยประมวลผลทั้งหมดภายในตัวอุปกรณ์

ระบบนี้ช่วยลดเวลาและแรงงานในการเดินสำรวจพื้นที่ปิด
เพิ่มความแม่นยำในการบันทึกตำแหน่งวัตถุลงบนแผนที่
และด้วยการประมวลผลแบบ Edge AI ที่จบในตัวอุปกรณ์
ทำให้ผู้ประกอบการสามารถนำไปใช้งานได้ด้วยต้นทุนที่เข้าถึงได้
โดยไม่ต้องลงทุนวางโครงสร้างพื้นฐานเครือข่ายที่ซับซ้อน
อีกทั้งสถาปัตยกรรมที่แยกส่วนเป็นโมดูลยังสามารถนำไปต่อยอดกับงานหุ่นยนต์บริการอื่น ๆ
ได้ในอนาคต

หุ่นยนต์ต้นแบบที่พัฒนาขึ้นเพื่อตอบโจทย์ดังกล่าวแสดงดังรูปที่ \ref{fig:robot-overview}

\begin{figure}[htbp]
\centering
\includegraphics[width=0.7\textwidth]{images/robot_overview.jpg}
\caption[หุ่นยนต์ต้นแบบที่พัฒนาขึ้นในโครงงาน (บทที่ 1)]{หุ่นยนต์ต้นแบบที่พัฒนาขึ้นในโครงงาน}
\label{fig:robot-overview}
\end{figure}

% ---------------------------------------------------------------------
\section{วัตถุประสงค์}

\begin{enumerate}[label=\arabic*., leftmargin=1.25cm, itemsep=6pt, topsep=6pt]
    \item เพื่อออกแบบและพัฒนาอัลกอริทึม 2D SLAM โดยใช้ไลดาร์
    (Light Detection and Ranging: LiDAR)
    สำหรับการสร้างแผนที่เชิงพื้นที่และระบุตำแหน่งหุ่นยนต์ภายในสภาพแวดล้อมปิด
    โดยระบบต้องมีความคลาดเคลื่อนเชิงตำแหน่งไม่เกิน 10 เซนติเมตร
    เมื่อทดสอบในพื้นที่ควบคุมขนาดไม่น้อยกว่า 5$\times$5 เมตร

    \item เพื่อพัฒนาและปรับปรุงโมเดลการเรียนรู้เชิงลึก
    สำหรับการตรวจจับวัตถุเป้าหมายบนอุปกรณ์ประมวลผลเอดจ์เอไอ
    โดยโมเดลต้องมีค่าความแม่นยำไม่น้อยกว่าร้อยละ 80
    และมีความเร็วในการประมวลผลไม่น้อยกว่า 5 เฟรมต่อวินาที

    \item เพื่อสร้างระบบผสานข้อมูลระหว่างพิกัดตำแหน่งจาก SLAM
    และข้อมูลการตรวจจับวัตถุจากกล้อง สำหรับการระบุพิกัดของวัตถุลงบนแผนที่
    โดยตำแหน่งของวัตถุที่ปรากฏบนแผนที่ต้องมีความคลาดเคลื่อนเทียบกับตำแหน่งจริง
    ไม่เกิน 50 เซนติเมตร

    \item เพื่อประเมินความสามารถในการทำงานแบบเรียลไทม์บนอุปกรณ์ฝังตัว
    (Embedded System) โดยระบบต้องสามารถประมวลผลกระบวนการ SLAM
    และการตรวจจับวัตถุควบคู่กันได้อย่างเสถียร
    ภายใต้อัตราการใช้งานทรัพยากรประมวลผลหลักเฉลี่ยไม่เกินร้อยละ 80
    ตลอดระยะเวลาการทำงาน
\end{enumerate}

% ---------------------------------------------------------------------
\section{ขอบเขตของโครงงาน}

\begin{enumerate}[label=\arabic*., leftmargin=1.25cm, itemsep=6pt, topsep=6pt]
    \item ระบบครอบคลุมตั้งแต่การเก็บข้อมูลแผนที่ด้วยการประมวลผลจาก SLAM
    การประมวลผลด้วยปัญญาประดิษฐ์ การระบุพิกัดของวัตถุที่ตรวจพบลงบนแผนที่
    และการแสดงผลบนหน้าต่างการควบคุม

    \item ระบบรองรับการทำงานในอาคารที่พื้นผิวเรียบมีสิ่งกีดขวางจำนวนน้อย
    เช่น คลังสินค้า โกดัง หรืออาคารสำนักงาน
    โดยออกแบบให้ผู้ปฏิบัติงานทั่วไปที่ไม่มีความเชี่ยวชาญเฉพาะด้านหุ่นยนต์
    สามารถติดตามการทำงานของระบบได้ผ่านเว็บแอปพลิเคชัน

    \item โครงงานส่งมอบ (1) หุ่นยนต์เคลื่อนที่อัตโนมัติ
    (2) เว็บแอปพลิเคชันสำหรับติดตามการทำงานของหุ่นยนต์ (3) แผนที่จากการสร้างด้วยระบบ SLAM

    \item โครงงานนี้ไม่ครอบคลุมการพัฒนาระบบให้เข้าถึงพื้นที่เปิดโล่ง
    หรือที่มีสิ่งกีดขวางจำนวนมาก และการประมวลผลแบบ 3 มิติ
\end{enumerate}

% ---------------------------------------------------------------------
\section{เครื่องมือที่ใช้}

เครื่องมือที่ใช้ในการพัฒนาโครงงานแบ่งตามส่วนของระบบ แสดงดังตารางที่ \ref{tab:tools}

\begin{table}[htbp]
\centering
\caption[เครื่องมือที่ใช้ในการพัฒนาโครงงาน (บทที่ 1)]{เครื่องมือที่ใช้ในการพัฒนาโครงงาน}
\label{tab:tools}
\small
\begin{tabularx}{\textwidth}{|>{\raggedright\arraybackslash}p{3.2cm}|X|}
\hline
\textbf{ส่วนของระบบ} & \textbf{เครื่องมือที่ใช้} \\
\hline
Hardware &
\textbf{Advantech MIC-711-ON (NVIDIA Jetson Orin NX):} ประมวลผลหลัก
รัน SLAM / Nav2 / YOLO26l บนชุดซอฟต์แวร์ JetPack 6.1 (L4T 36.4) \newline
\textbf{YDLIDAR X3 Pro:} สแกนพื้นที่ 360$^\circ$ สร้างแผนที่และตรวจสิ่งกีดขวาง \newline
\textbf{Logitech C615 (USB Camera):} รับภาพขนาด 640$\times$480 พิกเซล ที่ 15 เฟรมต่อวินาที
ป้อนเข้าระบบ AI ตรวจจับวัตถุ \newline
\textbf{ESP32 DOIT DevKit V1:} ควบคุมมอเตอร์ด้วย PID + อ่านค่า Encoder ด้วยวงจรนับพัลส์ในตัวชิป \newline
\textbf{L298N:} โมดูลขับมอเตอร์ \newline
\textbf{GY-521 (MPU6050):} หน่วยวัดความเฉื่อย ใช้ค่าอัตราการหมุนรอบแกนตั้งในการผสานข้อมูล \newline
\textbf{DC Motor + Encoder:} ขับเคลื่อนหุ่นยนต์แบบ Differential Drive \\
\hline
Software / \newline Framework &
\textbf{Ubuntu 22.04:} ระบบปฏิบัติการหลักบน MIC-711-ON \newline
\textbf{ROS 2 Humble:} Framework กลางสำหรับสื่อสารระหว่าง Node ต่าง ๆ \newline
\textbf{slam\_toolbox:} สร้างแผนที่ 2D และระบุตำแหน่งหุ่นยนต์แบบเรียลไทม์ \newline
\textbf{Nav2 Stack:} วางแผนเส้นทางและหลีกเลี่ยงสิ่งกีดขวางอัตโนมัติ \newline
\textbf{ros2\_control:} ชั้นเชื่อมต่อฮาร์ดแวร์ (Hardware Interface) ระหว่าง ROS 2 กับ ESP32 \newline
\textbf{robot\_localization:} ตัวกรองคาลมานแบบขยายสำหรับผสานโอดอมิเตอร์กับไจโรสโคป \newline
\textbf{YOLO26l + TensorRT (FP16):}
ตรวจจับวัตถุตามคลาสของชุดข้อมูล COCO บน Edge AI \newline
\textbf{FastAPI:} API Server สำหรับเชื่อมต่อ UI กับระบบหุ่นยนต์ \newline
\textbf{Python:} เขียน Node ROS 2 / ประมวลผล AI / ผสานข้อมูล \\
\hline
User Interface &
\textbf{React.js:} พัฒนา Web Dashboard แสดงแผนที่ / ตำแหน่งวัตถุ / ภาพจากกล้อง \newline
\textbf{WebSocket:} รับ-ส่งข้อมูลแบบ Real-time ระหว่าง UI กับหุ่นยนต์ \\
\hline
Data Management &
\textbf{ไฟล์ JSON:} บันทึกตำแหน่งวัตถุที่ตรวจพบลงไฟล์โดยตรง ไม่ใช้ระบบฐานข้อมูล \newline
เนื่องจากจำนวนวัตถุต่อรอบการสำรวจมีไม่มากพอที่จะต้องใช้ฐานข้อมูล \\
\hline
Development Tools &
\textbf{VS Code:} IDE หลักในการพัฒนา Software \newline
\textbf{Git / GitHub:} จัดการ Source Code และ Version Control \\
\hline
\end{tabularx}
\end{table}

% ---------------------------------------------------------------------
\section{ประโยชน์ที่ได้จากโครงงาน}

\begin{enumerate}[label=\arabic*., leftmargin=1.25cm, itemsep=6pt, topsep=6pt]
    \item เพิ่มประสิทธิภาพและลดระยะเวลาในการบริหารจัดการพื้นที่ภายในอาคาร
    ผู้ใช้งานสามารถติดตามตำแหน่งและประเภทของวัตถุภายในพื้นที่ปิดได้แบบเรียลไทม์
    ช่วยลดเวลาที่สูญเสียไปจากการเดินสำรวจหรือค้นหาสิ่งของเป้าหมายด้วยแรงงานคน
    และลดความผิดพลาดในการระบุพิกัดตำแหน่งในระบบบันทึกข้อมูล

    \item ลดภาระงานเดินสำรวจซ้ำ ๆ ของบุคลากร
    ทำให้นำเวลาไปทำงานที่มีมูลค่าสูงกว่าได้

    \item ลดภาระการลงทุนด้านโครงสร้างพื้นฐานและเครือข่าย
    ด้วยความสามารถของระบบเอดจ์เอไอที่ประมวลผลได้จบในตัวอุปกรณ์
    ทำให้สถานประกอบการลดค่าใช้จ่ายในการวางระบบอินเทอร์เน็ตความเร็วสูง
    และลดความเสี่ยงจากการที่ระบบหยุดชะงักเมื่อสัญญาณเครือข่ายขัดข้อง

    \item สร้างองค์ความรู้และต้นแบบระบบบูรณาการ ROS 2 และเอดจ์เอไอ
    โครงงานนี้เป็นกรณีศึกษาการผสานข้อมูลจากเซ็นเซอร์ไลดาร์เข้ากับการเรียนรู้เชิงลึก
    บนอุปกรณ์ฝังตัว ซึ่งสามารถใช้เป็นแนวทางอ้างอิงให้นักพัฒนาหรือนักวิจัย
    นำไปประยุกต์กับงานหุ่นยนต์เคลื่อนที่อัตโนมัติหรือหุ่นยนต์บริการในอนาคต

    \item เป็นต้นแบบที่ขยายผลต่อยอดได้
    ด้วยสถาปัตยกรรมที่แยกส่วนการทำงานเป็นโมดูล
    (SLAM, การตรวจจับวัตถุ, การนำทาง)
    ทำให้สามารถนำไปปรับใช้กับหุ่นยนต์รุ่นอื่นหรือเพิ่มความสามารถใหม่ในอนาคตได้
    โดยไม่ต้องพัฒนาระบบใหม่ทั้งหมด
\end{enumerate}

% ---------------------------------------------------------------------
\section{ขั้นตอนและแผนการดำเนินงานของโครงงาน}

ขั้นตอนและแผนการดำเนินงานของโครงงานตลอดภาคการศึกษา
แสดงดังตารางที่ \ref{tab:gantt}

\begin{table}[htbp]
\centering
\caption[ขั้นตอนและแผนการดำเนินงานของโครงงาน (บทที่ 1)]{ขั้นตอนและแผนการดำเนินงานของโครงงาน}
\label{tab:gantt}
\small
\begin{tabularx}{\textwidth}{|>{\raggedright\arraybackslash}X|*{7}{>{\centering\arraybackslash}p{0.86cm}|}}
\hline
\multirow{2}{=}{\textbf{ขั้นตอนการดำเนินงานของโครงงาน}}
 & \multicolumn{7}{c|}{\textbf{สัปดาห์ที่ (ภาคการศึกษาที่ 1/2569)}} \\
\cline{2-8}
 & 1--2 & 3--4 & 5--6 & 7--8 & 9--10 & 11--12 & 13--14 \\
\hline
1. ศึกษาความต้องการ ทบทวนเทคโนโลยี และออกแบบ System Architecture
 & $\checkmark$ & & & & & & \\
\hline
2. ติดตั้ง ROS 2 และพัฒนาระบบควบคุมการเคลื่อนที่ (Motor Control / Odometry)
 & & $\checkmark$ & & & & & \\
\hline
3. พัฒนาระบบ SLAM ด้วย LiDAR เพื่อสร้างแผนที่และระบุตำแหน่ง
 & & & $\checkmark$ & & & & \\
\hline
4. พัฒนาและนำโมเดล YOLO26 ไปใช้งานบน Edge Device (ONNX และ TensorRT)
 & & & & $\checkmark$ & & & \\
\hline
5. พัฒนาระบบนำทางอัตโนมัติ Path Planning, Obstacle Avoidance และ Sensor Fusion
 & & & & & $\checkmark$ & & \\
\hline
6. รวมระบบทั้งหมด (Integration) และพัฒนา Web Dashboard
 & & & & & & $\checkmark$ & \\
\hline
7. ทดสอบประสิทธิภาพภาคสนาม ปรับปรุงข้อบกพร่อง
 & & & & & & & $\checkmark$ \\
\hline
\end{tabularx}
\end{table}

% ---------------------------------------------------------------------
\chapter{เทคโนโลยีที่เกี่ยวข้อง}

บทนี้อธิบายทฤษฎีและเทคโนโลยีที่ถูกนำมาใช้ในการพัฒนาโครงงาน
ประกอบด้วยเฟรมเวิร์กกลางสำหรับหุ่นยนต์ กระบวนการสร้างแผนที่และระบุตำแหน่ง
การตรวจจับวัตถุด้วยการเรียนรู้เชิงลึกบนอุปกรณ์ประมวลผลขอบ
การแปลงพิกัดระหว่างระบบพิกัด และเทคโนโลยีฝั่งเว็บแอปพลิเคชัน

\section{ระบบปฏิบัติการหุ่นยนต์ ROS 2}

ROS 2 (Robot Operating System 2)~\cite{ros2} ไม่ใช่ระบบปฏิบัติการโดยตรง
แต่เป็นชุดเครื่องมือและมิดเดิลแวร์ (Middleware) ที่ทำงานอยู่บนระบบปฏิบัติการอีกชั้นหนึ่ง
ทำหน้าที่จัดการการสื่อสารระหว่างโปรแกรมย่อยของหุ่นยนต์
โครงงานนี้ใช้ ROS 2 รุ่น Humble Hawksbill ซึ่งทำงานบน Ubuntu 22.04

โครงสร้างการทำงานของ ROS 2 ประกอบด้วยองค์ประกอบหลักดังนี้

\begin{itemize}[leftmargin=1.25cm, itemsep=4pt, topsep=6pt]
    \item \textbf{Node} คือหน่วยประมวลผลย่อยที่รับผิดชอบงานหนึ่งอย่าง
    เช่น โหนดอ่านค่าไลดาร์ โหนดคำนวณตำแหน่ง หรือโหนดตรวจจับวัตถุ
    การแยกเป็นโหนดย่อยทำให้ระบบมีลักษณะเป็นโมดูล
    เปลี่ยนหรือทดสอบเฉพาะส่วนได้โดยไม่กระทบส่วนอื่น

    \item \textbf{Topic} คือช่องทางสื่อสารแบบ Publish/Subscribe
    โหนดผู้ส่งจะเผยแพร่ข้อความลงหัวข้อ และโหนดใดก็ตามที่สนใจสามารถสมัครรับได้
    รูปแบบนี้ทำให้ผู้ส่งและผู้รับไม่ผูกติดกันโดยตรง

    \item \textbf{Service} และ \textbf{Action} ใช้สำหรับการสื่อสารแบบร้องขอ-ตอบกลับ
    โดย Action รองรับงานที่ใช้เวลานานและสามารถรายงานความคืบหน้าหรือยกเลิกกลางคันได้
    เช่น การสั่งให้หุ่นยนต์เคลื่อนที่ไปยังเป้าหมาย

    \item \textbf{DDS (Data Distribution Service)} เป็นมิดเดิลแวร์ชั้นล่างที่ ROS 2 ใช้
    ทำหน้าที่ค้นหาโหนดในเครือข่ายโดยอัตโนมัติและกำหนดคุณภาพการส่งข้อมูล
    (Quality of Service) เช่น ความน่าเชื่อถือและขนาดคิวข้อความ
\end{itemize}

ความสัมพันธ์ระหว่างโหนดและหัวข้อในการสื่อสารแบบ Publish/Subscribe แสดงดังรูปที่ \ref{fig:ros2-pubsub}

\begin{figure}[htbp]
\centering
\includegraphics[width=\textwidth]{images/ros2_pubsub.png}
\caption[การสื่อสารระหว่างโหนดแบบ Publish/Subscribe ใน ROS 2 (บทที่ 2)]{การสื่อสารระหว่างโหนดแบบ Publish/Subscribe ใน ROS 2}
\label{fig:ros2-pubsub}
\end{figure}

\subsection{ros2\_control และการเชื่อมต่อกับไมโครคอนโทรลเลอร์}

สำหรับการเชื่อมต่อกับไมโครคอนโทรลเลอร์ โครงงานนี้ใช้เฟรมเวิร์ก \textbf{ros2\_control}~\cite{ros2control}
แทนการรัน ROS 2 บนไมโครคอนโทรลเลอร์โดยตรง
ros2\_control แยกระบบออกเป็นสองชั้น ชั้นแรกคือ \textbf{Hardware Interface}
ซึ่งเป็นโค้ดภาษา C++ ที่ทำงานบนหน่วยประมวลผลหลัก
ทำหน้าที่อ่านค่าและสั่งงานฮาร์ดแวร์จริง
และชั้นที่สองคือ \textbf{Controller} ซึ่งเป็นอัลกอริทึมควบคุมที่ทำงานอยู่เหนือชั้นแรกอีกที

ในโครงงานนี้ Hardware Interface ชื่อ \texttt{DiffDriveSerial} สื่อสารกับ ESP32
ผ่านพอร์ตอนุกรมยูเอสบีด้วยโพรโทคอลข้อความอย่างง่าย
โดยส่งคำสั่งความเร็วของล้อทั้งสองข้างไปและอ่านค่าเอนโคเดอร์สะสมกลับมาในทุกรอบการควบคุม
จากนั้นคอนโทรลเลอร์ \texttt{diff\_drive\_controller}
(ชื่ออินสแตนซ์ \texttt{diff\_cont})
จะทำหน้าที่แปลงคำสั่งความเร็วเชิงเส้นและเชิงมุมเป็นความเร็วล้อ คำนวณโอดอมิเตอร์
และกระจายการแปลงพิกัดจาก \texttt{odom} ไปยัง \texttt{base\_link}

ข้อดีของแนวทางนี้เมื่อเทียบกับการใช้ micro-ROS บนไมโครคอนโทรลเลอร์คือ
ไมโครคอนโทรลเลอร์ไม่ต้องรันสแตก DDS ซึ่งกินหน่วยความจำมาก
และไม่ต้องอาศัยโปรแกรมตัวกลาง (Agent) ที่เพิ่มจุดล้มเหลวอีกจุดหนึ่งในระบบ
อีกทั้งการคำนวณโอดอมิเตอร์ยังใช้โค้ดมาตรฐานที่ผ่านการทดสอบมาแล้วของ ROS 2
แทนที่จะต้องเขียนขึ้นเอง

โครงสร้างการทำงานสองชั้นของ ros2\_control แสดงดังรูปที่ \ref{fig:ros2-control}

\begin{figure}[htbp]
\centering
\includegraphics[height=11cm,keepaspectratio]{images/ros2_control_layers.png}
\caption[โครงสร้างสองชั้นของ ros2\_control และการเชื่อมต่อกับไมโครคอนโทรลเลอร์ (บทที่ 2)]{โครงสร้างสองชั้นของ ros2\_control และการเชื่อมต่อกับไมโครคอนโทรลเลอร์}
\label{fig:ros2-control}
\end{figure}

\section{การสร้างแผนที่และระบุตำแหน่งพร้อมกัน (SLAM)}

SLAM (Simultaneous Localization and Mapping)
คือกระบวนการที่หุ่นยนต์สร้างแผนที่ของสภาพแวดล้อมที่ไม่เคยรู้จักมาก่อน~\cite{probrobotics}
พร้อมกับประมาณตำแหน่งของตนเองภายในแผนที่นั้นไปพร้อมกัน
ปัญหานี้มีลักษณะเป็นไก่กับไข่ กล่าวคือ การสร้างแผนที่ที่ถูกต้องต้องรู้ตำแหน่งของหุ่นยนต์
แต่การรู้ตำแหน่งก็ต้องอาศัยแผนที่เป็นตัวอ้างอิง

\subsection{องค์ประกอบของกระบวนการ SLAM}

\begin{itemize}[leftmargin=1.25cm, itemsep=4pt, topsep=6pt]
    \item \textbf{Scan Matching} คือการเปรียบเทียบข้อมูลการสแกนชุดปัจจุบัน
    กับข้อมูลชุดก่อนหน้าหรือกับแผนที่ที่สร้างไว้
    เพื่อหาการเลื่อนและการหมุนที่ทำให้ข้อมูลทับกันได้ดีที่สุด
    ผลลัพธ์คือการประมาณการเคลื่อนที่ของหุ่นยนต์จากข้อมูลไลดาร์โดยตรง

    \item \textbf{Odometry} คือการประมาณตำแหน่งจากการนับรอบหมุนของล้อ
    ใช้เป็นค่าตั้งต้นให้กับ Scan Matching เพื่อลดขอบเขตการค้นหาและเพิ่มความเร็ว
    ข้อจำกัดคือความคลาดเคลื่อนจะสะสมขึ้นเรื่อย ๆ ตามระยะทาง
    โดยเฉพาะเมื่อล้อเกิดการลื่นไถล

    \item \textbf{Pose Graph} คือการเก็บตำแหน่งของหุ่นยนต์ ณ ช่วงเวลาต่าง ๆ เป็นปม (Node)
    และเก็บความสัมพันธ์เชิงตำแหน่งระหว่างปมเป็นเส้นเชื่อม (Edge)
    เมื่อมีข้อมูลใหม่เข้ามา ระบบจะปรับตำแหน่งของปมทั้งกราฟให้สอดคล้องกันมากที่สุด

    \item \textbf{Loop Closure} คือการตรวจพบว่าหุ่นยนต์กลับมายังบริเวณที่เคยผ่านมาแล้ว
    เมื่อตรวจพบ ระบบจะเพิ่มเส้นเชื่อมใหม่เข้าไปในกราฟและปรับแก้ทั้งกราฟย้อนหลัง
    ซึ่งช่วยลบล้างความคลาดเคลื่อนสะสมได้อย่างมีนัยสำคัญ
\end{itemize}

ความสัมพันธ์ระหว่างองค์ประกอบทั้งสี่ของกระบวนการ SLAM แสดงดังรูปที่ \ref{fig:posegraph}

\begin{figure}[htbp]
\centering
\includegraphics[width=0.85\textwidth]{images/pose_graph_loop_closure.png}
\caption[โครงสร้าง Pose Graph และการปรับแก้ด้วย Loop Closure (บทที่ 2)]{โครงสร้าง Pose Graph และการปรับแก้ด้วย Loop Closure}
\label{fig:posegraph}
\end{figure}

\subsection{แผนที่แบบตารางความน่าจะเป็นและ slam\_toolbox}

ผลลัพธ์ของ SLAM ในโครงงานนี้อยู่ในรูปแผนที่แบบตารางความน่าจะเป็นในการครอบครองพื้นที่
(Occupancy Grid Map) ซึ่งแบ่งพื้นที่ออกเป็นช่องสี่เหลี่ยมขนาดเท่ากัน
แต่ละช่องเก็บค่าความน่าจะเป็นว่าช่องนั้นมีสิ่งกีดขวางอยู่หรือไม่
โดยแบ่งสถานะเป็นพื้นที่ว่าง พื้นที่มีสิ่งกีดขวาง และพื้นที่ที่ยังไม่ทราบข้อมูล

โครงงานนี้ใช้แพ็กเกจ slam\_toolbox~\cite{slamtoolbox} ซึ่งเป็นการนำอัลกอริทึมแบบ Pose Graph
มาพัฒนาต่อสำหรับ ROS 2 รองรับการทำงานหลายโหมด
โดยโหมดที่ใช้ในโครงงานคือโหมดออนไลน์แบบอะซิงโครนัส (Online Asynchronous)
ซึ่งประมวลผลข้อมูลไลดาร์ทันทีที่ได้รับและยอมข้ามข้อมูลบางเฟรมได้
เพื่อรักษาอัตราการทำงานแบบเรียลไทม์บนอุปกรณ์ที่มีทรัพยากรจำกัด

ลักษณะของแผนที่แบบตารางความน่าจะเป็นแสดงดังรูปที่ \ref{fig:occgrid}

\begin{figure}[htbp]
\centering
\includegraphics[width=0.85\textwidth]{images/occupancy_grid.png}
\caption[ตัวอย่างแผนที่แบบตารางความน่าจะเป็นในการครอบครองพื้นที่ (บทที่ 2)]{ตัวอย่างแผนที่แบบตารางความน่าจะเป็นในการครอบครองพื้นที่}
\label{fig:occgrid}
\end{figure}

\section{การผสานข้อมูลเซนเซอร์ด้วยตัวกรองคาลมานแบบขยาย}

โอดอมิเตอร์ที่คำนวณจากเอนโคเดอร์เพียงอย่างเดียวมีจุดอ่อนที่สำคัญคือ
เมื่อล้อลื่นไถลหรือหมุนอยู่กับที่โดยที่ตัวหุ่นยนต์ไม่ได้หมุนจริง
เอนโคเดอร์จะรายงานการหมุนที่ไม่ได้เกิดขึ้น
ความคลาดเคลื่อนของทิศทางที่เกิดขึ้นนี้เป็นความคลาดเคลื่อนแบบสะสมที่ SLAM แก้ให้ไม่ทัน
หากเกิดขึ้นเร็วกว่าขอบเขตการค้นหาของการจับคู่ข้อมูลสแกน

ตัวกรองคาลมานแบบขยาย (Extended Kalman Filter: EKF)
เป็นวิธีมาตรฐานในการรวมค่าที่วัดได้จากเซนเซอร์หลายตัวที่มีความไม่แน่นอนต่างกัน~\cite{probrobotics}
โดยถ่วงน้ำหนักแต่ละแหล่งข้อมูลตามค่าความแปรปรวนที่กำหนดให้
เซนเซอร์ที่ประกาศความแปรปรวนต่ำจะมีอิทธิพลต่อค่าประมาณสุดท้ายมากกว่า

ในงานหุ่นยนต์เคลื่อนที่ การผสานที่ให้ผลคุ้มค่าที่สุดคือการรวม
\textbf{อัตราการหมุนรอบแกนตั้งจากไจโรสโคป} เข้ากับโอดอมิเตอร์จากล้อ
เนื่องจากไจโรสโคปวัดการหมุนของตัวหุ่นยนต์โดยตรงและไม่ได้รับผลจากการลื่นไถลของล้อ
ขณะที่ล้อให้ข้อมูลความเร็วเชิงเส้นที่ไจโรสโคปให้ไม่ได้
ข้อแลกเปลี่ยนคือไจโรสโคปแบบไมโครอิเล็กโตรแมคานิกส์มีค่าเบี่ยงเบนศูนย์ (Bias)
ที่ต้องวัดและหักออก และมีการเดินสุ่ม (Random Walk) ของค่าที่อ่านได้
ทำให้ขณะหุ่นยนต์หยุดนิ่งสนิท ค่าประมาณทิศทางจากไจโรสโคปจะแย่กว่าจากเอนโคเดอร์เสมอ

\section{ระบบนำทางอัตโนมัติ Nav2}

Nav2~\cite{nav2} เป็นชุดซอฟต์แวร์นำทางมาตรฐานของ ROS 2 ที่รับตำแหน่งเป้าหมายบนแผนที่
แล้ววางแผนเส้นทางและสั่งความเร็วให้หุ่นยนต์เคลื่อนที่ไปถึงโดยหลบสิ่งกีดขวางระหว่างทาง
องค์ประกอบหลักที่เกี่ยวข้องกับโครงงานนี้ได้แก่

\begin{itemize}[leftmargin=1.25cm, itemsep=4pt, topsep=6pt]
    \item \textbf{Costmap} คือแผนที่ต้นทุนที่สร้างทับบนแผนที่จาก SLAM
    โดยขยายขอบของสิ่งกีดขวางออกมาเป็นบริเวณที่หุ่นยนต์ไม่ควรเข้าใกล้
    แบ่งเป็นแผนที่ต้นทุนส่วนรวม (Global) สำหรับวางแผนเส้นทางทั้งเส้น
    และแผนที่ต้นทุนเฉพาะที่ (Local) สำหรับหลบสิ่งกีดขวางที่เพิ่งปรากฏ

    \item \textbf{Planner} และ \textbf{Controller} ทำหน้าที่คำนวณเส้นทางจากตำแหน่งปัจจุบัน
    ไปยังเป้าหมาย และแปลงเส้นทางนั้นเป็นคำสั่งความเร็วในแต่ละช่วงเวลาตามลำดับ

    \item \textbf{Behavior Tree} เป็นโครงสร้างต้นไม้ที่กำหนดลำดับการตัดสินใจของระบบนำทาง
    รวมถึงพฤติกรรมกู้คืนสถานการณ์ เช่น การหมุนตัวหรือการถอยหลังเมื่อวางแผนเส้นทางไม่สำเร็จ

    \item \textbf{รูปร่างของหุ่นยนต์ (Footprint)} ใช้กำหนดขอบเขตเชิงกายภาพที่ระบบต้องหลีกเลี่ยง
    การระบุเป็นรูปหลายเหลี่ยมให้ผลดีกว่าการระบุเป็นรัศมีวงกลม
    เมื่อหุ่นยนต์มีรูปทรงที่ไม่สมมาตร เพราะวงกลมที่ครอบรูปทรงทั้งหมดจะใหญ่เกินจริง
\end{itemize}


\section{การตรวจจับวัตถุด้วยการเรียนรู้เชิงลึกบนอุปกรณ์เอดจ์}

\subsection{การตรวจจับวัตถุด้วยโครงข่ายประสาทเทียมเชิงลึก}

การตรวจจับวัตถุ (Object Detection) คือการระบุทั้งชนิดและตำแหน่งของวัตถุในภาพ
ซึ่งต่างจากการจำแนกภาพที่ตอบเพียงว่าภาพทั้งภาพเป็นวัตถุชนิดใด
ผลลัพธ์ของการตรวจจับวัตถุคือกรอบสี่เหลี่ยมล้อมรอบวัตถุ (Bounding Box)
พร้อมชื่อคลาสและค่าความเชื่อมั่น

แบบจำลองการตรวจจับวัตถุแบ่งได้เป็นสองกลุ่มใหญ่
กลุ่มแรกคือแบบสองขั้นตอน (Two-stage) ที่เสนอบริเวณที่น่าสนใจก่อนแล้วจึงจำแนก
ให้ความแม่นยำสูงแต่ช้า และกลุ่มที่สองคือแบบขั้นตอนเดียว (One-stage)
ที่ทำนายกรอบและคลาสพร้อมกันในการประมวลผลรอบเดียว
ตระกูล YOLO (You Only Look Once) จัดอยู่ในกลุ่มหลัง
จึงเหมาะกับงานที่ต้องการความเร็วสูงบนอุปกรณ์ฝังตัว

ลักษณะผลลัพธ์ของการตรวจจับวัตถุแสดงดังรูปที่ \ref{fig:yolo-example}

\begin{figure}[htbp]
\centering
\includegraphics[height=8cm,keepaspectratio]{images/yolo_bbox_example.jpg}
\caption[ตัวอย่างผลลัพธ์ของการตรวจจับวัตถุด้วยแบบจำลองแบบขั้นตอนเดียว (บทที่ 2)]{ตัวอย่างผลลัพธ์ของการตรวจจับวัตถุด้วยแบบจำลองแบบขั้นตอนเดียว}
\label{fig:yolo-example}
\end{figure}

\subsection{แบบจำลอง YOLO26}

โครงงานนี้ใช้แบบจำลอง YOLO26~\cite{yolo26} ผ่านไลบรารี Ultralytics~\cite{ultralytics}
ซึ่งเป็นตระกูลแบบจำลองที่ออกแบบมาสำหรับอุปกรณ์ประมวลผลขอบและอุปกรณ์กำลังไฟต่ำ
โดยเริ่มจากรุ่น Small (YOLO26s) และเปลี่ยนเป็นรุ่น Large (YOLO26l)
หลังการประเมินความแม่นยำในพื้นที่ใช้งานจริง

การเลือกรุ่น Small แทนรุ่น Nano ที่เล็กกว่ามาจากผลการวัดบนอุปกรณ์จริง
ซึ่งพบว่าเวลาประมวลผลของทั้งสองรุ่นต่างกันเพียงเล็กน้อย
เนื่องจากคอขวดของระบบอยู่ที่ค่าใช้จ่ายคงที่ในการเรียกใช้แบบจำลองแต่ละครั้ง
ไม่ได้อยู่ที่ปริมาณการคำนวณของตัวแบบจำลอง
การใช้รุ่นที่ใหญ่ขึ้นจึงแทบไม่มีต้นทุนด้านความเร็วเพิ่มขึ้น
ต่อมาการประเมินบนภาพจากพื้นที่ใช้งานจริงพบว่ารุ่น Small ตรวจไม่พบวัตถุระยะไกลเป็นส่วนใหญ่
จึงเปลี่ยนเป็นรุ่น Large ซึ่งเมื่อทำงานบน TensorRT ใช้เวลาอนุมานใกล้เคียงกับรุ่น Small เดิม
(รายละเอียดในหัวข้อ \ref{sec:eval-det})

จุดเด่นทางสถาปัตยกรรมที่เกี่ยวข้องกับโครงงานนี้ได้แก่

\begin{itemize}[leftmargin=1.25cm, itemsep=4pt, topsep=6pt]
    \item การอนุมานแบบต่อเนื่องปลายทางถึงปลายทาง (End-to-End)
    โดยไม่ต้องใช้ขั้นตอน Non-Maximum Suppression แยกต่างหาก
    ช่วยลดเวลาประมวลผลหลังการอนุมานและทำให้การแปลงแบบจำลองไปยังรูปแบบอื่นทำได้ง่ายขึ้น

    \item การตัดส่วน Distribution Focal Loss ออกจากส่วนหัวของการตรวจจับ
    ทำให้โครงสร้างเรียบง่ายลงและเอื้อต่อการทำ Quantization

    \item ประสิทธิภาพต่อหน่วยพลังงานที่ดีขึ้นบนอุปกรณ์ตระกูล NVIDIA Jetson
\end{itemize}

แบบจำลองที่ผ่านการฝึกมาแล้วบนชุดข้อมูล COCO~\cite{coco} รองรับการจำแนกวัตถุจำนวน 80 คลาส
ซึ่งเป็นชุดคลาสแบบปิด กล่าวคือระบบจะตรวจจับได้เฉพาะวัตถุที่อยู่ในรายการคลาสเหล่านี้เท่านั้น
หากต้องการตรวจจับวัตถุนอกเหนือจากนี้ จำเป็นต้องเก็บชุดข้อมูลเพิ่มและฝึกแบบจำลองซ้ำ

\subsection{การเร่งการอนุมาน: TensorRT และ ONNX Runtime}

TensorRT~\cite{tensorrt} เป็นชุดเครื่องมือของ NVIDIA สำหรับปรับแต่งแบบจำลองให้ทำงานเร็วขึ้นบนหน่วยประมวลผลกราฟิก
กระบวนการทำงานคือการแปลงแบบจำลองจากรูปแบบมาตรฐานให้กลายเป็นเครื่องยนต์อนุมาน (Engine)
ที่ผูกกับสถาปัตยกรรมของอุปกรณ์เป้าหมายโดยเฉพาะ
เทคนิคที่ใช้ในการปรับแต่งประกอบด้วยการรวมชั้นการคำนวณที่ต่อเนื่องกันให้เป็นชั้นเดียว (Layer Fusion)
การเลือกอัลกอริทึมที่เร็วที่สุดสำหรับฮาร์ดแวร์นั้น
และการลดความละเอียดของการคำนวณจากทศนิยม 32 บิตเหลือ 16 บิต (FP16)
ซึ่งช่วยเพิ่มความเร็วและลดการใช้หน่วยความจำโดยที่ความแม่นยำลดลงเพียงเล็กน้อย

ข้อควรระวังคือเครื่องยนต์อนุมานที่สร้างขึ้นจะใช้ได้เฉพาะกับอุปกรณ์และรุ่นของไลบรารีที่สร้างมันเท่านั้น
จึงต้องสร้างบนอุปกรณ์เป้าหมายจริงเสมอ

\textbf{ในช่วงแรกโครงงานพิจารณา TensorRT แล้วเลือกไม่ใช้}
ด้วยเหตุผลจากประสบการณ์จริงระหว่างการพัฒนา
กล่าวคือเมื่อจำเป็นต้องเปลี่ยนรุ่นของชุดซอฟต์แวร์ JetPack บนอุปกรณ์
เครื่องยนต์อนุมานที่สร้างไว้เดิมไม่สามารถโหลดได้อีกต่อไป
เพราะผูกกับรุ่นของไลบรารี TensorRT ที่สร้างมัน
ซึ่งเป็นข้อจำกัดเดียวกับที่อธิบายไว้ในย่อหน้าก่อนหน้า

ระบบจึงเปลี่ยนมาใช้รูปแบบ \textbf{ONNX (Open Neural Network Exchange)}
ร่วมกับ ONNX Runtime~\cite{onnxruntime} ที่เรียกใช้หน่วยประมวลผลกราฟิกผ่าน CUDA Execution Provider แทน
ไฟล์ ONNX เป็นรูปแบบกลางที่อธิบายโครงสร้างของแบบจำลอง ไม่ผูกกับรุ่นของไลบรารีอนุมาน
จึงยังใช้งานได้แม้เปลี่ยนรุ่นชุดซอฟต์แวร์ของอุปกรณ์
โดยยังคงส่งออกด้วยความละเอียด 16 บิต (FP16) เช่นเดียวกัน

จากการวัดบนอุปกรณ์จริงที่ล็อกสัญญาณนาฬิกาของหน่วยประมวลผลกราฟิก (ตารางที่ \ref{tab:modelbench} ในบทที่ 4)
ONNX แบบ FP16 ใช้เวลาประมวลผลน้อยกว่าการเรียกใช้แบบจำลอง
ผ่าน PyTorch โดยตรงเล็กน้อย ขณะที่ ONNX แบบ FP32 ช้ากว่าอย่างมีนัยสำคัญ
การส่งออกเป็น FP16 จึงเป็นเงื่อนไขที่จำเป็น ไม่ใช่ทางเลือก

\textbf{เมื่อเปลี่ยนเป็นรุ่น Large ระบบจึงกลับมาใช้ TensorRT}
เพราะรุ่น Large ผ่าน PyTorch ใช้เวลาอนุมานประมาณสองเท่าของรุ่น Small
ขณะที่เครื่องยนต์อนุมาน TensorRT แบบ FP16 ลดเวลาลงร้อยละ 51
และใช้หน่วยประมวลผลกราฟิกน้อยกว่ารุ่น Small แบบ ONNX ที่ใช้เดิม
ปัญหาเรื่องการผูกกับรุ่นของไลบรารีจัดการได้โดยเก็บไฟล์แบบจำลองต้นฉบับไว้เป็นแหล่งที่มา
หากเปลี่ยนรุ่นชุดซอฟต์แวร์ของอุปกรณ์ ต้องสร้างเครื่องยนต์ใหม่บนอุปกรณ์ซึ่งใช้เวลาประมาณ 13 นาที
ไม่ได้สูญเสียแบบจำลอง

ผลการวัดเวลาอนุมานของแบบจำลองที่พิจารณาทั้งหมด รวมถึงรุ่นที่ใช้งานจริง
ภายใต้เงื่อนไขเดียวกันบนอุปกรณ์จริง แสดงดังรูปที่ \ref{fig:infer-bench}

\begin{figure}[htbp]
\centering
\includegraphics[width=0.85\textwidth]{images/inference_time_comparison.png}
\caption[เปรียบเทียบเวลาอนุมานต่อเฟรมของแบบจำลองและรูปแบบที่ต่างกัน (บทที่ 2)]{เปรียบเทียบเวลาอนุมานต่อเฟรมของแบบจำลองและรูปแบบที่ต่างกันบน Jetson Orin NX
ตัวเลขภายในแท่งคืออัตราการใช้งานหน่วยประมวลผลกราฟิก
YOLO26l บน TensorRT แบบ FP16 ใช้เวลาเท่ากับ YOLO26s แบบ ONNX ที่ใช้เดิม และใช้หน่วยประมวลผลกราฟิกน้อยกว่า
(วัดที่โหมดพลังงาน 15 วัตต์ด้วยภาพจริงจากกล้อง ขณะไม่มีงานอื่นใช้หน่วยประมวลผลกราฟิก
ตัวเลขจึงต่างจากตารางที่ \ref{tab:modelbench} ซึ่งวัดด้วยภาพทดสอบและล็อกสัญญาณนาฬิกา)}
\label{fig:infer-bench}
\end{figure}

\section{การแปลงพิกัดและการระบุตำแหน่งวัตถุบนแผนที่}

ระบบหุ่นยนต์ประกอบด้วยระบบพิกัดหลายระบบที่ต้องอ้างอิงถึงกัน
ROS 2 จัดการเรื่องนี้ด้วยไลบรารี TF2 ซึ่งเก็บความสัมพันธ์ของระบบพิกัดทั้งหมด
ในรูปโครงสร้างต้นไม้พร้อมประทับเวลา ทำให้สามารถสอบถามการแปลงพิกัด
ระหว่างเฟรมใดกับเฟรมใดก็ได้ ณ เวลาที่ต้องการ

ระบบพิกัดหลักที่ใช้ในโครงงานนี้ตั้งชื่อตามมาตรฐาน REP 105~\cite{rep105} ประกอบด้วย \texttt{map} ซึ่งเป็นเฟรมอ้างอิงของแผนที่
\texttt{odom} ซึ่งเป็นเฟรมที่ต่อเนื่องแต่มีความคลาดเคลื่อนสะสม
\texttt{base\_footprint} ซึ่งเป็นเงาของหุ่นยนต์บนพื้นและเป็นเฟรมที่ slam\_toolbox ใช้อ้างอิง
\texttt{base\_link} ซึ่งยึดติดกับตัวหุ่นยนต์
และเฟรมของเซนเซอร์แต่ละตัวคือ \texttt{laser\_frame} และ \texttt{camera\_link}

การระบุตำแหน่งของวัตถุลงบนแผนที่อาศัยหลักการฉายจุดบนภาพกลับไปเป็นทิศทางในปริภูมิสามมิติ
โดยใช้พารามิเตอร์ภายในของกล้อง (Intrinsic Parameters) ที่ได้จากการสอบเทียบ~\cite{zhang2000}
จากนั้นใช้ระยะทางจากไลดาร์ในทิศทางเดียวกันเป็นตัวกำหนดความลึก
แล้วแปลงพิกัดผลลัพธ์จากเฟรมกล้องไปยังเฟรมแผนที่ผ่าน TF2

\section{เทคโนโลยีฝั่งเว็บแอปพลิเคชัน}

ส่วนแสดงผลของระบบพัฒนาด้วยสถาปัตยกรรมแบบแยกส่วนหน้าและส่วนหลัง
ส่วนหลังใช้ FastAPI~\cite{fastapi} ซึ่งเป็นเว็บเฟรมเวิร์กภาษาไพทอนที่รองรับการทำงานแบบอะซิงโครนัส
ทำหน้าที่เปิดปลายทาง REST API สำหรับการเรียกข้อมูลแบบครั้งต่อครั้ง
และเปิดช่องทาง WebSocket สำหรับการส่งข้อมูลแบบสองทางอย่างต่อเนื่อง
ซึ่งจำเป็นสำหรับการแสดงตำแหน่งหุ่นยนต์และผลการตรวจจับแบบเรียลไทม์
โดยไม่ต้องให้เบราว์เซอร์คอยร้องขอซ้ำ ๆ

ส่วนหน้าใช้ React.js~\cite{react} สำหรับสร้างส่วนติดต่อผู้ใช้แบบหน้าเดียว
โดยรับข้อมูลทั้งหมดจากส่วนหลังผ่าน REST API และ WebSocket
ระบบไม่ได้ใช้ฐานข้อมูล แต่บันทึกรายการวัตถุที่พบลงไฟล์ JSON บนตัวอุปกรณ์โดยตรง
ซึ่งเพียงพอกับปริมาณข้อมูลระดับหลักสิบรายการต่อรอบการสำรวจ

% ---------------------------------------------------------------------
\chapter{การออกแบบระบบงาน}

บทนี้อธิบายการออกแบบระบบตั้งแต่การรวบรวมความต้องการ
ภาพรวมสถาปัตยกรรมและลำดับการไหลของข้อมูล
ไปจนถึงการออกแบบเชิงวัตถุในระดับคลาส

\section{Requirement ของระบบงาน}

ความต้องการของระบบแบ่งออกเป็นความต้องการเชิงหน้าที่ (Functional Requirement)
ซึ่งระบุว่าระบบต้องทำอะไรได้บ้าง และความต้องการที่ไม่ใช่เชิงหน้าที่
(Non-Functional Requirement) ซึ่งระบุคุณภาพและข้อจำกัดของระบบ
รายละเอียดแสดงดังตารางที่ \ref{tab:fr} และตารางที่ \ref{tab:nfr}

\begin{table}[htbp]
\centering
\caption[ความต้องการเชิงหน้าที่ของระบบ (บทที่ 3)]{ความต้องการเชิงหน้าที่ของระบบ}
\label{tab:fr}
\small
\begin{tabularx}{\textwidth}{|c|X|}
\hline
\textbf{รหัส} & \textbf{รายละเอียด} \\
\hline
FR-01 & ระบบต้องสร้างแผนที่สองมิติของพื้นที่ปิดจากข้อมูลไลดาร์ได้แบบเรียลไทม์ \\
\hline
FR-02 & ระบบต้องประมาณตำแหน่งและทิศทางของหุ่นยนต์บนแผนที่ได้ตลอดเวลาที่ทำงาน \\
\hline
FR-03 & ระบบต้องขับเคลื่อนหุ่นยนต์ตามคำสั่งความเร็วเชิงเส้นและเชิงมุมที่ได้รับ \\
\hline
FR-04 & ระบบต้องหยุดมอเตอร์โดยอัตโนมัติเมื่อขาดการสื่อสารกับหน่วยประมวลผลหลัก \\
\hline
FR-05 & ระบบต้องตรวจจับและจำแนกชนิดวัตถุจากภาพกล้องด้วยแบบจำลองที่ทำงานในตัวอุปกรณ์ \\
\hline
FR-06 & ระบบต้องแปลงตำแหน่งวัตถุที่ตรวจพบจากพิกัดภาพให้เป็นพิกัดบนแผนที่ \\
\hline
FR-07 & ระบบต้องบันทึกชนิด พิกัด และเวลาที่พบวัตถุลงไฟล์ข้อมูล \\
\hline
FR-08 & ระบบต้องแสดงแผนที่ ตำแหน่งหุ่นยนต์ และตำแหน่งวัตถุที่พบบนเว็บแอปพลิเคชัน \\
\hline
FR-09 & ผู้ใช้ต้องสามารถสั่งล้างรายการวัตถุที่บันทึกไว้ผ่านเว็บแอปพลิเคชันได้ \\
\hline
FR-10 & ระบบต้องส่งออกแผนที่ที่สร้างเสร็จแล้วเป็นไฟล์เพื่อนำไปใช้อ้างอิงภายหลังได้
        ผ่านบริการบันทึกแผนที่ของ slam\_toolbox \\
\hline
FR-11 & ระบบต้องนำทางไปยังเป้าหมายที่ผู้ปฏิบัติงานกำหนดบนแผนที่
        พร้อมหลบสิ่งกีดขวางระหว่างทางได้ \\
\hline
FR-12 & ระบบต้องรองรับการผสานข้อมูลเอนโคเดอร์กับไจโรสโคปเพื่อประมาณทิศทางของหุ่นยนต์
        โดยเปิดใช้เป็นตัวเลือกเสริม \\
\hline
FR-13 & ระบบต้องสำรวจพื้นที่ได้เองโดยเลือกเป้าหมายที่ขอบเขตของพื้นที่ที่ยังไม่ทราบข้อมูล
        และกลับสู่จุดเริ่มต้นเมื่อสำรวจเสร็จ \\
\hline
\end{tabularx}
\end{table}

\begin{table}[htbp]
\centering
\caption[ความต้องการที่ไม่ใช่เชิงหน้าที่ของระบบ (บทที่ 3)]{ความต้องการที่ไม่ใช่เชิงหน้าที่ของระบบ}
\label{tab:nfr}
\small
\begin{tabularx}{\textwidth}{|c|X|}
\hline
\textbf{รหัส} & \textbf{รายละเอียด} \\
\hline
NFR-01 & ความคลาดเคลื่อนเชิงตำแหน่งของ SLAM ต้องไม่เกิน 10 เซนติเมตร
         ในพื้นที่ทดสอบขนาดไม่น้อยกว่า 5$\times$5 เมตร \\
\hline
NFR-02 & แบบจำลองตรวจจับวัตถุต้องมีความแม่นยำไม่น้อยกว่าร้อยละ 80
         และประมวลผลได้ไม่น้อยกว่า 5 เฟรมต่อวินาที \\
\hline
NFR-03 & ตำแหน่งวัตถุบนแผนที่ต้องคลาดเคลื่อนจากตำแหน่งจริงไม่เกิน 50 เซนติเมตร \\
\hline
NFR-04 & อัตราการใช้งานหน่วยประมวลผลหลักเฉลี่ยต้องไม่เกินร้อยละ 80
         ขณะรัน SLAM และการตรวจจับวัตถุพร้อมกัน \\
\hline
NFR-05 & การประมวลผลทั้งหมดต้องเกิดขึ้นภายในตัวอุปกรณ์ โดยไม่ส่งข้อมูลไปประมวลผลบนคลาวด์ \\
\hline
NFR-06 & ผู้ใช้ที่ไม่มีความเชี่ยวชาญด้านหุ่นยนต์ต้องสามารถใช้งานผ่านเว็บเบราว์เซอร์ได้
         โดยไม่ต้องติดตั้งโปรแกรมเพิ่ม \\
\hline
NFR-07 & ระบบต้องทำงานต่อเนื่องได้อย่างเสถียรตลอดรอบการสำรวจหนึ่งรอบ
         โดยไม่ต้องรีสตาร์ตโหนดใด \\
\hline
\end{tabularx}
\end{table}

\section{การทำงานของระบบ}

สถาปัตยกรรมภาพรวมของระบบแสดงดังรูปที่ \ref{fig:arch}

\begin{figure}[htbp]
\centering
\includegraphics[width=\textwidth]{system_architecture.jpg}
\caption[สถาปัตยกรรมภาพรวมของระบบ (บทที่ 3)]{สถาปัตยกรรมภาพรวมของระบบ (System Architecture)}
\label{fig:arch}
\end{figure}

การเชื่อมต่อระหว่างอุปกรณ์ฮาร์ดแวร์แต่ละส่วนแสดงดังรูปที่ \ref{fig:hwblock}

\begin{figure}[htbp]
\centering
\includegraphics[width=\textwidth]{images/hardware_block_diagram.png}
\caption[แผนภาพการเชื่อมต่อฮาร์ดแวร์ของระบบ (บทที่ 3)]{แผนภาพการเชื่อมต่อฮาร์ดแวร์ของระบบ}
\label{fig:hwblock}
\end{figure}

ลำดับการทำงานของระบบ (Flow) เป็นดังนี้

\begin{enumerate}[label=\arabic*., leftmargin=1.25cm, itemsep=4pt, topsep=6pt]
    \item YDLiDAR + USB Camera $\rightarrow$
          เก็บข้อมูลสภาพแวดล้อมและภาพแบบ Real-time
    \item MIC-711-ON (ROS 2 + slam\_toolbox) $\rightarrow$
          ประมวลผลสร้างแผนที่และระบุตำแหน่งหุ่นยนต์
    \item YOLO26l + TensorRT $\rightarrow$
          ตรวจจับและระบุชนิดวัตถุจากภาพกล้อง
    \item FastAPI + WebSocket $\rightarrow$
          ส่งข้อมูลแบบ Real-time ผ่าน WiFi ไปยัง Web Dashboard
    \item ไฟล์ JSON $\rightarrow$
          บันทึกข้อมูลการพบวัตถุและตำแหน่งบนแผนที่ลงไฟล์
    \item React.js $\rightarrow$
          แสดงผลแผนที่ ตำแหน่งหุ่นยนต์ และตำแหน่งวัตถุที่พบผ่านเบราว์เซอร์
\end{enumerate}

เมื่อพิจารณาในระดับโหนดและหัวข้อของ ROS 2 การไหลของข้อมูลทั้งระบบแสดงดังรูปที่ \ref{fig:nodegraph}

\begin{figure}[htbp]
\centering
\includegraphics[width=\textwidth]{images/ros2_node_graph.png}
\caption[แผนภาพการไหลของข้อมูลระหว่างโหนดของระบบ (บทที่ 3)]{แผนภาพการไหลของข้อมูลระหว่างโหนดของระบบ
อ่านจากระบบจริงขณะเปิดครบทุกส่วน (วงรีคือโหนด สี่เหลี่ยมคือหัวข้อ ลูกศรคือทิศทางการไหลของข้อมูล)
แสดงเฉพาะหัวข้อที่มีทั้งผู้เผยแพร่และผู้รับ
และไม่แสดง \texttt{/tf} \texttt{/rosout} หัวข้อภายในของ Action และแผนที่ต้นทุนของ Nav2}
\label{fig:nodegraph}
\end{figure}

\section{Use Case Diagram ของระบบ}

ผู้เกี่ยวข้องกับระบบ (Actor) แบ่งเป็นสองกลุ่ม
กลุ่มแรกคือผู้ปฏิบัติงานซึ่งเป็นผู้ใช้ทั่วไปที่ติดตามการทำงานผ่านเว็บแอปพลิเคชัน
และสั่งงานหุ่นยนต์ผ่านคำสั่งของระบบ
และกลุ่มที่สองคือตัวหุ่นยนต์ซึ่งทำหน้าที่เป็นผู้กระทำในระบบเช่นกัน
เพราะเป็นฝ่ายเริ่มต้นกระบวนการสร้างแผนที่และตรวจจับวัตถุด้วยตนเอง
แผนภาพยูสเคสของระบบแสดงดังรูปที่ \ref{fig:usecase}
และรายละเอียดของแต่ละยูสเคสแสดงดังตารางที่ \ref{tab:usecase}

\begin{figure}[htbp]
\centering
\begin{adjustbox}{max width=\textwidth, max totalheight=0.72\textheight}
\begin{tikzpicture}[
    font=\small,
    uc/.style={draw=black, ellipse, text width=4.4cm, align=center,
               minimum height=1.1cm, inner sep=1pt}
]
% ── กรอบขอบเขตระบบ ────────────────────────────────────────────────
\node[draw=black, minimum width=7.6cm, minimum height=14.0cm,
      anchor=north] (sys) at (0,0.6) {};
\node[anchor=north, align=center, font=\small\bfseries]
      at ([yshift=-4pt]sys.north) {ระบบหุ่นยนต์สำรวจพื้นที่\\และระบุตำแหน่งวัตถุ};

% ── ยูสเคส ───────────────────────────────────────────────────────
\node[uc] (u1) at (0,-1.80)  {UC-01 สั่งเคลื่อนที่และ\\กำหนดเป้าหมายนำทาง};
\node[uc] (u2) at (0,-3.60)  {UC-02 ล้างรายการวัตถุ\\ที่บันทึกไว้};
\node[uc] (u3) at (0,-5.20)  {UC-03 ติดตามสถานะ};
\node[uc] (u4) at (0,-6.60)  {UC-04 ดูรายการวัตถุที่พบ};
\node[uc] (u5) at (0,-8.00)  {UC-05 บันทึกแผนที่};
\node[uc] (u6) at (0,-9.40)  {UC-06 สร้างแผนที่};
\node[uc] (u7) at (0,-10.80) {UC-07 ตรวจจับวัตถุ};
\node[uc] (u8) at (0,-12.20) {UC-08 ระบุพิกัดวัตถุ};

% ── ผู้กระทำ (Actor) ─────────────────────────────────────────────
%  #1 = ตำแหน่งจุดกึ่งกลางลำตัว, #2 = ชื่อ (อยู่ใต้รูป)
\def\actor#1#2{%
  \begin{scope}[shift={#1}]
    \draw[black] (0,0.55) circle (0.22);
    \draw[black] (0,0.33) -- (0,-0.25);
    \draw[black] (-0.32,0.10) -- (0.32,0.10);
    \draw[black] (0,-0.25) -- (-0.26,-0.72);
    \draw[black] (0,-0.25) -- (0.26,-0.72);
    \node[below, font=\small, align=center] at (0,-0.78) {#2};
  \end{scope}}

\actor{(-5.2,-5.0)}{ผู้ปฏิบัติงาน}
\actor{(5.2,-10.8)}{หุ่นยนต์}

% ── เส้นเชื่อม (จากข้างลำตัวของ Actor) ──────────────────────────────
\foreach \u in {u1,u2,u3,u4,u5} \draw[black] (-4.85,-5.0) -- (\u.west);
\foreach \u in {u6,u7,u8}       \draw[black] (4.85,-10.8) -- (\u.east);
\end{tikzpicture}
\end{adjustbox}
\caption[แผนภาพยูสเคสของระบบ (บทที่ 3)]{แผนภาพยูสเคสของระบบ}
\label{fig:usecase}
\end{figure}

\begin{table}[htbp]
\centering
\caption[รายละเอียดยูสเคสของระบบ (บทที่ 3)]{รายละเอียดยูสเคสของระบบ}
\label{tab:usecase}
\small
\begin{tabularx}{\textwidth}{|c|>{\raggedright\arraybackslash}p{3.6cm}|c|X|}
\hline
\textbf{รหัส} & \textbf{ชื่อยูสเคส} & \textbf{ผู้กระทำ} & \textbf{คำอธิบาย} \\
\hline
UC-01 & สั่งเคลื่อนที่และกำหนดเป้าหมายนำทาง & ผู้ปฏิบัติงาน
      & ขับหุ่นยนต์ด้วยคำสั่งจากแป้นพิมพ์ กำหนดตำแหน่งเป้าหมายบนแผนที่
        ให้ระบบนำทางวางแผนเส้นทางไปเอง หรือสั่งให้สำรวจพื้นที่อัตโนมัติ \\
\hline
UC-02 & ล้างรายการวัตถุที่บันทึกไว้ & ผู้ปฏิบัติงาน
      & สั่งล้างรายการวัตถุทั้งหมดผ่านเว็บแอปพลิเคชันเพื่อเริ่มการสำรวจรอบใหม่ \\
\hline
UC-03 & ติดตามสถานะ & ผู้ปฏิบัติงาน & ดูแผนที่ ตำแหน่งหุ่นยนต์ และภาพจากกล้องแบบเรียลไทม์ \\
\hline
UC-04 & ดูรายการวัตถุที่พบ & ผู้ปฏิบัติงาน & เรียกดูชนิด พิกัด และเวลาของวัตถุที่ระบบตรวจพบ \\
\hline
UC-05 & บันทึกแผนที่ & ผู้ปฏิบัติงาน & สั่งบันทึกแผนที่ปัจจุบันเป็นไฟล์ \\
\hline
UC-06 & สร้างแผนที่ & หุ่นยนต์ & ประมวลผลข้อมูลไลดาร์ร่วมกับโอดอมิเตอร์เพื่อสร้างแผนที่ \\
\hline
UC-07 & ตรวจจับวัตถุ & หุ่นยนต์ & อนุมานภาพจากกล้องด้วยแบบจำลองบนอุปกรณ์เอดจ์ \\
\hline
UC-08 & ระบุพิกัดวัตถุ & หุ่นยนต์ & แปลงผลการตรวจจับเป็นพิกัดบนแผนที่และบันทึกลงไฟล์ข้อมูล \\
\hline
\end{tabularx}
\end{table}

\section{Class Diagram ระดับ Design ของระบบ}

การออกแบบระดับคลาสของระบบสะท้อนโครงสร้างแบบโมดูลของ ROS 2
โดยแต่ละโหนดถูกออกแบบเป็นคลาสที่สืบทอดจากคลาสฐานของ ROS 2
และรับผิดชอบงานเพียงด้านเดียว แผนภาพคลาสแสดงดังรูปที่ \ref{fig:class}

\begin{figure}[htbp]
\centering
\begin{tikzpicture}[
    font=\footnotesize,
    cls/.style={draw, rectangle split, rectangle split parts=2,
                rectangle split part align=left,
                text width=4.6cm, align=left, inner sep=4pt},
    dep/.style={-{Latex[length=2mm]}, dashed},
    use/.style={-{Latex[length=2mm]}}
]

% ── ชั้นควบคุมการเคลื่อนที่ ────────────────────────────────────────
\node[cls] (hw) at (-3.6,0) {
    \textbf{DiffDriveSerial} \nodepart{two}
    - serialPort, baudRate \\ + read() \\ + write() \\ + stopMotors()};

\node[cls] (dc) at (3.6,0) {
    \textbf{diff\_drive\_controller} \nodepart{two}
    - wheelRadius, wheelSeparation \\
    + update() \\ + computeOdometry() \\ + broadcastTF()};

% ── ชั้นรับรู้ ────────────────────────────────────────────────────
\node[cls] (dn) at (-3.6,-4.4) {
    \textbf{DetectorNode} \nodepart{two}
    - trtEngine, confThreshold \\
    + imageCallback() \\ + infer() \\ + track()};

\node[cls] (ol) at (3.6,-4.4) {
    \textbf{ObjectLocalizer} \nodepart{two}
    - cameraIntrinsics \\ - minReturns, maxSpread \\
    + boxToBearingRange() \\ + pickLidarBeams() \\ + toMapFrame()};

% ── ชั้นข้อมูลและแสดงผล ──────────────────────────────────────────
\node[cls] (ls) at (3.6,-8.8) {
    \textbf{LandmarkStore} \nodepart{two}
    - landmarks, mergeRadius \\ + matchOrCreate() \\
    + updateEma() \\ + writeJson()};

\node[cls] (api) at (-3.6,-8.8) {
    \textbf{ApiServer} \nodepart{two}
    + getLandmarks() / getMap() \\ + streamStatus() \\
    + streamCamera() \\ + clearLandmarks()};

% ── ความสัมพันธ์ ─────────────────────────────────────────────────
\draw[use] (hw.east) -- node[midway, above, font=\scriptsize, inner sep=2pt] {encoder} (dc.west);
\draw[use] (dn.east) -- node[midway, above, font=\scriptsize, inner sep=2pt] {detections} (ol.west);
\draw[use] (dc.south) -- node[midway, right, font=\scriptsize, inner sep=3pt] {TF / odom} (ol.north);
\draw[use] (ol.south) -- node[midway, right, font=\scriptsize, inner sep=3pt] {บันทึก} (ls.north);
\draw[use] (ls.west) -- node[midway, above, font=\scriptsize, inner sep=2pt] {query} (api.east);
\draw[dep] (dn.south) -- node[midway, left, font=\scriptsize, inner sep=3pt] {ภาพพร้อมกรอบวัตถุ} (api.north);
\end{tikzpicture}
\caption[แผนภาพคลาสระดับการออกแบบของระบบ (บทที่ 3)]{แผนภาพคลาสระดับการออกแบบของระบบ}
\label{fig:class}
\end{figure}

คลาสหลักในระบบมีดังนี้

\begin{itemize}[leftmargin=1.25cm, itemsep=4pt, topsep=6pt]
    \item \texttt{DiffDriveSerial} เป็น Hardware Interface ของ ros2\_control
    ที่เขียนด้วยภาษา C++ ทำงานบนหน่วยประมวลผลหลัก สื่อสารกับ ESP32 ผ่านพอร์ตอนุกรม
    โดยเขียนคำสั่งความเร็วล้อและอ่านค่าเอนโคเดอร์สะสมกลับมาในทุกรอบการควบคุม

    \item \texttt{diff\_drive\_controller} เป็นคอนโทรลเลอร์มาตรฐานของ ros2\_control
    รับคำสั่งความเร็วเชิงเส้นและเชิงมุม แปลงเป็นความเร็วของล้อทั้งสองข้าง
    คำนวณโอดอมิเตอร์จากค่าเอนโคเดอร์
    แล้วเผยแพร่ข้อความโอดอมิเตอร์พร้อมกระจายการแปลงพิกัดจาก \texttt{odom}
    ไปยัง \texttt{base\_link}

    \item \texttt{DetectorNode} รับภาพจากกล้อง
    ส่งเข้าแบบจำลองที่ทำงานบน TensorRT
    แล้วเผยแพร่ผลการตรวจจับในรูปกรอบสี่เหลี่ยม ชื่อคลาส ค่าความเชื่อมั่น
    และหมายเลขติดตามวัตถุ (Track ID)

    \item \texttt{ObjectLocalizer} รับผลการตรวจจับร่วมกับข้อมูลไลดาร์และการแปลงพิกัดจาก TF2
    คำนวณพิกัดของวัตถุบนแผนที่ แล้วส่งต่อให้ชั้นจัดเก็บข้อมูล

    \item \texttt{LandmarkStore} เก็บรายการวัตถุที่ระบบยืนยันแล้วไว้ในหน่วยความจำ
    รับผิดชอบการจับคู่ว่าผลการตรวจจับใหม่เป็นวัตถุเดิมหรือวัตถุใหม่
    การปรับตำแหน่งด้วยค่าเฉลี่ยเคลื่อนที่แบบถ่วงน้ำหนักเลขชี้กำลัง
    (Exponential Moving Average) ที่ค่าคงที่ $\alpha$ = 0.3
    และการเขียนรายการทั้งหมดลงไฟล์ JSON

    \item \texttt{ApiServer} เปิดปลายทาง REST API และช่องทาง WebSocket
    เป็นสะพานเชื่อมระหว่างระบบ ROS 2 กับส่วนติดต่อผู้ใช้บนเบราว์เซอร์
\end{itemize}

% ---------------------------------------------------------------------
\chapter{รายละเอียดการทำงานของระบบ}

บทนี้อธิบายรายละเอียดการทำงานจริงของแต่ละโมดูล
โดยเรียงตามลำดับการพัฒนาและการพึ่งพากันของระบบ
เริ่มจากชั้นควบคุมการเคลื่อนที่ซึ่งเป็นฐานของทุกส่วน
ไปจนถึงชั้นแสดงผลบนเว็บแอปพลิเคชัน

\section{ระบบควบคุมการเคลื่อนที่และการคำนวณ Odometry}

\subsection{สถาปัตยกรรมของชั้นควบคุมการเคลื่อนที่}

การควบคุมมอเตอร์และการอ่านค่าเอนโคเดอร์ถูกแยกออกไปทำงานบนไมโครคอนโทรลเลอร์ ESP32
แทนที่จะทำบนหน่วยประมวลผลหลัก เหตุผลคืองานเหล่านี้ต้องการการตอบสนองที่แน่นอนตามเวลา
ซึ่งระบบปฏิบัติการทั่วไปที่ไม่ใช่ระบบเรียลไทม์ไม่รับประกันให้

ESP32 รันเฟิร์มแวร์ของตนเองที่ไม่ได้ใช้ ROS 2
และสื่อสารกับหน่วยประมวลผลหลักผ่านพอร์ตอนุกรมยูเอสบีที่อัตราบอด 57600
ด้วยโพรโทคอลข้อความอย่างง่าย โดยรับคำสั่งความเร็วของล้อทั้งสองข้าง
และตอบกลับด้วยค่าเอนโคเดอร์สะสมของแต่ละข้าง

ฝั่งหน่วยประมวลผลหลัก Hardware Interface ของ ros2\_control ชื่อ \texttt{DiffDriveSerial}
ทำหน้าที่อ่านและเขียนพอร์ตอนุกรมนี้ในทุกรอบการควบคุม (30 ครั้งต่อวินาที)
แล้วส่งต่อค่าที่อ่านได้ให้คอนโทรลเลอร์ \texttt{diff\_cont} นำไปคำนวณโอดอมิเตอร์
ด้วยเหตุนี้ระบบจึงไม่มีหัวข้อจำนวนพัลส์ปรากฏใน ROS 2
เพราะค่าเอนโคเดอร์ถูกอ่านผ่านชั้น Hardware Interface โดยตรง ไม่ได้ส่งผ่านหัวข้อ

เฟิร์มแวร์มีกลไกเฝ้าระวัง (Watchdog) ที่จะสั่งหยุดมอเตอร์ทันที
หากไม่ได้รับคำสั่งความเร็วใหม่ภายในช่วงเวลาที่กำหนด
เพื่อป้องกันไม่ให้หุ่นยนต์วิ่งต่อไปเมื่อการสื่อสารขาดหาย

การต่อสัญญาณระหว่างไมโครคอนโทรลเลอร์กับชุดขับเคลื่อนแสดงดังรูปที่ \ref{fig:wiring-motor}

\begin{figure}[htbp]
\centering
\includegraphics[width=\textwidth]{images/motor_wiring_diagram.png}
\caption[การต่อสัญญาณระหว่าง ESP32 โมดูลขับมอเตอร์ และเอนโคเดอร์ (บทที่ 4)]{การต่อสัญญาณระหว่าง ESP32 โมดูลขับมอเตอร์ และเอนโคเดอร์}
\label{fig:wiring-motor}
\end{figure}

ส่วนประกอบเชิงกลของชุดขับเคลื่อนที่ติดตั้งจริงแสดงดังรูปที่ \ref{fig:drivetrain-photo}

\begin{figure}[htbp]
\centering
\includegraphics[height=10cm,keepaspectratio]{images/drivetrain_photo.jpg}
\caption[ชุดขับเคลื่อนและการติดตั้งเอนโคเดอร์บนหุ่นยนต์ต้นแบบ (บทที่ 4)]{ชุดขับเคลื่อนและการติดตั้งเอนโคเดอร์บนหุ่นยนต์ต้นแบบ}
\label{fig:drivetrain-photo}
\end{figure}

\subsection{เฟิร์มแวร์บนไมโครคอนโทรลเลอร์และโพรโทคอลสื่อสาร}

เฟิร์มแวร์บน ESP32 เขียนขึ้นใหม่ทั้งหมดแทนการดัดแปลงโค้ดเดิมที่เขียนไว้สำหรับชิปตระกูล AVR
เนื่องจากโค้ดเดิมอ้างอิงรีจิสเตอร์และชุดคำสั่งขัดจังหวะเฉพาะของ AVR ซึ่งไม่มีอยู่บน ESP32
การทำงานหลักของเฟิร์มแวร์มีสามส่วน

\begin{itemize}[leftmargin=1.25cm, itemsep=6pt, topsep=6pt]
    \item \textbf{การอ่านค่าเอนโคเดอร์ด้วยวงจรนับพัลส์ในตัวชิป}
    ESP32 มีวงจรนับพัลส์ (Pulse Counter) ที่ถอดรหัสสัญญาณควอเดรเจอร์แบบสี่เท่าได้ในระดับฮาร์ดแวร์
    พร้อมตัวกรองสัญญาณรบกวนที่ตั้งไว้ที่ 1 ไมโครวินาที
    วิธีนี้ไม่ใช้เวลาประมวลผลของหน่วยประมวลผลกลางเลย และไม่พลาดการนับเมื่อล้อหมุนเร็ว
    ซึ่งต่างจากการใช้ชุดคำสั่งขัดจังหวะที่อาจนับตกหล่นได้
    เนื่องจากวงจรนับเป็นแบบ 16 บิต เฟิร์มแวร์จึงตั้งจุดเฝ้าระวังที่ขอบทั้งสองด้าน
    แล้วขยายค่าที่นับได้ให้กว้างขึ้นด้วยซอฟต์แวร์

    \item \textbf{การสร้างสัญญาณควบคุมความเร็วด้วยตัวควบคุมเฉพาะของชิป}
    เฟิร์มแวร์ใช้ชุดคำสั่งควบคุมสัญญาณ (LEDC) ของ ESP32 โดยตรง
    ที่ความถี่ 20 กิโลเฮิรตซ์ ความละเอียด 8 บิต
    แทนฟังก์ชันมาตรฐานของไลบรารี Arduino ที่ทำงานที่ 1 กิโลเฮิรตซ์
    เหตุผลคือความถี่ 1 กิโลเฮิรตซ์อยู่ในย่านที่หูได้ยินและทำให้มอเตอร์ส่งเสียงหวีด
    ขณะที่ 20 กิโลเฮิรตซ์อยู่เหนือย่านการได้ยินแต่ยังต่ำพอที่โมดูลขับมอเตอร์จะรับไหว
    สัญญาณควบคุมถูกป้อนเข้าที่ขากำหนดทิศทางโดยตรึงขาเปิดใช้งานไว้ที่แรงดันสูง
    ซึ่งตรงกับวิธีเดินสายบนหุ่นยนต์ต้นแบบ

    \item \textbf{ตัวควบคุมแบบ PID และกลไกเฝ้าระวัง} เฟิร์มแวร์รันรอบควบคุม PID
    ที่อัตรา 30 ครั้งต่อวินาที ซึ่งตรงกับอัตรารอบการควบคุมของฝั่งหน่วยประมวลผลหลัก
    และสั่งหยุดมอเตอร์ทันทีหากไม่ได้รับคำสั่งความเร็วใหม่ภายใน 2 วินาที
\end{itemize}

โพรโทคอลสื่อสารเป็นข้อความตัวอักษรเดียวตามด้วยตัวแปรคั่นด้วยช่องว่าง
ปิดท้ายด้วยอักขระขึ้นบรรทัดใหม่ ที่อัตราบอด 57600 รูปแบบ 8N1
รายละเอียดคำสั่งทั้งหมดแสดงดังตารางที่ \ref{tab:serialproto}

\begin{table}[htbp]
\centering
\caption[คำสั่งในโพรโทคอลสื่อสารระหว่างหน่วยประมวลผลหลักกับ ESP32 (บทที่ 4)]{คำสั่งในโพรโทคอลสื่อสารระหว่างหน่วยประมวลผลหลักกับ ESP32}
\label{tab:serialproto}
\small
\begin{tabularx}{\textwidth}{|c|>{\raggedright\arraybackslash}p{3.4cm}|X|}
\hline
\textbf{คำสั่ง} & \textbf{ตัวแปร} & \textbf{ความหมาย} \\
\hline
\texttt{e} & --- & อ่านค่าเอนโคเดอร์สะสมของล้อซ้ายและล้อขวา \\
\hline
\texttt{r} & --- & รีเซ็ตค่าเอนโคเดอร์สะสมทั้งสองข้างเป็นศูนย์ \\
\hline
\texttt{o} & ค่าสัญญาณซ้าย ขวา & สั่งค่าสัญญาณควบคุมความเร็วดิบ ($-$255 ถึง 255) ข้ามตัวควบคุม PID
             ใช้สำหรับการตรวจสอบฮาร์ดแวร์ \\
\hline
\texttt{m} & จำนวนพัลส์ซ้าย ขวา & สั่งความเร็วเป้าหมายเป็นจำนวนพัลส์ต่อรอบการควบคุม
             ผ่านตัวควบคุม PID \\
\hline
\texttt{u} & $K_p$:$K_d$:$K_i$:$K_o$ & ปรับค่าพารามิเตอร์ของตัวควบคุม PID \\
\hline
\texttt{i} & --- & อ่านค่าดิบของมาตรความเร่งสามแกนและไจโรสโคปสามแกนจากหน่วยวัดความเฉื่อย \\
\hline
\end{tabularx}
\end{table}


\subsection{การแปลงคำสั่งความเร็วเป็นความเร็วล้อ}

คำสั่งที่ระบบนำทางส่งมาอยู่ในรูปความเร็วเชิงเส้น $v$ และความเร็วเชิงมุม $\omega$
ซึ่งต้องแปลงเป็นความเร็วของล้อซ้ายและล้อขวาตามแบบจำลองการขับเคลื่อนแบบดิฟเฟอเรนเชียล ดังนี้

\begin{equation}
v_L = v - \frac{\omega L}{2}, \qquad v_R = v + \frac{\omega L}{2}
\end{equation}

ความสัมพันธ์เชิงเรขาคณิตของสมการข้างต้นแสดงดังรูปที่ \ref{fig:diffdrive-kin}

\begin{figure}[htbp]
\centering
\includegraphics[width=0.9\textwidth]{images/diff_drive_kinematics.png}
\caption[แบบจำลองจลนศาสตร์ของการขับเคลื่อนแบบดิฟเฟอเรนเชียล (บทที่ 4)]{แบบจำลองจลนศาสตร์ของการขับเคลื่อนแบบดิฟเฟอเรนเชียล}
\label{fig:diffdrive-kin}
\end{figure}

โดยที่ $L$ คือระยะห่างระหว่างล้อทั้งสอง (Track Width)
ซึ่งจากการสอบเทียบของหุ่นยนต์ต้นแบบมีค่าเท่ากับ \textbf{0.25168 เมตร}
ค่านี้ไม่ได้มาจากการวัดด้วยไม้บรรทัดโดยตรง
เพราะจุดสัมผัสจริงของล้อกับพื้นไม่ใช่จุดกึ่งกลางของหน้ายาง
แต่ได้จากการสอบเทียบย้อนกลับจากมุมที่หุ่นยนต์หมุนได้จริงเทียบกับมุมที่โอดอมิเตอร์รายงาน
และต่อมาได้รับการยืนยันอีกครั้งจากข้อมูลเอนโคเดอร์ขณะหมุนอยู่กับที่
ซึ่งให้ค่า \textbf{0.25169 เมตร} ต่างจากค่าที่ตั้งไว้เพียงร้อยละ 0.003

ผลการสอบเทียบย้อนกลับที่ใช้กำหนดค่าระยะห่างระหว่างล้อแสดงดังรูปที่ \ref{fig:calib-track}

\begin{figure}[htbp]
\centering
\includegraphics[width=\textwidth]{images/calibration_track_width.png}
\caption[ผลการสอบเทียบระยะห่างระหว่างล้อของหุ่นยนต์ต้นแบบ (บทที่ 4)]{ผลการสอบเทียบระยะห่างระหว่างล้อของหุ่นยนต์ต้นแบบ}
\label{fig:calib-track}
\end{figure}

\subsection{การคำนวณ Odometry}

คอนโทรลเลอร์ \texttt{diff\_cont} บนหน่วยประมวลผลหลักทำหน้าที่แปลงจำนวนพัลส์สะสม
ที่อ่านผ่านชั้น Hardware Interface ให้เป็นตำแหน่งและทิศทางของหุ่นยนต์
ระบบไม่มีโหนดคำนวณโอดอมิเตอร์แยกต่างหาก
พารามิเตอร์ของหุ่นยนต์ต้นแบบที่ได้จากการวัดและสอบเทียบ
แสดงดังตารางที่ \ref{tab:robotparam}

\begin{table}[htbp]
\centering
\caption[พารามิเตอร์เชิงกลของหุ่นยนต์ต้นแบบ (บทที่ 4)]{พารามิเตอร์เชิงกลของหุ่นยนต์ต้นแบบ}
\label{tab:robotparam}
\begin{tabularx}{\textwidth}{|X|c|c|}
\hline
\textbf{พารามิเตอร์} & \textbf{สัญลักษณ์} & \textbf{ค่าที่ใช้} \\
\hline
จำนวนพัลส์ต่อรอบของล้อซ้าย & $N_L$ & 2475 พัลส์ต่อรอบ \\
\hline
จำนวนพัลส์ต่อรอบของล้อขวา & $N_R$ & 2470 พัลส์ต่อรอบ \\
\hline
เส้นผ่านศูนย์กลางล้อประสิทธิผล (ขณะรับน้ำหนัก) & $D$ & 0.0654 เมตร \\
\hline
เส้นผ่านศูนย์กลางล้อขณะไม่รับน้ำหนัก & --- & 0.068 เมตร \\
\hline
ระยะห่างระหว่างล้อ & $L$ & 0.25168 เมตร \\
\hline
รัศมีล้อที่ตั้งค่าในคอนโทรลเลอร์ & $r$ & 0.0327 เมตร \\
\hline
ตัวคูณแก้รัศมีล้อซ้าย / ขวา & --- & 1.002982 / 0.997018 \\
\hline
อัตรารอบการควบคุม (Control Loop) & --- & 30 ครั้งต่อวินาที \\
\hline
อัตราการเผยแพร่ข้อความโอดอมิเตอร์ & --- & ตั้งค่าไว้ 50 วัดได้จริง 30 ครั้งต่อวินาที \\
\hline
ขีดจำกัดความเร็วเชิงเส้นของคอนโทรลเลอร์ & --- & $\pm$0.30 เมตรต่อวินาที \\
\hline
ขีดจำกัดความเร็วเชิงมุมของคอนโทรลเลอร์ & --- & $\pm$0.5 เรเดียนต่อวินาที \\
\hline
\end{tabularx}
\end{table}

ค่าสองค่าในตารางนี้ต้องอธิบายเพิ่มเติมเพื่อไม่ให้เข้าใจผิด
ค่าแรกคืออัตราการเผยแพร่ข้อความโอดอมิเตอร์ ซึ่งแม้จะตั้งไว้ที่ 50 ครั้งต่อวินาที
แต่คอนโทรลเลอร์ไม่สามารถเผยแพร่ได้เร็วกว่าอัตรารอบการทำงานของตัวจัดการคอนโทรลเลอร์
อัตราที่วัดได้จริงบนระบบจึงเป็น 30 ครั้งต่อวินาทีเท่ากับอัตรารอบการควบคุม
ค่าที่สองคือเส้นผ่านศูนย์กลางล้อ ซึ่งค่าที่ใช้ในการคำนวณไม่ใช่ค่าที่วัดด้วยเวอร์เนียร์
แต่เป็นค่าประสิทธิผลที่ได้จากการสอบเทียบระยะทางจริงเทียบกับระยะทางที่โอดอมิเตอร์รายงาน
ซึ่งเล็กกว่าเส้นผ่านศูนย์กลางขณะไม่รับน้ำหนักประมาณร้อยละ 4
เนื่องจากยางยุบตัวเมื่อรับน้ำหนักของหุ่นยนต์

ระยะทางที่ล้อเคลื่อนที่ได้ต่อหนึ่งพัลส์คำนวณได้จาก

\begin{equation}
d_{\text{tick}} = \frac{\pi D}{N}
\end{equation}

เมื่อแทนค่าจะได้ระยะทางต่อพัลส์ประมาณ 8.30$\times$10$^{-5}$ เมตร

อนึ่ง จำนวนพัลส์ต่อรอบของล้อทั้งสองข้างที่ต่างกันเล็กน้อยนั้น
ไม่ได้สะท้อนความแตกต่างของตัวเอนโคเดอร์หรืออัตราทดเกียร์
ซึ่งเป็นชิ้นส่วนรุ่นเดียวกันทั้งสองข้าง
แต่เป็นผลของการสอบเทียบที่ดูดซับความต่างของเส้นผ่านศูนย์กลางยางจริงเข้าไว้
ระบบจึงแยกการแก้ความต่างของรัศมีล้อออกมาเป็นตัวคูณต่างหากตามที่แสดงในตาราง

จากนั้นระยะทางที่ล้อแต่ละข้างเคลื่อนที่ในช่วงเวลาหนึ่งคือ

\begin{equation}
\Delta s_L = d_{\text{tick}} \cdot \Delta n_L, \qquad
\Delta s_R = d_{\text{tick}} \cdot \Delta n_R
\end{equation}

โดย $\Delta n_L$ และ $\Delta n_R$ คือผลต่างของจำนวนพัลส์สะสมระหว่างสองรอบการคำนวณ
ระยะทางและมุมที่ตัวหุ่นยนต์เปลี่ยนไปคำนวณได้จาก

\begin{equation}
\Delta s = \frac{\Delta s_R + \Delta s_L}{2}, \qquad
\Delta \theta = \frac{\Delta s_R - \Delta s_L}{L}
\end{equation}

และปรับตำแหน่งสะสมโดยใช้มุมเฉลี่ยของช่วงเวลานั้นเพื่อลดความคลาดเคลื่อนขณะเลี้ยว

\begin{equation}
\begin{aligned}
x_{k+1} &= x_k + \Delta s \cos\left(\theta_k + \frac{\Delta \theta}{2}\right) \\
y_{k+1} &= y_k + \Delta s \sin\left(\theta_k + \frac{\Delta \theta}{2}\right) \\
\theta_{k+1} &= \theta_k + \Delta \theta
\end{aligned}
\end{equation}

ผลลัพธ์ถูกเผยแพร่เป็นข้อความ \texttt{nav\_msgs/Odometry}
บนหัวข้อ \texttt{/diff\_cont/odom}
(ซึ่งเป็นหัวข้อตามชื่ออินสแตนซ์ของคอนโทรลเลอร์ ไม่ใช่ \texttt{/odom})
พร้อมกับกระจายการแปลงพิกัดจากเฟรม \texttt{odom} ไปยัง \texttt{base\_link} ผ่าน TF2

การตรวจสอบความถูกต้องของโอดอมิเตอร์ทำเป็นสองขั้น
ขั้นแรกคือการเข็นหุ่นยนต์ด้วยมือโดยไม่สั่งงานมอเตอร์
เพื่อให้ความผิดพลาดของทิศทางการนับเอนโคเดอร์ไม่ถูกตัวควบคุม PID กลบไป
ผลคือการเข็นตรงเป็นระยะประมาณ 1 เมตรให้ระยะทางในแนวเส้นตรง 0.980 เมตร
โดยมุมเปลี่ยนไปเพียง 0.1 องศา และการหมุนอยู่กับที่ประมาณ 90 องศา
ให้การเลื่อนตำแหน่งเพียง 0.9 เซนติเมตร
การที่การเคลื่อนที่ตรงไม่ทำให้เกิดมุม และการหมุนไม่ทำให้เกิดการเลื่อนตำแหน่ง
คือหลักฐานว่าทิศทางการนับของเอนโคเดอร์ทั้งสองข้างถูกต้องและสมการจลนศาสตร์ไม่สลับข้าง

ขั้นที่สองคือการสอบเทียบความไม่สมมาตรของล้อทั้งสองข้าง
ด้วยการเดินตรงระยะ 3 เมตรแบบวงปิดที่ความเร็ว 0.10 เมตรต่อวินาที
แล้วเปรียบเทียบตำแหน่งปลายทางจริงบนพื้นกับค่าที่โอดอมิเตอร์รายงาน
ผลคือโอดอมิเตอร์รายงานการเบี่ยงออกด้านข้างเพียง 0.9 มิลลิเมตร
ขณะที่หุ่นยนต์เบี่ยงไปทางขวาจริง 10.9 เซนติเมตร
ช่องว่างนี้คือความลำเอียงของโอดอมิเตอร์ล้วน ๆ คิดเป็นความคลาดเคลื่อนของทิศทาง 4.07 องศา
ตลอดระยะ 3 เมตร หรือความไม่สมมาตรของรัศมีล้อสองข้างร้อยละ 0.60
เมื่อนำค่านี้ไปใส่เป็นตัวคูณแก้รัศมีล้อตามที่แสดงในตารางที่ \ref{tab:robotparam}
แล้วทดสอบซ้ำในเงื่อนไขเดียวกัน ความคลาดเคลื่อนของทิศทางลดลงเหลือประมาณ 0.1 องศา
ตรงกับค่าที่คำนวณไว้ล่วงหน้า
ข้อสังเกตเชิงวิธีการคือค่าแก้นี้ใส่ไว้ที่ตัวคูณรัศมีล้อ ไม่ใช่ที่จำนวนพัลส์ต่อรอบ
เพราะจำนวนพัลส์ต่อรอบเป็นคุณสมบัติของจานเอนโคเดอร์และชุดเกียร์ซึ่งเป็นชิ้นส่วนรุ่นเดียวกันทั้งสองข้าง
มีเพียงเส้นผ่านศูนย์กลางของยางเท่านั้นที่ต่างกันได้

ผลการทดสอบความถูกต้องของโอดอมิเตอร์ที่คำนวณได้แสดงดังรูปที่ \ref{fig:odom-test}

\begin{figure}[htbp]
\centering
\includegraphics[width=\textwidth]{images/odometry_test_path.png}
\caption[ผลการทดสอบความถูกต้องของการคำนวณโอดอมิเตอร์ (บทที่ 4)]{ผลการทดสอบความถูกต้องของการคำนวณโอดอมิเตอร์}
\label{fig:odom-test}
\end{figure}

\subsection{การผสานข้อมูลไจโรสโคปกับโอดอมิเตอร์ด้วยตัวกรองคาลมานแบบขยาย}

โอดอมิเตอร์จากล้อเพียงอย่างเดียวมีข้อจำกัดที่เห็นได้ชัดเมื่อหุ่นยนต์ติดขัดและล้อหมุนฟรี
เอนโคเดอร์จะรายงานการหมุนที่ตัวหุ่นยนต์ไม่ได้หมุนจริง
ระบบจึงติดตั้งหน่วยวัดความเฉื่อย GY-521 (MPU6050) เพิ่มเติม
และผสานอัตราการหมุนรอบแกนตั้งเข้ากับโอดอมิเตอร์ด้วยตัวกรองคาลมานแบบขยาย
ของแพ็กเกจ robot\_localization~\cite{robotlocalization}
โครงสร้างการผสานข้อมูลแสดงดังรูปที่ \ref{fig:ekf-block}

\begin{figure}[htbp]
\centering
\includegraphics[height=12cm,keepaspectratio]{images/ekf_block_diagram.png}
\caption[โครงสร้างการผสานข้อมูลโอดอมิเตอร์กับไจโรสโคป (บทที่ 4)]{โครงสร้างการผสานข้อมูลโอดอมิเตอร์กับไจโรสโคป}
\label{fig:ekf-block}
\end{figure}

ประเด็นสำคัญของการออกแบบมีดังนี้

\begin{itemize}[leftmargin=1.25cm, itemsep=6pt, topsep=6pt]
    \item \textbf{หน่วยวัดความเฉื่อยเชื่อมต่อผ่านชั้น Hardware Interface เดิม}
    เนื่องจากกระบวนการของ ros2\_control ยึดครองพอร์ตอนุกรมของ ESP32 ไว้แต่เพียงผู้เดียว
    จึงไม่สามารถสร้างโหนดแยกเพื่ออ่านค่าหน่วยวัดความเฉื่อยได้
    ชั้น Hardware Interface จึงส่งคำสั่งอ่านค่าต่อท้ายการอ่านเอนโคเดอร์ในทุกรอบการควบคุม
    แล้วประกาศค่าที่ได้เป็นเซนเซอร์ของ ros2\_control
    ให้ตัวกระจายข้อมูลเซนเซอร์เผยแพร่ต่อในระบบ

    \item \textbf{การแบ่งน้ำหนักระหว่างสองแหล่งข้อมูล}
    ล้อให้ความเร็วเชิงเส้นและอัตราการหมุนที่ค่าความแปรปรวน 1$\times$10$^{-2}$
    ส่วนไจโรสโคปให้อัตราการหมุนที่ค่าความแปรปรวน 5$\times$10$^{-5}$
    ทำให้ในกรณีที่ล้อลื่นไถล ไจโรสโคปมีอิทธิพลเหนือกว่าล้อประมาณ 200 เท่า
    ขณะที่หากหน่วยวัดความเฉื่อยขัดข้อง ค่าประมาณยังคงดำเนินต่อได้จากล้อ
    ทั้งนี้ระบบไม่นำตำแหน่งจากล้อเข้าสู่ตัวกรอง
    การลื่นไถลจึงไม่เข้าสู่ระบบในรูปของการกระโดดของตำแหน่ง

    \item \textbf{การย้ายผู้กระจายการแปลงพิกัด}
    เมื่อเปิดใช้งานตัวกรอง คอนโทรลเลอร์ \texttt{diff\_cont}
    จะหยุดกระจายการแปลงพิกัดจาก \texttt{odom} ไปยัง \texttt{base\_link}
    และให้ตัวกรองทำหน้าที่นั้นแทน
    เหตุผลคือหากมีผู้กระจายสองรายบนเส้นเชื่อมเดียวกัน
    ระบบพิกัดจะสั่นไปมาระหว่างสองค่าและแสดงอาการเหมือนสัญญาณรบกวนของโอดอมิเตอร์
    ทั้งที่ไม่ใช่

    \item \textbf{การตั้งค่าตัวกรองสัญญาณของชิปตามอัตราการอ่าน}
    หน่วยประมวลผลหลักอ่านค่าหนึ่งตัวอย่างต่อรอบการควบคุมที่ 30 ครั้งต่อวินาที
    ความถี่สูงสุดที่แทนค่าได้จึงเป็น 15 เฮิรตซ์
    สัญญาณที่สูงกว่านั้นจะพับกลับมาปรากฏเป็นสัญญาณรบกวน
    ระบบจึงตั้งตัวกรองความถี่ต่ำผ่านในตัวชิปไว้ที่แถบความถี่ 20 เฮิรตซ์
    การเปลี่ยนค่านี้จากแถบ 42 เฮิรตซ์มาเป็น 20 เฮิรตซ์
    ทำให้ความแปรปรวนของไจโรสโคปขณะระบบทำงานจริงลดลงประมาณ 4.8 เท่า
    ซึ่งเป็นลักษณะเฉพาะของการพับกลับของสัญญาณโดยตรง
\end{itemize}

ค่าเบี่ยงเบนศูนย์ของไจโรสโคปถูกวัดขณะหุ่นยนต์หยุดนิ่งและนำไปหักออกในชั้น Hardware Interface
โดยวัดจากข้อมูล 10 วินาทีขณะหุ่นยนต์วางนิ่งบนพื้นราบและมอเตอร์ไม่ทำงาน
ผลการวัดทั้งหมดของส่วนนี้สรุปดังตารางที่ \ref{tab:imuresult}

\begin{table}[htbp]
\centering
\caption[ผลการวัดของหน่วยวัดความเฉื่อยและตัวกรองคาลมานแบบขยาย (บทที่ 4)]{ผลการวัดของหน่วยวัดความเฉื่อยและตัวกรองคาลมานแบบขยาย}
\label{tab:imuresult}
\small
\begin{tabularx}{\textwidth}{|X|c|}
\hline
\textbf{รายการ} & \textbf{ค่าที่วัดได้} \\
\hline
ค่าเบี่ยงเบนศูนย์ของไจโรสโคปรอบแกนตั้ง & $-$0.62 องศาต่อวินาที \\
\hline
อัตราการหมุนเฉลี่ยหลังหักค่าเบี่ยงเบนศูนย์ (ขณะหยุดนิ่ง) & $-$0.009 องศาต่อวินาที \\
\hline
ส่วนเบี่ยงเบนมาตรฐานของอัตราการหมุน ขณะมอเตอร์ไม่มีไฟเลี้ยง & 0.00114 เรเดียนต่อวินาที \\
\hline
ส่วนเบี่ยงเบนมาตรฐาน ขณะระบบทำงานจริง ที่ตัวกรองในชิป 42 เฮิรตซ์ & 0.01053 เรเดียนต่อวินาที \\
\hline
ส่วนเบี่ยงเบนมาตรฐาน ขณะระบบทำงานจริง ที่ตัวกรองในชิป 20 เฮิรตซ์ & 0.00486 เรเดียนต่อวินาที \\
\hline
การเลื่อนของทิศทางสุทธิขณะหุ่นยนต์หยุดนิ่ง & $+$0.01 องศาต่อนาที \\
\hline
อัตราการเผยแพร่ข้อมูลหน่วยวัดความเฉื่อย & 30.00 ครั้งต่อวินาที \\
\hline
อัตราการเผยแพร่สถานะข้อต่อเมื่อเปิดการอ่านหน่วยวัดความเฉื่อย & 29.996 ครั้งต่อวินาที \\
\hline
เวลาที่ใช้ในการอ่านหน่วยวัดความเฉื่อยหนึ่งครั้ง (ค่ามัธยฐาน) & 6.6 มิลลิวินาที \\
\hline
\end{tabularx}
\end{table}

ประเด็นที่ต้องอ่านจากตารางนี้มีสองข้อ
ข้อแรกคือ \textbf{ค่าสัญญาณรบกวนที่วัดขณะระบบไม่ทำงานใช้แทนค่าขณะระบบทำงานจริงไม่ได้}
ชิปตัวเดียวกันในห้องเดียวกันให้ความแปรปรวนต่างกันถึง 85 เท่าระหว่างสองเงื่อนไขนี้
การตั้งค่าความแปรปรวนจากค่าขณะหยุดนิ่งทำให้ตัวกรองเชื่อไจโรสโคปมากเกินจริง
และไม่ได้ทำการกรองสัญญาณเลย ค่าที่ติดตั้งจริงจึงมาจากการวัดขณะระบบทำงาน
ข้อที่สองคือการอ่านหน่วยวัดความเฉื่อยใช้เวลา 6.6 มิลลิวินาทีจากงบประมาณเวลา 33.3 มิลลิวินาที
ต่อหนึ่งรอบการควบคุม ซึ่งเมื่อรวมกับการอ่านเอนโคเดอร์และการสั่งมอเตอร์แล้วยังพอดีกับรอบ
ยืนยันได้จากอัตราการเผยแพร่สถานะข้อต่อที่ยังคงอยู่ที่ 30 ครั้งต่อวินาที

การตรวจสอบความถูกต้องของการผสานข้อมูลทำโดยการเข็นหุ่นยนต์ด้วยมือประมาณ 1.1 เมตร
และหมุนประมาณ 90 องศา แล้วเปรียบเทียบผลของโอดอมิเตอร์จากล้อกับผลของตัวกรอง
ได้ระยะทาง 1.127 เมตรเทียบกับ 1.116 เมตร ต่างกันร้อยละ 0.98
และมุมที่เปลี่ยนไป $-$87.9 องศาเทียบกับ $-$92.2 องศา
\textbf{ผลที่สำคัญที่สุดของการทดสอบนี้คือเครื่องหมายของมุมตรงกันทั้งสองแหล่ง}
เพราะการผสานไจโรสโคปที่มีเครื่องหมายกลับด้านจะไม่แสดงอาการผิดปกติใด ๆ ในระบบพิกัด
แต่ทำให้ค่าประมาณทิศทางแย่กว่าการไม่ผสานข้อมูลเลย
ส่วนมุมที่ต่างกัน 4.3 องศานั้นยังไม่สามารถระบุได้ว่าเป็นความคลาดเคลื่อนของแหล่งใด
เพราะมุมจริงของการหมุนด้วยมือไม่ทราบค่าแน่นอน
และค่าความคลาดเคลื่อนของตัวคูณสเกลของไจโรสโคปตามสเปกที่ร้อยละ 3 ก็ครอบคลุมส่วนต่างนี้อยู่แล้ว

ผลการวัดหลังการหักค่าเบี่ยงเบนศูนย์และการยืนยันการทำงานแสดงดังรูปที่ \ref{fig:imu-drift}

\begin{figure}[htbp]
\centering
\includegraphics[width=\textwidth]{images/imu_yaw_drift.png}
\caption[การเลื่อนของทิศทางขณะหุ่นยนต์หยุดนิ่ง (บทที่ 4)]{การเลื่อนของทิศทางขณะหุ่นยนต์หยุดนิ่ง}
\label{fig:imu-drift}
\end{figure}

\textbf{ข้อจำกัดที่ต้องระบุไว้อย่างตรงไปตรงมาคือ ขณะหุ่นยนต์หยุดนิ่งสนิท
ตัวกรองให้ผลแย่กว่าโอดอมิเตอร์จากล้อ} เนื่องจากเอนโคเดอร์ที่ไม่หมุนย่อมรายงานว่าไม่มีการหมุน
ซึ่งเป็นคำตอบที่ถูกต้องพอดี ขณะที่ไจโรสโคปยังคงมีการเดินสุ่มของค่าที่อ่านได้
เมื่อวัดในช่วงเวลา 30 วินาทีขณะหุ่นยนต์หยุดนิ่ง
ผลรวมของมุมที่เปลี่ยนไปตลอดช่วงของโอดอมิเตอร์จากล้อเท่ากับ 0.000 องศาพอดี
ขณะที่ของตัวกรองเท่ากับ 5.68 องศา แม้ว่าการเลื่อนสุทธิของตัวกรองจะเหลือเพียง 0.01 องศาต่อนาทีก็ตาม
ประโยชน์ทั้งหมดของการผสานข้อมูลนี้จึงอยู่ที่กรณีหุ่นยนต์กำลังเคลื่อนที่และกรณีล้อลื่นไถล
และด้วยเหตุนี้ระบบจึงตั้งค่าให้การผสานข้อมูลเป็นตัวเลือกเสริมที่ปิดไว้เป็นค่าเริ่มต้น
เพื่อให้ชุดการทำงานที่ผ่านการทดสอบแล้วยังคงเริ่มทำงานได้เหมือนเดิมเสมอ


\section{ระบบสร้างแผนที่และระบุตำแหน่งด้วย SLAM}

\subsection{การกำหนดโครงสร้างระบบพิกัด}

ก่อนที่ slam\_toolbox จะทำงานได้ ระบบต้องมีต้นไม้ของระบบพิกัดที่สมบูรณ์
โดยเฉพาะการแปลงพิกัดจาก \texttt{base\_link} ไปยัง \texttt{laser\_frame} (ผ่าน \texttt{chassis})
ซึ่งเป็นค่าคงที่ที่ได้จากการวัดตำแหน่งติดตั้งไลดาร์บนตัวหุ่นยนต์
ค่านี้ถูกประกาศไว้ในไฟล์ URDF ที่บรรยายโครงสร้างเชิงกลของหุ่นยนต์
เมื่อระบบทำงานครบทุกส่วน เส้นทางหลักของต้นไม้ระบบพิกัดมีลำดับดังนี้

\begin{center}
\texttt{map} $\rightarrow$ \texttt{odom} $\rightarrow$ \texttt{base\_link}
$\rightarrow$ \texttt{chassis}
$\rightarrow$ \{\texttt{laser\_frame}, \texttt{camera\_link}, \texttt{imu\_link}\}
\end{center}

ส่วน \texttt{base\_footprint} เป็นกรอบพิกัดลูกของ \texttt{base\_link} ที่ระดับพื้น
และล้อทั้งสองเป็นกรอบพิกัดลูกที่ขยับตามค่าเอนโคเดอร์
ต้นไม้ของระบบพิกัดที่ตรวจสอบได้จากระบบจริงแสดงดังรูปที่ \ref{fig:tftree}
โดย slam\_toolbox เผยแพร่ \texttt{map}$\rightarrow$\texttt{odom} ที่ประมาณ 12 ครั้งต่อวินาทีขณะเปิดครบทุกส่วน
(ตั้งค่าไว้ 50 ครั้งต่อวินาที และวัดได้ 50.1 ครั้งต่อวินาทีเมื่อเปิดเฉพาะฐานหุ่นยนต์กับ SLAM)
และ \texttt{diff\_cont} เผยแพร่ \texttt{odom}$\rightarrow$\texttt{base\_link} ที่ 30 ครั้งต่อวินาที
ส่วนเส้นเชื่อมที่แสดงอัตรา 10,000 คือการแปลงพิกัดคงที่ (static) จาก URDF

\begin{figure}[htbp]
\centering
\includegraphics[width=\textwidth]{images/tf_tree.png}
\caption[ต้นไม้ของระบบพิกัดขณะระบบทำงานครบทุกโหนด (บทที่ 4)]{ต้นไม้ของระบบพิกัดขณะระบบทำงานครบทุกโหนด จาก \texttt{tf2\_tools view\_frames} (24 ก.ย.) แสดงผู้เผยแพร่และอัตราการเผยแพร่ของแต่ละเส้นเชื่อม}
\label{fig:tftree}
\end{figure}

โดยการแปลงพิกัดจาก \texttt{map} ไปยัง \texttt{odom}
เป็นค่าที่ slam\_toolbox คำนวณและแก้ไขให้ตลอดเวลา
ซึ่งก็คือค่าชดเชยความคลาดเคลื่อนสะสมของโอดอมิเตอร์นั่นเอง

ตำแหน่งติดตั้งจริงของไลดาร์และกล้องที่ใช้กำหนดค่าคงที่ในไฟล์ URDF แสดงดังรูปที่ \ref{fig:lidar-mount}

\begin{figure}[htbp]
\centering
\includegraphics[width=0.85\textwidth]{images/lidar_mounting.jpg}
\caption[ตำแหน่งติดตั้งไลดาร์และกล้องบนหุ่นยนต์ต้นแบบ (บทที่ 4)]{ตำแหน่งติดตั้งไลดาร์และกล้องบนหุ่นยนต์ต้นแบบ}
\label{fig:lidar-mount}
\end{figure}

\subsection{คุณลักษณะของไลดาร์และผลการวัดจังหวะเวลาการสแกน}

ไลดาร์ที่ใช้ในโครงงานคือ \textbf{YDLIDAR X3 Pro}
ซึ่งเป็นไลดาร์สองมิติแบบช่องสัญญาณเดียว เชื่อมต่อผ่านพอร์ตอนุกรมยูเอสบีที่อัตราบอด 115200
ค่าที่ตั้งไว้ในไฟล์กำหนดค่าและค่าที่วัดได้จากระบบจริงแสดงดังตารางที่ \ref{tab:lidar}

\begin{table}[htbp]
\centering
\caption[คุณลักษณะของไลดาร์และผลการวัดบนระบบจริง (บทที่ 4)]{คุณลักษณะของไลดาร์และผลการวัดบนระบบจริง}
\label{tab:lidar}
\small
\begin{tabularx}{\textwidth}{|X|c|>{\raggedright\arraybackslash}p{3.6cm}|}
\hline
\textbf{รายการ} & \textbf{ค่า} & \textbf{ที่มา} \\
\hline
ระยะวัดต่ำสุด & 0.12 เมตร & ค่าตามสเปก \\
\hline
ระยะวัดสูงสุด & 8.0 เมตร & ค่าตามสเปก \\
\hline
จำนวนลำแสงต่อหนึ่งรอบการสแกน & 350 ลำแสง & วัดจากระบบจริง \\
\hline
ความละเอียดเชิงมุมโดยประมาณ & 1.03$^\circ$ & คำนวณจากจำนวนลำแสง \\
\hline
คาบการสแกนหนึ่งรอบ & 86.2 มิลลิวินาที & วัดจากระบบจริง \\
\hline
อัตราการสแกน & 11.61 เฮิรตซ์ & วัดจากระบบจริง (ตั้งค่าไว้ 10.0) \\
\hline
สัดส่วนลำแสงที่ไม่ได้ค่ากลับมา & ร้อยละ 19--28 & วัดจากระบบจริง (ขึ้นกับตำแหน่ง) \\
\hline
\end{tabularx}
\end{table}

ประเด็นที่มีผลต่อการออกแบบส่วนอื่นของระบบมีสองข้อ

ข้อแรกคือ \textbf{ลำแสงที่ไม่ได้ค่ากลับมามีสัดส่วนสูง}
ซึ่งเป็นลักษณะปกติของไลดาร์ระดับนี้เมื่อใช้ในอาคาร
การวัดด้วยการนับลำแสงจาก 100 รอบการสแกน (35{,}000 ลำแสง) ขณะหุ่นยนต์หยุดนิ่ง
ให้ค่าร้อยละ 27.9 ร้อยละ 25.7 และร้อยละ 19.0 ในสามตำแหน่งที่ต่างกันภายในห้องเดียวกัน
\textbf{ค่านี้จึงเป็นคุณสมบัติของห้องมากกว่าของตัวเซนเซอร์}
เพราะทิศทางที่ไม่มีวัตถุอยู่ภายในระยะวัดสูงสุดจะคืนค่าเป็นศูนย์เหมือนกับลำแสงที่ไม่ได้ค่าทุกประการ
และไม่สามารถแยกออกจากกันได้จากข้อความการสแกน
การออกแบบชั้นผสานข้อมูลจึงใช้ค่าที่แย่ที่สุดคือประมาณร้อยละ 28 เป็นเกณฑ์
และเป็นเหตุผลโดยตรงที่ชั้นผสานข้อมูลกำหนดจำนวนลำแสงขั้นต่ำไว้ที่ 3 ลำแสง
เพราะวัตถุกว้าง 0.3 เมตรที่ระยะ 3 เมตรครอบคลุมลำแสงประมาณ 6 ลำแสง
และจะเหลือลำแสงที่วัดค่าได้จริงราว 4 ลำแสงเท่านั้น
อนึ่ง สัดส่วนนี้ไม่ได้กระจายเท่ากันทุกทิศทาง เมื่อแบ่งเป็นสิบสองส่วนตามมุม
ส่วนที่แย่ที่สุดในการวัดครั้งหนึ่งสูงถึงร้อยละ 43.8 ในทิศทางตรงหน้าหุ่นยนต์
ซึ่งที่ระยะ 3 เมตรจะเหลือลำแสงที่ใช้ได้เพียงประมาณ 3.4 ลำแสง
นับเป็นกรณีที่คาบเส้นเกณฑ์พอดีและเป็นทิศทางที่หุ่นยนต์วิ่งเข้าหาเสมอ
นอกจากนี้ไดรเวอร์ถูกตั้งให้คืนค่าลำแสงที่ไม่ได้ค่าเป็นศูนย์แทนค่าอนันต์
ผู้ใช้ข้อมูลจึงต้องกรองด้วยค่าระยะต่ำสุดที่มากับข้อความการสแกนเอง
มิเช่นนั้นจะตีความลำแสงที่ไม่ได้ค่าเป็นสิ่งกีดขวางที่ระยะศูนย์เมตร

ข้อที่สองคือ \textbf{หัวข้อข้อมูลการสแกนใช้คุณภาพการส่งแบบไม่รับประกันการส่งถึง}
(Best Effort) ตามมาตรฐานของข้อมูลเซนเซอร์
โปรแกรมใดที่สมัครรับข้อมูลด้วยคุณภาพการส่งแบบรับประกันการส่งถึงจะไม่ได้รับข้อความเลย
โดยระบบแจ้งเตือนเพียงครั้งเดียวตอนเริ่มต้นแล้วเงียบไป
ซึ่งมีอาการเหมือนหัวข้อที่ไม่มีผู้เผยแพร่ทุกประการ


\subsection{การตั้งค่าการทำงานของ slam\_toolbox}

โครงงานนี้ใช้โหมดออนไลน์แบบอะซิงโครนัส
เพื่อให้ระบบยังคงตอบสนองได้แม้เมื่อหน่วยประมวลผลถูกใช้งานร่วมกับการอนุมานแบบจำลอง
พารามิเตอร์ที่มีผลต่อคุณภาพของแผนที่และภาระการประมวลผลมากที่สุดได้แก่
ความละเอียดของแผนที่ ระยะทางและมุมที่หุ่นยนต์ต้องเคลื่อนที่ก่อนจึงจะประมวลผลข้อมูลชุดใหม่
และช่วงเวลาที่ใช้ในการตรวจหาการวนกลับมาที่เดิม

การกำหนดระยะทางขั้นต่ำก่อนประมวลผลชุดข้อมูลใหม่มีผลโดยตรงต่อจำนวนปมในกราฟ
หากตั้งค่าน้อยเกินไปจะได้แผนที่ละเอียดแต่กราฟใหญ่และประมวลผลช้า
หากตั้งมากเกินไปจะเบาแต่แผนที่หยาบและมีโอกาสที่การจับคู่ข้อมูลสแกนจะล้มเหลว

แผนที่ที่ได้เมื่อสำรวจพื้นที่ทดสอบครบหนึ่งรอบแสดงดังรูปที่ \ref{fig:map-result}

\begin{figure}[htbp]
\centering
\includegraphics[width=\textwidth]{images/slam_map_result.png}
\caption[แผนที่ของพื้นที่ทดสอบที่ได้จากกระบวนการ SLAM (บทที่ 4)]{แผนที่ของพื้นที่ทดสอบที่ได้จากกระบวนการ SLAM}
\label{fig:map-result}
\end{figure}

พารามิเตอร์หลักที่ใช้จริงในโครงงานแสดงดังตารางที่ \ref{tab:slamparam}

\begin{table}[htbp]
\centering
\caption[พารามิเตอร์หลักของ slam\_toolbox ที่ใช้ในโครงงาน (บทที่ 4)]{พารามิเตอร์หลักของ slam\_toolbox ที่ใช้ในโครงงาน}
\label{tab:slamparam}
\small
\begin{tabularx}{\textwidth}{|>{\raggedright\arraybackslash}p{5.2cm}|c|X|}
\hline
\textbf{พารามิเตอร์} & \textbf{ค่าที่ใช้} & \textbf{ผลต่อระบบ} \\
\hline
โหมดการทำงาน & Mapping & สร้างแผนที่ใหม่ทุกครั้งที่เริ่มทำงาน \\
\hline
ตัวแก้ปัญหากราฟ & Ceres & ปรับตำแหน่งปมทั้งกราฟเมื่อมีข้อมูลใหม่ \\
\hline
เฟรมฐานที่ใช้อ้างอิง & \texttt{base\_footprint} & เงาของหุ่นยนต์บนพื้น \\
\hline
ความละเอียดของแผนที่ & 0.05 เมตรต่อช่อง & ความละเอียดเทียบกับขนาดหน่วยความจำ \\
\hline
ระยะทางขั้นต่ำก่อนประมวลผลชุดใหม่ & 0.20 เมตร & จำนวนปมในกราฟและภาระประมวลผล \\
\hline
มุมขั้นต่ำก่อนประมวลผลชุดใหม่ & 0.20 เรเดียน & เช่นเดียวกับข้างต้น \\
\hline
ระยะไลดาร์สูงสุดที่นำมาสร้างแผนที่ & 8.0 เมตร & ต้องไม่เกินระยะวัดจริงของไลดาร์ \\
\hline
คาบการปรับปรุงแผนที่ & 2.0 วินาที & ความถี่ในการเผยแพร่แผนที่ \\
\hline
คาบการกระจายการแปลงพิกัด & 0.02 วินาที & ความต่อเนื่องของเฟรม \texttt{map} \\
\hline
ระยะค้นหาการวนกลับที่เดิม & 3.0 เมตร & ขอบเขตการตรวจหา Loop Closure \\
\hline
จำนวนสแกนขั้นต่ำของโซ่ที่ยอมรับ & 15 & ลดการจับคู่ผิดในทางเดินที่มีลักษณะซ้ำกัน \\
\hline
ขอบเขตการค้นหาเชิงมุมแบบหยาบ & 0.175 เรเดียน & ค้นหา $\pm$10 องศารอบค่าคาดการณ์จากโอดอมิเตอร์ \\
\hline
\end{tabularx}
\end{table}


\subsection{การจัดการอุปกรณ์ต่อพ่วง}

เนื่องจากทั้งไลดาร์และไมโครคอนโทรลเลอร์เชื่อมต่อผ่านพอร์ตอนุกรมยูเอสบี
ชื่ออุปกรณ์ที่ระบบปฏิบัติการกำหนดให้อาจสลับกันได้ในแต่ละครั้งที่เปิดเครื่อง
ระบบจึงกำหนดกฎ udev เพื่อสร้างชื่ออุปกรณ์ถาวรคือ \texttt{/dev/ydlidar}
และ \texttt{/dev/esp32}

ประเด็นสำคัญคือ
\textbf{ไม่สามารถอ้างอิงจากหมายเลขผู้ผลิตและหมายเลขรุ่นของชิปแปลงสัญญาณได้}
เพราะอุปกรณ์ทั้งสองใช้ชิปแปลงสัญญาณตระกูลเดียวกัน (CP210x)
จึงมีหมายเลขผู้ผลิตและหมายเลขรุ่นตรงกันทุกประการ
และทั้งคู่ไม่มีหมายเลขประจำเครื่อง (Serial Number) ที่ใช้แยกความแตกต่างได้
กฎ udev ของระบบจึงอ้างอิงจาก
\textbf{ตำแหน่งพอร์ตยูเอสบีที่อุปกรณ์เสียบอยู่} (\texttt{KERNELS}) แทน
ซึ่งใช้ได้เพราะอุปกรณ์ถูกเสียบไว้ที่พอร์ตเดิมเสมอ
ข้อแลกเปลี่ยนคือหากย้ายพอร์ตหรือเปลี่ยนบอร์ดหลัก จะต้องสร้างกฎใหม่

\section{ระบบนำทางอัตโนมัติ}

ระบบนำทางใช้ Nav2 ทำงานร่วมกับแผนที่ที่ slam\_toolbox สร้างขึ้นในรอบเดียวกัน
โดยไม่ใช้ระบบระบุตำแหน่งบนแผนที่ที่บันทึกไว้ล่วงหน้า
เนื่องจากการสาธิตของโครงงานสร้างแผนที่และนำทางภายในการทำงานรอบเดียว
การออกแบบจึงลดจำนวนระบบย่อยลงหนึ่งระบบ
แผนที่ต้นทุนเฉพาะที่ (Local Costmap) ที่ระบบนำทางใช้หลบสิ่งกีดขวางรอบตัวหุ่นยนต์
แสดงดังรูปที่ \ref{fig:nav2-costmap}

\begin{figure}[htbp]
\centering
\includegraphics[width=0.85\textwidth]{images/nav2_costmap_path.png}
\caption[แผนที่ต้นทุนเฉพาะที่ของระบบนำทาง (บทที่ 4)]{แผนที่ต้นทุนเฉพาะที่ของระบบนำทางใน RViz ขณะเปิดทั้งระบบ
ครอบคลุมพื้นที่ 3$\times$3 เมตรรอบหุ่นยนต์ (สี่เหลี่ยมสีส้มคือรูปร่างของหุ่นยนต์) ที่ความละเอียด 0.05 เมตรต่อช่อง
สีม่วงเข้มคือช่องที่ไลดาร์พบสิ่งกีดขวาง สีฟ้าคือช่องที่หุ่นยนต์จะชนแน่นอนหากศูนย์กลางเข้าไป
และสีแดงถึงม่วงคือบริเวณขยายขอบที่ต้นทุนลดลงตามระยะห่าง จุดสีแดงคือข้อมูลจากไลดาร์}
\label{fig:nav2-costmap}
\end{figure}

ประเด็นการตั้งค่าที่มีผลมากที่สุดสองข้อมีดังนี้

\begin{itemize}[leftmargin=1.25cm, itemsep=6pt, topsep=6pt]
    \item \textbf{การระบุรูปร่างของหุ่นยนต์เป็นรูปหลายเหลี่ยมแทนรัศมีวงกลม}
    หุ่นยนต์ต้นแบบมีรูปทรงสี่เหลี่ยมที่ไม่สมมาตรรอบจุดกึ่งกลางการหมุน
    ระบบจึงระบุขอบเขตเป็นรูปสี่เหลี่ยมที่กินพื้นที่ด้านหน้า 0.09 เมตร
    ด้านหลัง 0.265 เมตร และด้านข้างข้างละ 0.147 เมตร
    หากระบุเป็นรัศมีวงกลมที่ครอบรูปทรงทั้งหมด หุ่นยนต์จะไม่สามารถผ่านช่องทางแคบที่ผ่านได้จริง
    ข้อควรระวังคือหากกำหนดทั้งสองค่าพร้อมกัน ระบบจะใช้ค่ารูปหลายเหลี่ยมและละเลยค่ารัศมีโดยไม่แจ้งเตือน

    \item \textbf{การจำกัดความเร็วให้ต่ำกว่าที่คอนโทรลเลอร์ยอมรับมาก}
    ระบบนำทางถูกจำกัดความเร็วเชิงเส้นไว้ที่ 0.10 เมตรต่อวินาที
    และความเร็วเชิงมุมที่ 0.25 เรเดียนต่อวินาที
    ซึ่งต่ำกว่าขีดจำกัดของคอนโทรลเลอร์ที่ 0.30 เมตรต่อวินาทีมาก
    เหตุผลอยู่ที่ผลการวัดจังหวะเวลาการสแกนในหัวข้อก่อนหน้า
    ที่ความเร็ว 0.10 เมตรต่อวินาที ข้อมูลการสแกนหนึ่งรอบเหลื่อมกันเพียงประมาณ 0.9 เซนติเมตร
    และอธิบายไว้ในหัวข้อปัญหาและอุปสรรคของบทที่ 5
    ทั้งนี้ค่าดังกล่าวถูกปรับขึ้นจาก 0.055 เมตรต่อวินาทีและ 0.125 เรเดียนต่อวินาที
    ในช่วงท้ายของโครงงาน ผลการทดสอบการนำทางและการหมุนตัวเพื่อกู้คืนสถานการณ์ด้านล่าง
    \textbf{จึงเป็นผลที่วัดได้ที่ค่าเดิม}
\end{itemize}

ผลการทดสอบระบบนำทางบนพื้นที่จริงแสดงดังตารางที่ \ref{tab:navresult}
โดยกำหนดเป้าหมายจากเครื่องมือแสดงผลของ ROS 2 บนแผนที่ที่ slam\_toolbox
กำลังสร้างอยู่ในรอบเดียวกัน

\begin{table}[htbp]
\centering
\caption[ผลการทดสอบการนำทางไปยังเป้าหมายบนแผนที่ (บทที่ 4)]{ผลการทดสอบการนำทางไปยังเป้าหมายบนแผนที่}
\label{tab:navresult}
\small
\begin{tabularx}{\textwidth}{|c|X|c|c|c|}
\hline
\textbf{เป้าหมาย} & \textbf{จากตำแหน่ง $\rightarrow$ ถึงตำแหน่ง (เมตร)}
 & \textbf{ระยะตรง} & \textbf{เวลา} & \textbf{ผล} \\
\hline
A & ($-$0.83, $-$2.04) $\rightarrow$ (3.43, 0.52) & 4.97 ม. & 115.5 วินาที & สำเร็จ \\
\hline
B & (3.35, 0.41) $\rightarrow$ (2.90, 0.24) & 0.48 ม. & 31.8 วินาที & สำเร็จ \\
\hline
C & (3.01, 0.30) $\rightarrow$ (3.08, 0.37) & 0.10 ม. & 6.9 วินาที & สำเร็จ \\
\hline
\end{tabularx}
\end{table}

ความเร็วเฉลี่ยของเป้าหมาย A คือ 0.043 เมตรต่อวินาที
ซึ่งต่ำกว่าขีดจำกัด 0.055 เมตรต่อวินาทีที่ตั้งไว้ขณะทดสอบเล็กน้อยตามที่คาดไว้สำหรับเส้นทางที่มีการเลี้ยว
และยืนยันว่าขีดจำกัดความเร็วที่ตั้งไว้เป็นตัวกำหนดความเร็วจริงของระบบ
\textbf{ข้อสังเกตที่ต้องระบุไว้คือเป้าหมาย C ไม่ใช่การทดสอบที่สมบูรณ์}
เพราะระยะ 0.10 เมตรอยู่ภายในเกณฑ์ความคลาดเคลื่อนที่ยอมรับได้ของเป้าหมาย (0.15 เมตร)
หุ่นยนต์จึงอยู่ในเกณฑ์ตั้งแต่ต้นและเพียงปรับทิศทางเท่านั้น
การทดสอบที่ใช้วางแผนเส้นทางจริงจึงมีสองครั้งจากสามครั้ง

นอกจากนี้ได้ทดสอบพฤติกรรมกู้คืนสถานการณ์ด้วยการสั่งให้หุ่นยนต์หมุนตัว 1.57 เรเดียน
ผลคือหมุนได้ 1.5762 เรเดียน (90.31 องศา) ในเวลา 16.40 วินาที
ตรงกับเวลาที่คำนวณไว้ล่วงหน้าจากขีดจำกัดความเร็วเชิงมุมของพฤติกรรมกู้คืนที่
0.1 เรเดียนต่อวินาที ซึ่งอยู่ในช่วง 15.7--16.5 วินาที
ค่าขีดจำกัดนี้ถูกปรับเป็น 0.2 เรเดียนต่อวินาทีในภายหลัง
เวลาการหมุนจึงคาดว่าจะลดลงเหลือประมาณครึ่งหนึ่ง แต่ยังไม่ได้วัดซ้ำที่ค่าใหม่

\subsection{การสำรวจพื้นที่อัตโนมัติ}
\label{sec:explore}

นอกจากการนำทางไปยังเป้าหมายที่ผู้ใช้กำหนดแล้ว ระบบยังสำรวจพื้นที่ได้เองโดยไม่ต้องมีผู้ควบคุม
ด้วยวิธีการสำรวจตามขอบเขตของพื้นที่ที่ยังไม่ทราบข้อมูล (Frontier-based Exploration)~\cite{yamauchi1997}
ขอบเขตในที่นี้คือแนวรอยต่อระหว่างช่องที่ทราบแล้วว่าว่างกับช่องที่ยังไม่ทราบข้อมูลบนแผนที่
การขับหุ่นยนต์ไปยังขอบเขตทำให้ไลดาร์มองเห็นพื้นที่ใหม่และแผนที่ขยายออกไปเรื่อย ๆ
จนไม่เหลือขอบเขตที่เข้าถึงได้

ระบบใช้แพ็กเกจ explore\_lite จากโครงการ m-explore-ros2~\cite{mexplore}
ซึ่งเป็นชั้นบนสุดของระบบนำทางและไม่ได้สั่งมอเตอร์โดยตรง
โหนดนี้อ่านแผนที่แบบตารางความน่าจะเป็นจาก slam\_toolbox
ค้นหาขอบเขตทั้งหมด เลือกขอบเขตที่ดีที่สุด แล้วส่งเป็นเป้าหมายปกติให้ Nav2
ลำดับการทำงานจึงเป็น slam\_toolbox $\rightarrow$ explore\_lite $\rightarrow$ Nav2 $\rightarrow$ twist\_mux $\rightarrow$ ล้อ
ผลที่ตามมาคือหากระบบนำทางไปยังเป้าหมายที่กำหนดด้วยมือไม่ได้ การสำรวจอัตโนมัติก็จะทำงานไม่ได้เช่นกัน
ผู้ใช้เปิดการสำรวจด้วยคำสั่ง \texttt{make explore} หลังจากเปิดฐานหุ่นยนต์ SLAM และ Nav2 แล้ว

ขอบเขตแต่ละแห่งถูกให้คะแนนจากระยะทางและขนาด
โดยค่าต้นทุนเท่ากับ $3.0 \times$ ระยะทาง $-$ $1.0 \times$ ขนาดของขอบเขต
ระบบจึงเลือกขอบเขตที่อยู่ใกล้ก่อน ซึ่งเหมาะกับหุ่นยนต์ที่เคลื่อนที่ช้า
เพราะไม่ต้องข้ามพื้นที่ไปมาทุกครั้งที่เลือกเป้าหมายใหม่
การตั้งค่าที่ปรับให้เข้ากับหุ่นยนต์ต้นแบบมีดังนี้

\begin{itemize}[leftmargin=1.25cm, itemsep=6pt, topsep=6pt]
    \item \textbf{ใช้แผนที่จาก SLAM โดยตรง ไม่ใช่แผนที่ต้นทุนของ Nav2}
    เพราะแผนที่ต้นทุนผ่านการขยายขอบสิ่งกีดขวางแล้ว และอาจระบุช่องที่ยังไม่ทราบข้อมูลเป็นสิ่งกีดขวาง
    ซึ่งจะทำให้มองไม่เห็นขอบเขตใดเลย

    \item \textbf{ลดขนาดขอบเขตขั้นต่ำจาก 0.75 เมตรเป็น 0.3 เมตร}
    ค่าเริ่มต้น 0.75 เมตรกว้างกว่าประตูภายในอาคาร
    ช่องทางที่หุ่นยนต์ (กว้าง 0.29 เมตร) ผ่านได้จริงจึงถูกมองว่าเล็กเกินกว่าจะสำรวจ
    และทำให้บางห้องไม่ถูกสำรวจเลย ค่า 0.3 เมตรยังกรองสัญญาณรบกวนระดับช่องเดียวที่ขอบแผนที่ออกได้

    \item \textbf{กำหนดเวลาที่ไม่มีความคืบหน้าตามระยะทาง ไม่ใช่ตามเวลา}
    หากหุ่นยนต์ไม่เข้าใกล้ขอบเขตที่เลือกภายในเวลาที่กำหนด ขอบเขตนั้นจะถูกตัดทิ้งและเลือกขอบเขตใหม่
    ค่าเริ่มต้น 30 วินาทีที่ความเร็ว 0.22 เมตรต่อวินาทีเทียบเท่าระยะประมาณ 6.6 เมตร
    ระบบจึงคงระยะนี้ไว้และตั้งเวลาเป็น 66 วินาทีที่ความเร็ว 0.10 เมตรต่อวินาที
    หากใช้ 30 วินาทีเท่าเดิม หุ่นยนต์ที่กำลังอ้อมสิ่งกีดขวางจะถูกตัดขอบเขตที่ดีทิ้งไปทีละแห่ง
    และหยุดสำรวจทั้งที่แผนที่ยังไม่ครบ

    \item \textbf{กลับไปยังจุดเริ่มต้นเมื่อไม่เหลือขอบเขต}
    ทำให้เห็นชัดว่าการสำรวจเสร็จสิ้นแล้ว เพราะหุ่นยนต์ที่เคลื่อนกลับมาที่จุดเดิมคือสัญญาณว่าจบรอบ
\end{itemize}

ระหว่างการสำรวจ โหนดจะเผยแพร่ตำแหน่งของขอบเขตที่พบเป็นเครื่องหมายในเครื่องมือแสดงผล
ผู้ใช้จึงดูได้ว่าหุ่นยนต์กำลังมุ่งไปยังขอบเขตใดและเพราะเหตุใด
ทั้งนี้การสำรวจอัตโนมัติได้ทดลองใช้งานในพื้นที่ทดสอบแล้ว
แต่ยังไม่ได้วัดผลเชิงตัวเลข เช่น เวลาที่ใช้ หรือสัดส่วนของพื้นที่ที่สำรวจได้เทียบกับพื้นที่จริง


\section{ระบบตรวจจับวัตถุบนอุปกรณ์เอดจ์เอไอ}

\subsection{การเตรียมแบบจำลอง: ONNX และ TensorRT}

ในรุ่นแรกของระบบ แบบจำลอง YOLO26s ที่ผ่านการฝึกบนชุดข้อมูล COCO ถูกส่งออกเป็นไฟล์รูปแบบ ONNX
โดยกำหนดความละเอียดของการคำนวณเป็นทศนิยม 16 บิต (FP16)
และกำหนดขนาดภาพเข้าคงที่ที่ 640$\times$640 พิกเซล
การส่งออกทำบนอุปกรณ์เป้าหมายด้วยสคริปต์ที่เขียนขึ้นเอง
และออกแบบให้ข้ามขั้นตอนนี้โดยอัตโนมัติหากไฟล์ที่มีอยู่ตรงกับเงื่อนไขที่ร้องขอแล้ว
เพื่อไม่ให้การเปิดระบบแต่ละครั้งต้องเสียเวลาส่งออกซ้ำ

ข้อควรระวังคือขนาดภาพเข้าถูกผนึกไว้ในโครงสร้างของไฟล์ ONNX ตั้งแต่ตอนส่งออก
การเปลี่ยนขนาดภาพที่ใช้อนุมานจึงต้องส่งออกไฟล์ใหม่
สคริปต์ส่งออกจึงอ่านข้อมูลกำกับ (Metadata) ที่ฝังอยู่ในไฟล์เดิมมาตรวจสอบก่อนเสมอ

ในรุ่นที่ใช้งานจริง แบบจำลอง YOLO26l ถูกส่งออกเป็นเครื่องยนต์ TensorRT แบบ FP16
โดยกำหนดขนาดภาพเข้าเป็น 480$\times$640 พิกเซลให้ตรงกับภาพของกล้อง
เนื่องจากการใช้ขนาดสี่เหลี่ยมจัตุรัส 640$\times$640 ทำให้ภาพถูกเติมขอบว่าง
และค่า F1 ลดลงประมาณ 2 จุด (หัวข้อ \ref{sec:eval-det})
เครื่องยนต์อนุมานถูกสร้างจากไฟล์น้ำหนักต้นฉบับ \texttt{yolo26l.pt} บนอุปกรณ์เป้าหมายโดยตรง
แบบขนาดภาพเข้าคงที่และขนาดชุดข้อมูลเท่ากับหนึ่งภาพ
และโหนดตรวจจับอ่านขนาดภาพเข้าจากข้อมูลกำกับของแบบจำลองเองโดยไม่ต้องกำหนดซ้ำ

การสร้างเครื่องยนต์เป็นขั้นตอนที่มีต้นทุนสูงที่สุดของระบบนี้
คือใช้เวลาประมาณ 13--14 นาที และใช้หน่วยความจำประมาณ 3.1 กิกะไบต์
จากหน่วยความจำ 7.4 กิกะไบต์ที่หน่วยประมวลผลกลางและหน่วยประมวลผลกราฟิกใช้ร่วมกัน
โดยอุปกรณ์ไม่มีพื้นที่แลกเปลี่ยนบนดิสก์ สคริปต์สร้างเครื่องยนต์จึงออกแบบไว้สามข้อ
ข้อแรก ข้ามการสร้างหากมีไฟล์เครื่องยนต์อยู่แล้ว คำสั่ง \texttt{make yolo} จึงเสียเวลาสร้างเพียงครั้งแรก
ข้อสอง ปฏิเสธการสร้างหากหน่วยความจำว่างน้อยกว่า 4.5 กิกะไบต์
ซึ่งหมายความว่าต้องปิดส่วนอื่นของระบบก่อน เพื่อไม่ให้ระบบปฏิบัติการหยุดโปรแกรมที่กำลังทำงานเพราะหน่วยความจำไม่พอ
ข้อสาม สร้างไฟล์ในโฟลเดอร์ชั่วคราวแล้วจึงย้ายไปยังตำแหน่งจริง
เพื่อไม่ให้มีไฟล์ระหว่างทางค้างอยู่ข้างไฟล์ต้นฉบับ
เนื่องจากเครื่องยนต์ใช้ได้เฉพาะกับรุ่นของ TensorRT ที่สร้างมัน
ไฟล์ \texttt{.pt} จึงเป็นแหล่งที่มาหลัก ส่วนเครื่องยนต์เป็นเพียงไฟล์ที่สร้างใหม่ได้เสมอ
ส่วนเส้นทาง ONNX ของรุ่นแรกยังคงอยู่เป็นทางเลือกสำรอง ซึ่งเปลี่ยนกลับได้ด้วยการระบุไฟล์แบบจำลองและค่าความเชื่อมั่นในคำสั่ง

สภาพแวดล้อมการทำงานของไลบรารีเรียนรู้เชิงลึกถูกจัดเก็บเป็นสภาพแวดล้อมเสมือนของภาษาไพทอน
(Virtual Environment) แทนการใช้คอนเทนเนอร์
ประเด็นที่ต้องระวังที่สุดคือไลบรารีที่ติดตั้งจากแหล่งมาตรฐานมักเป็นรุ่นที่รองรับเฉพาะ
หน่วยประมวลผลกลาง
เมื่อนำมาใช้บนอุปกรณ์เป้าหมายแบบจำลองจะไม่ได้ใช้หน่วยประมวลผลกราฟิกเลย
และจะทำงานช้าลงหลายเท่าโดยไม่มีข้อความแจ้งเตือนใด ๆ
จึงต้องติดตั้งไลบรารีรุ่นที่จัดทำมาสำหรับอุปกรณ์ตระกูลนี้โดยเฉพาะ
และตรวจสอบทุกครั้งว่าระบบเลือกใช้หน่วยประมวลผลกราฟิกจริง

รุ่นของซอฟต์แวร์และขนาดของไฟล์แบบจำลองที่ใช้จริงแสดงดังตารางที่ \ref{tab:aienv}

\begin{table}[htbp]
\centering
\caption[สภาพแวดล้อมซอฟต์แวร์ของระบบตรวจจับวัตถุ (บทที่ 4)]{สภาพแวดล้อมซอฟต์แวร์ของระบบตรวจจับวัตถุ}
\label{tab:aienv}
\small
\begin{tabularx}{\textwidth}{|X|c|}
\hline
\textbf{รายการ} & \textbf{รุ่น / ค่าที่ใช้} \\
\hline
ชุดซอฟต์แวร์ของอุปกรณ์ (JetPack) & 6.1 (L4T 36.4.0) \\
\hline
CUDA / cuDNN / TensorRT & 12.6.68 / 9.3.0 / 10.3.0 \\
\hline
PyTorch / TorchVision (รุ่นสำหรับสถาปัตยกรรม aarch64) & 2.11.0 / 0.26.0 \\
\hline
ONNX Runtime (GPU) & 1.24.0 \\
\hline
ไลบรารีแบบจำลองตรวจจับวัตถุ & ultralytics 8.4.144 \\
\hline
ตัวติดตามวัตถุข้ามเฟรม & ByteTrack \\
\hline
\textbf{แบบจำลองที่ใช้งานจริง} & \textbf{\texttt{yolo26l\_480x640.engine}} \\
& \textbf{(TensorRT FP16, 480$\times$640)} \\
\hline
เวลาและหน่วยความจำที่ใช้สร้างเครื่องยนต์หนึ่งครั้ง & 13--14 นาที, ประมาณ 3.1 กิกะไบต์ \\
\hline
ค่าความเชื่อมั่นขั้นต่ำของการตรวจจับ & 0.4 \\
\hline
ขนาดไฟล์น้ำหนักต้นทาง \texttt{yolo26s.pt} (รุ่นแรก) & 20.4 เมกะไบต์ \\
\hline
ขนาดไฟล์ \texttt{yolo26s.onnx} แบบ 16 บิต (เทียบกับ 32 บิต) (รุ่นแรก) & 19.2 (38.3) เมกะไบต์ \\
\hline
เวลาที่ใช้ในการส่งออกไฟล์ ONNX หนึ่งครั้ง (รุ่นแรก) & 5.5--6.0 วินาที \\
\hline
\end{tabularx}
\end{table}

ขั้นตอนการเตรียมแบบจำลองเมื่อสั่ง \texttt{make yolo} แสดงดังรูปที่ \ref{fig:onnx-export}

\begin{figure}[htbp]
\centering
\includegraphics[width=\textwidth]{images/model_prep_flow.png}
\caption[ขั้นตอนการเตรียมแบบจำลองก่อนเปิดโหนดตรวจจับ (บทที่ 4)]{ขั้นตอนการเตรียมแบบจำลองก่อนเปิดโหนดตรวจจับ
เส้นทางหลักคือเครื่องยนต์ TensorRT ของ YOLO26l ซึ่งสร้างเพียงครั้งแรก
ส่วนเส้นทาง ONNX ของ YOLO26s เป็นทางเลือกสำรองของรุ่นแรก}
\label{fig:onnx-export}
\end{figure}

\subsection{กระบวนการอนุมานและผลลัพธ์}

โหนดตรวจจับวัตถุรับภาพจากกล้องยูเอสบี ปรับขนาดและทำให้เป็นมาตรฐานตามที่แบบจำลองต้องการ
ส่งเข้าแบบจำลองแล้วเผยแพร่ผลลัพธ์เป็นรายการของกรอบสี่เหลี่ยม
โดยแต่ละรายการประกอบด้วยพิกัดมุมของกรอบในหน่วยพิกเซล ชื่อคลาส และค่าความเชื่อมั่น
ผลการตรวจจับที่มีค่าความเชื่อมั่นต่ำกว่าเกณฑ์ที่กำหนดจะถูกคัดออกก่อนส่งต่อ

นอกจากนี้โหนดยังติดหมายเลขติดตามวัตถุ (Track ID) ให้กับกรอบแต่ละกรอบ
โดยใช้ตัวติดตาม ByteTrack~\cite{bytetrack} ที่มากับไลบรารีของแบบจำลอง
หมายเลขนี้เป็นสิ่งที่ทำให้ชั้นผสานข้อมูลสามารถแยกได้ว่า
``เก้าอี้ตัวเดิมที่เห็นอีกครั้ง'' ต่างจาก ``เก้าอี้ตัวที่สอง''
และเป็นกุญแจหลักในการป้องกันการบันทึกวัตถุซ้ำ

อีกประเด็นหนึ่งที่มีผลต่อความถูกต้องคือ ประทับเวลาของข้อความผลการตรวจจับ
ถูกกำหนดให้เท่ากับ \textbf{เวลาที่ถ่ายภาพ} ไม่ใช่เวลาที่เผยแพร่ผลลัพธ์
เพราะชั้นผสานข้อมูลจะใช้ประทับเวลานี้ไปสอบถามการแปลงพิกัดย้อนหลัง

อัตราเฟรมที่วัดได้ขณะระบบทำงานจริงแสดงดังรูปที่ \ref{fig:fps-gpu}

\begin{figure}[htbp]
\centering
\includegraphics[width=\textwidth]{images/inference_fps_gpu.png}
\caption[อัตราเฟรมและการใช้งานหน่วยประมวลผลกราฟิกขณะอนุมาน (บทที่ 4)]{อัตราเฟรมและการใช้งานหน่วยประมวลผลกราฟิกขณะอนุมาน
(ก) อัตราเฟรมของการตรวจจับแต่ละแบบจำลองและแต่ละรอบการวัด เทียบกับเกณฑ์ 5 เฟรมต่อวินาที
สีน้ำเงินคือ YOLO26l บน TensorRT ที่ใช้งานจริง สีส้มคือ YOLO26s แบบ ONNX ของรุ่นแรก
(ข) อัตราข้อความของหัวข้อหลักที่ทำงานพร้อมกันในรอบการสาธิตวันที่ 22 ก.ย. ซึ่งยืนยันว่า SLAM ทำงานอยู่ระหว่างการตรวจจับ}
\label{fig:fps-gpu}
\end{figure}

\subsection{ผลการวัดสมรรถนะบนอุปกรณ์จริง}

ผลการวัดเวลาประมวลผลของแต่ละรูปแบบแบบจำลองบนอุปกรณ์เป้าหมาย
โดยใช้ภาพทดสอบขนาด 640$\times$480 พิกเซล เฉลี่ยจาก 100 เฟรมหลังการอุ่นเครื่อง 20 เฟรม
แยกกระบวนการทดสอบของแต่ละรูปแบบ และล็อกสัญญาณนาฬิกาของหน่วยประมวลผลกราฟิกไว้ที่ 624.75 MHz
แสดงดังตารางที่ \ref{tab:modelbench}
เมื่อใช้การปรับสัญญาณนาฬิกาตามค่าเริ่มต้นแบบที่หุ่นยนต์ใช้งานจริง
ผลของรูปแบบ ONNX ต่างจากตารางไม่เกินประมาณ 2 มิลลิวินาที
ส่วนรูปแบบ PyTorch ได้เวลาเท่าเดิม

\begin{table}[htbp]
\centering
\caption[เวลาประมวลผลต่อเฟรมของรูปแบบแบบจำลองที่พิจารณา (บทที่ 4)]{เวลาประมวลผลต่อเฟรมของรูปแบบแบบจำลองที่พิจารณา}
\label{tab:modelbench}
\small
\begin{tabularx}{\textwidth}{|X|c|c|c|}
\hline
\textbf{รูปแบบแบบจำลอง} & \textbf{ต่อเฟรม (ms)} & \textbf{รวมตัวติดตาม (ms)} & \textbf{GPU} \\
\hline
YOLO26n รูปแบบ PyTorch & 36.3 & 59.4 & ร้อยละ 45 \\
\hline
YOLO26s รูปแบบ PyTorch & 36.7 & 60.3 & ร้อยละ 85 \\
\hline
YOLO26s รูปแบบ ONNX ความละเอียด 32 บิต & 43.1 & 66.4 & ร้อยละ 94 \\
\hline
\textbf{YOLO26s รูปแบบ ONNX ความละเอียด 16 บิต} & \textbf{33.2} & \textbf{56.5} & \textbf{ร้อยละ 93} \\
\hline
YOLO26n รูปแบบ ONNX ความละเอียด 16 บิต & 21.9 & 46.8 & ร้อยละ 89 \\
\hline
\end{tabularx}
\end{table}

หมายเหตุ: ``ต่อเฟรม'' คือเวลาเตรียมภาพ อนุมาน และประมวลผลหลังการอนุมานรวมกัน
ส่วน ``รวมตัวติดตาม'' คือเวลาเมื่อรวมการติดตามวัตถุด้วย ByteTrack ซึ่งเป็นวิธีที่ระบบใช้งานจริง

ตารางนี้ยืนยันข้อสรุปสามข้อ
ข้อแรกคือ \textbf{การเรียกใช้ผ่าน PyTorch ถูกจำกัดด้วยการสั่งงานทีละคำสั่งจากภาษาไพทอน}
ไม่ใช่ด้วยการคำนวณ รุ่น Nano และรุ่น Small ใช้เวลาเกือบเท่ากัน (36.3 และ 36.7 มิลลิวินาที)
ทั้งที่รุ่น Small มีปริมาณการคำนวณมากกว่าประมาณ 3.5 เท่า
และหน่วยประมวลผลกราฟิกยังทำงานไม่เต็มที่
เมื่อรันผ่าน ONNX Runtime ซึ่งสั่งงานทั้งกราฟโดยไม่ผ่านภาษาไพทอน
ต้นทุนจริงของแต่ละรุ่นจึงปรากฏ คือ 21.9 มิลลิวินาทีสำหรับรุ่น Nano และ 33.2 มิลลิวินาทีสำหรับรุ่น Small
ข้อที่สองคือ \textbf{การส่งออกเป็นความละเอียด 16 บิตเป็นเงื่อนไขที่จำเป็น}
เพราะรูปแบบ ONNX ความละเอียด 32 บิตช้ากว่าการเรียกใช้ผ่าน PyTorch โดยตรง
ขณะที่ความละเอียด 16 บิตเร็วที่สุดในบรรดารูปแบบของรุ่น Small
ข้อที่สามคือทุกรูปแบบแม้รวมตัวติดตามแล้วยังใช้เวลาน้อยกว่า 66.7 มิลลิวินาทีต่อเฟรม
ซึ่งเป็นคาบเวลาของกล้องที่ 15 เฟรมต่อวินาที
รุ่น Nano แบบ ONNX 16 บิตแม้จะเร็วที่สุดจึงไม่ได้ถูกเลือก
เพราะความเร็วไม่ใช่ข้อจำกัดของระบบ และรุ่น Small ให้ความแม่นยำที่สูงกว่า
ด้วยเหตุผลเดียวกันนี้ ระบบรุ่นสุดท้ายจึงเปลี่ยนเป็นรุ่น Large บน TensorRT (หัวข้อ \ref{sec:eval-det})
ซึ่งใช้เวลาอนุมาน 36.4 มิลลิวินาที ใกล้เคียงกับ YOLO26s แบบ ONNX ที่ใช้เดิม
เมื่อวัดในเงื่อนไขเดียวกัน (รูปที่ \ref{fig:infer-bench} และตารางที่ \ref{tab:det-cost})

อนึ่ง ผลการวัดชุดนี้ทำซ้ำเพื่อแทนผลการวัดครั้งก่อน
ซึ่งวัดเวลารวมตัวติดตามเพียงอย่างเดียวจาก 30 เฟรม
ทำให้เวลาของทุกรูปแบบดูใกล้เคียงกันจนแยกไม่ออก

ผลการวัดขณะระบบทำงานจริงของทั้งสองรุ่นแสดงดังตารางที่ \ref{tab:modellive}
คอลัมน์แรกคือ YOLO26s แบบ ONNX ของรุ่นแรก ทำงานลำพังต่อเนื่อง 301 วินาทีในสภาพแวดล้อมที่มีวัตถุหลายชนิด
คอลัมน์ที่สองคือ YOLO26l บน TensorRT ที่ใช้งานจริง ทำงานพร้อมกับ SLAM และชั้นผสานข้อมูลเป็นเวลา 120 วินาที
ขณะหุ่นยนต์จอดนิ่ง

\begin{table}[htbp]
\centering
\caption[ผลการวัดสมรรถนะของระบบตรวจจับวัตถุขณะทำงานจริง (บทที่ 4)]{ผลการวัดสมรรถนะของระบบตรวจจับวัตถุขณะทำงานจริง
(--- คือไม่ได้วัดในรอบนั้น)}
\label{tab:modellive}
\small
\begin{tabularx}{\textwidth}{|X|c|c|}
\hline
\textbf{รายการ} & \textbf{YOLO26s ONNX} & \textbf{YOLO26l TensorRT} \\
 & \textbf{(ทำงานลำพัง)} & \textbf{(พร้อม SLAM)} \\
\hline
อัตราการเผยแพร่ผลการตรวจจับ (เฟรมต่อวินาที) & 15.13 & \textbf{15.16} \\
\hline
อัตราเฟรมของกล้อง (เฟรมต่อวินาที) & 15.15 & --- \\
\hline
จำนวนวัตถุที่ตรวจพบเฉลี่ยต่อเฟรม & 4.41 ชิ้น & 1.74 ชิ้น \\
\hline
สัดส่วนเฟรมที่มีผลการตรวจจับแต่ไม่ได้รับหมายเลขติดตาม & ร้อยละ 0 & --- \\
\hline
ความล่าช้าตั้งแต่ถ่ายภาพจนผู้รับได้ข้อมูล (ค่ามัธยฐาน) & 65 มิลลิวินาที & 57 มิลลิวินาที \\
\hline
อัตราการใช้งานหน่วยประมวลผลกราฟิกเฉลี่ย & ร้อยละ 57 & ร้อยละ 51 \\
\hline
อุณหภูมิสูงสุดของชิปประมวลผล & 52.0 องศาเซลเซียส & 54.6 องศาเซลเซียส \\
\hline
\end{tabularx}
\end{table}

ข้อสรุปคือ \textbf{ระบบถูกจำกัดด้วยอัตราเฟรมของกล้อง ไม่ใช่ด้วยความเร็วของการอนุมาน}
โดยโหนดประมวลผลได้ทุกเฟรมที่กล้องส่งมาโดยไม่มีการทิ้งเฟรม
และอัตรา 15.13 เฟรมต่อวินาทีเมื่อทำงานลำพังสูงกว่าเกณฑ์ 5 เฟรมต่อวินาทีที่กำหนดไว้ในวัตถุประสงค์ถึงสามเท่า

\textbf{เกณฑ์ในวัตถุประสงค์กำหนดให้วัดขณะที่ SLAM ทำงานอยู่ด้วย}
สำหรับรุ่นแรก จึงวัดซ้ำจากข้อมูลที่บันทึกด้วย rosbag ขณะระบบทำงานครบทุกส่วนเป็นเวลา 215.3 วินาที
ได้อัตราการเผยแพร่ผลการตรวจจับ \textbf{13.05 เฟรมต่อวินาที}
โดยในช่วงเวลาเดียวกัน ข้อมูลการสแกนของไลดาร์ยังเผยแพร่ที่ 11.60 ครั้งต่อวินาที
และแผนที่ถูกปรับปรุงต่อเนื่อง ซึ่งยืนยันว่ากระบวนการ SLAM ทำงานอยู่จริง
ส่วนต่างประมาณ 2 เฟรมต่อวินาทีจากการทำงานลำพังคือภาระของ SLAM ชั้นผสานข้อมูล และตัวบันทึกข้อมูล
ทั้งนี้ค่าที่อ่านจาก rosbag เป็นจำนวนข้อความที่ถูกบันทึกได้
จึงเป็นขอบล่างของอัตราที่ระบบเผยแพร่จริง

\textbf{สำหรับ YOLO26l บน TensorRT ที่ใช้งานจริง วัดขณะ SLAM ทำงานอยู่โดยตรง}
ได้อัตรา \textbf{15.16 เฟรมต่อวินาที} (1{,}787 ข้อความใน 120 วินาที ช่วงห่างเฉลี่ย 66.0 มิลลิวินาที)
ซึ่งเท่ากับอัตราเฟรมของกล้อง และสูงกว่าเกณฑ์ประมาณ 3 เท่า
ขณะที่ข้อมูลการสแกนของไลดาร์เผยแพร่ที่ 11.5 ครั้งต่อวินาที
เมื่อทำงานลำพังก่อนหน้านั้นได้ 15.1--15.2 เฟรมต่อวินาทีเช่นกัน
โดยเวลาอนุมานรวมการติดตามวัตถุมีค่ามัธยฐาน 41--51 มิลลิวินาที ยังน้อยกว่าคาบของกล้อง (66.7 มิลลิวินาที)
แม้แบบจำลองจะใหญ่ขึ้นจากรุ่น Small เป็นรุ่น Large
ทั้งนี้การวัดสองรุ่นทำคนละเงื่อนไข รอบของ YOLO26l ไม่ได้เปิดระบบนำทางและตัวบันทึกข้อมูล
และหุ่นยนต์จอดนิ่ง ผลจึงยืนยันว่าผ่านเกณฑ์ ไม่ใช่ว่า YOLO26l เร็วกว่า YOLO26s

\textbf{ข้อควรระวังในการตีความผลการวัดคือการเปรียบเทียบระหว่างรอบ
มีตัวแปรเปลี่ยนไปพร้อมกันหลายตัว} ทั้งรูปแบบแบบจำลองและความหนาแน่นของวัตถุในฉาก
(4.41 เทียบกับ 1.74 ชิ้นต่อเฟรม)
ค่าการใช้งานหน่วยประมวลผลกราฟิกและอุณหภูมิที่ต่างกันจึงไม่สามารถสรุปว่าเกิดจากแบบจำลองเพียงอย่างเดียว
เพราะขั้นตอนหลังการอนุมานและการติดตามวัตถุล้วนแปรผันตามจำนวนวัตถุที่ตรวจพบ
การเปรียบเทียบทรัพยากรของแบบจำลองในเงื่อนไขเดียวกันแสดงไว้ในรูปที่ \ref{fig:infer-bench}
และตารางที่ \ref{tab:det-cost}

\textbf{ผลการวัดอีกข้อหนึ่งที่ต้องรายงานไว้คือพฤติกรรมของหมายเลขติดตามวัตถุ}
ตลอดการทดสอบ 301 วินาทีของรุ่นแรก ระบบสร้างหมายเลขติดตามขึ้นทั้งสิ้น 144 หมายเลข
โดยวัตถุที่อยู่นิ่งและเห็นชัดรักษาหมายเลขเดิมไว้ได้สมบูรณ์
(คอมพิวเตอร์โน้ตบุ๊กหมายเลข 37 ปรากฏใน 4{,}459 จาก 4{,}553 เฟรม คิดเป็นร้อยละ 100
ของช่วงที่มองเห็น) แต่มีหลายคลาสที่ได้รับหมายเลขมากกว่าหนึ่งหมายเลข
เช่น เก้าอี้ได้หมายเลข 329 และ 230
\textbf{ผลการวัดนี้ไม่สามารถแยกได้ว่า ``มีเก้าอี้สองตัวในห้อง'' หรือ ``มีเก้าอี้ตัวเดียวแต่ได้สองหมายเลข''}
ซึ่งเป็นประเด็นสำคัญเพราะชั้นผสานข้อมูลใช้หมายเลขติดตามเป็นกุญแจหลัก
ในการตัดสินว่าวัตถุที่เห็นเป็นวัตถุเดิมหรือไม่
หากวัตถุชิ้นเดียวได้รับหมายเลขใหม่ ระบบจะบันทึกเป็นวัตถุสองชิ้นบนแผนที่
ข้อจำกัดนี้ถูกบันทึกไว้อีกครั้งในหัวข้อข้อจำกัดของระบบในบทที่ 5


\section{ระบบผสานข้อมูลและเว็บแอปพลิเคชันแสดงผล}

\subsection{การแปลงพิกัดวัตถุจากภาพสู่แผนที่}

โมดูลผสานข้อมูลรับผลการตรวจจับจากโหนดตรวจจับวัตถุ
แล้วคำนวณพิกัดของวัตถุบนแผนที่ตามลำดับขั้นดังนี้

\begin{enumerate}[label=\arabic*., leftmargin=1.25cm, itemsep=4pt, topsep=6pt]
    \item พิจารณาเฉพาะ \textbf{ขอบซ้ายและขอบขวา} ของกรอบสี่เหลี่ยม
    เนื่องจากระบบนี้เป็นการผสานข้อมูลในระนาบสองมิติ พิกัดแนวตั้งของกรอบจึงไม่ถูกใช้
    ระบบไม่ได้ใช้จุดกึ่งกลางด้านล่างของกรอบเป็นตำแหน่งของวัตถุบนพื้น
    ซึ่งเป็นวิธีที่ต้องอาศัยสมมติฐานว่าพื้นเรียบสนิทและกล้องมีมุมก้มที่ทราบค่าแน่นอน

    \item แปลงตำแหน่งพิกเซลของขอบทั้งสองเป็นมุมกวาดในเฟรมกล้อง
    โดยใช้พารามิเตอร์ภายในของกล้องที่ได้จากการสอบเทียบด้วยกระดานหมากรุก
    ผลลัพธ์คือ \textbf{ช่วงมุม} ที่วัตถุครอบคลุม ไม่ใช่ทิศทางเดียว

    \item แปลงช่วงมุมดังกล่าวจากเฟรมกล้องไปเป็นช่วงมุมในเฟรมของหุ่นยนต์
    โดยคิดทั้งมุมติดตั้งของกล้องและระยะเยื้องระหว่างกล้องกับไลดาร์
    แล้วรวบรวม \textbf{ทุกลำแสงของไลดาร์ที่ตกอยู่ในช่วงมุมนั้น}

    \item ตรวจสอบคุณภาพของกลุ่มลำแสงที่ได้ก่อนนำไปใช้ ด้วยเงื่อนไขสองข้อ คือ
    ต้องมีลำแสงที่วัดค่าได้ไม่น้อยกว่าจำนวนที่กำหนด (\texttt{min\_returns})
    และค่าระยะทางของลำแสงเหล่านั้นต้องกระจายตัวไม่เกินเกณฑ์ที่กำหนด
    (\texttt{max\_spread})
    \textbf{หากไม่ผ่านเงื่อนไขใดเงื่อนไขหนึ่ง ระบบจะปฏิเสธการตรวจจับนั้นทิ้งไป
    ไม่ใช้วิธีเดาระยะทางจากขนาดของกรอบแทน} เพราะการกระจายตัวที่มากผิดปกติ
    หมายความว่าช่วงมุมนั้นคร่อมทั้งตัววัตถุและฉากหลังพร้อมกัน
    ค่าที่ได้จึงไม่ใช่ระยะทางถึงวัตถุ

    \item แปลงพิกัดที่ได้จากเฟรมหุ่นยนต์ไปยังเฟรมแผนที่
    โดยสอบถามการแปลงพิกัดจาก TF2 ที่ประทับเวลาตรงกับเวลาที่ถ่ายภาพ

    \item ตรวจสอบว่าวัตถุนี้เคยถูกบันทึกไว้แล้วหรือไม่
    โดยพิจารณาจากหมายเลขติดตามวัตถุ (Track ID) เป็นลำดับแรก
    หากเป็นวัตถุที่เพิ่งปรากฏหรือหมายเลขติดตามขาดช่วง
    จึงค่อยพิจารณาจากระยะทางถึงวัตถุคลาสเดียวกันที่บันทึกไว้แล้ว
    เพื่อป้องกันการบันทึกซ้ำจากการมองเห็นวัตถุเดียวกันหลายเฟรม
\end{enumerate}

ขั้นตอนการแปลงพิกัดตามที่อธิบายข้างต้นแสดงเป็นแผนภาพดังรูปที่ \ref{fig:bearing-range}

\begin{figure}[htbp]
\centering
\includegraphics[width=\textwidth]{images/bearing_range_projection.png}
\caption[การแปลงขอบของกรอบภาพเป็นช่วงมุมและการเลือกลำแสงไลดาร์ (บทที่ 4)]{การแปลงขอบของกรอบภาพเป็นช่วงมุมและการเลือกลำแสงไลดาร์}
\label{fig:bearing-range}
\end{figure}

ในเชิงคณิตศาสตร์ การแปลงตำแหน่งพิกเซลแนวนอน $u$ เป็นมุมกวาดในเฟรมกล้อง
ใช้แบบจำลองกล้องรูเข็มดังนี้

\begin{equation}
\alpha(u) = \tan^{-1}\!\left(\frac{u - c_x}{f_x}\right)
\end{equation}

โดยมุมที่ได้มีค่าเป็นบวกเมื่อวัตถุอยู่ทางขวาของภาพ
ซึ่งตรงข้ามกับระเบียบของ ROS ที่กำหนดให้ทิศทางทวนเข็มนาฬิกา (ด้านซ้าย) เป็นบวก
การแปลงมุมดังกล่าวไปเป็นมุมในเฟรมของหุ่นยนต์จึงต้องกลับเครื่องหมาย
แล้วบวกด้วยมุมติดตั้งของกล้อง $\psi$ ดังนี้

\begin{equation}
\beta = -\alpha + \psi
\end{equation}

จากนั้นตำแหน่งของวัตถุในเฟรมหุ่นยนต์คือจุดตัดระหว่างรังสีสายตาที่ออกจากจุดกำเนิดของกล้อง
$\mathbf{C}$ ในทิศทาง $\mathbf{d} = (\cos\beta,\ \sin\beta)$
กับวงกลมรัศมี $r$ รอบจุดกำเนิดของไลดาร์ $\mathbf{L}$
เมื่อให้ $\mathbf{m} = \mathbf{C} - \mathbf{L}$ จะได้สมการกำลังสอง

\begin{equation}
t^2 + 2\,(\mathbf{m}\cdot\mathbf{d})\,t + |\mathbf{m}|^2 - r^2 = 0
\end{equation}

และเลือกรากที่ไกลกว่าเสมอ

\begin{equation}
t = -\,\mathbf{m}\cdot\mathbf{d} + \sqrt{(\mathbf{m}\cdot\mathbf{d})^2 - |\mathbf{m}|^2 + r^2},
\qquad
\mathbf{P}_{\text{base}} = \mathbf{C} + t\,\mathbf{d}
\end{equation}

หากค่าภายใต้รากที่สองติดลบ แสดงว่าระยะทางที่วัดได้สั้นกว่าที่เรขาคณิตของการติดตั้งจะเป็นไปได้
ระบบจะปฏิเสธการตรวจจับนั้น และเมื่อกล้องกับไลดาร์อยู่ตำแหน่งเดียวกัน
($\mathbf{C} = \mathbf{L}$) สมการนี้จะลดรูปกลับไปเป็น $\mathbf{C} + r\,\mathbf{d}$ พอดี
สุดท้ายจึงแปลงพิกัดจากเฟรมหุ่นยนต์ไปยังเฟรมแผนที่ด้วยท่าทางของหุ่นยนต์
$(x_{r}, y_{r}, \theta_{r})$ ที่สอบถามจาก TF2 ณ เวลาที่ถ่ายภาพ

\begin{equation}
\begin{bmatrix} x_{\text{map}} \\ y_{\text{map}} \end{bmatrix}
=
\begin{bmatrix} \cos\theta_{r} & -\sin\theta_{r} \\
                \sin\theta_{r} & \cos\theta_{r} \end{bmatrix}
\mathbf{P}_{\text{base}}
+
\begin{bmatrix} x_{r} \\ y_{r} \end{bmatrix}
\end{equation}

การสอบเทียบพารามิเตอร์ภายในของกล้องที่ใช้ในขั้นตอนที่ 2 แสดงดังรูปที่ \ref{fig:cam-calib}

\begin{figure}[htbp]
\centering
\includegraphics[height=10cm,keepaspectratio]{images/camera_calibration.jpg}
\caption[การสอบเทียบพารามิเตอร์ภายในของกล้องด้วยกระดานหมากรุก (บทที่ 4)]{การสอบเทียบพารามิเตอร์ภายในของกล้องด้วยกระดานหมากรุก}
\label{fig:cam-calib}
\end{figure}

การใช้ประทับเวลาในการสอบถามการแปลงพิกัดมีความสำคัญมาก
เพราะหากใช้การแปลงพิกัด ณ เวลาปัจจุบันแทนเวลาที่ถ่ายภาพ
ตำแหน่งของวัตถุจะคลาดเคลื่อนตามระยะทางที่หุ่นยนต์เคลื่อนที่ไประหว่างการอนุมาน

ค่าพารามิเตอร์ภายในของกล้องที่ได้จากการสอบเทียบและนำไปใช้จริง
แสดงดังตารางที่ \ref{tab:camcalib}

\begin{table}[htbp]
\centering
\caption[ผลการสอบเทียบพารามิเตอร์ภายในของกล้อง (บทที่ 4)]{ผลการสอบเทียบพารามิเตอร์ภายในของกล้อง}
\label{tab:camcalib}
\small
\begin{tabularx}{\textwidth}{|X|c|}
\hline
\textbf{รายการ} & \textbf{ค่าที่ได้} \\
\hline
ความละเอียดของภาพ & 640$\times$480 พิกเซล \\
\hline
ความยาวโฟกัสแนวนอน $f_x$ & 667.874 พิกเซล \\
\hline
ความยาวโฟกัสแนวตั้ง $f_y$ & 669.846 พิกเซล \\
\hline
จุดศูนย์กลางภาพ $c_x$ & 321.569 พิกเซล \\
\hline
จุดศูนย์กลางภาพ $c_y$ & 234.502 พิกเซล \\
\hline
ความคลาดเคลื่อนของการฉายกลับ & 0.3403 พิกเซล \\
\hline
มุมมองภาพแนวนอนที่วัดได้ & 51.2$^\circ$ \\
\hline
กระดานหมากรุกที่ใช้ & 9$\times$6 ช่อง ขนาดช่อง 20 มิลลิเมตร \\
\hline
จำนวนภาพที่ใช้ในการคำนวณ & 59 ภาพ \\
\hline
\end{tabularx}
\end{table}

ประเด็นที่ต้องบันทึกไว้จากการสอบเทียบมีสองข้อ
ข้อแรกคือ \textbf{ต้องล็อกระบบโฟกัสอัตโนมัติของกล้องไว้ที่ค่าเดียว}
และต้องเป็นค่าเดียวกับที่ใช้ขณะทำงานจริง
เพราะการเปลี่ยนระยะโฟกัสทำให้ความยาวโฟกัสที่สอบเทียบไว้ใช้ไม่ได้
ข้อที่สองคือ \textbf{มุมมองภาพที่วัดได้จริงคือ 51.2$^\circ$
ซึ่งต่างจากค่าที่ผู้ผลิตระบุไว้อย่างมีนัยสำคัญ}
ค่าที่ผู้ผลิตระบุมักเป็นมุมมองตามแนวทแยงของโหมดภาพอัตราส่วน 16 ต่อ 9
ซึ่งใช้กับโหมดภาพอัตราส่วน 4 ต่อ 3 ที่ระบบนี้ใช้ไม่ได้
การใช้ค่าจากผู้ผลิตจะทำให้พิกัดวัตถุคลาดเคลื่อนอย่างเป็นระบบ

พารามิเตอร์ที่ควบคุมการยอมรับหรือปฏิเสธผลการตรวจจับแสดงดังตารางที่ \ref{tab:fuseparam}

\begin{table}[htbp]
\centering
\caption[พารามิเตอร์ของชั้นผสานข้อมูลและเงื่อนไขการยอมรับผลการตรวจจับ (บทที่ 4)]{พารามิเตอร์ของชั้นผสานข้อมูลและเงื่อนไขการยอมรับผลการตรวจจับ}
\label{tab:fuseparam}
\small
\begin{tabularx}{\textwidth}{|>{\raggedright\arraybackslash}p{4.4cm}|c|X|}
\hline
\textbf{พารามิเตอร์} & \textbf{ค่า} & \textbf{เหตุผล} \\
\hline
ค่าความเชื่อมั่นขั้นต่ำ & 0.5 & คัดผลการตรวจจับที่ไม่น่าเชื่อถือออกก่อนผสานข้อมูล \\
\hline
จำนวนลำแสงขั้นต่ำ & 3 ลำแสง & รองรับสัดส่วนลำแสงที่ไม่ได้ค่าร้อยละ 28 ที่ระยะ 3 เมตร \\
\hline
การกระจายตัวของระยะทางสูงสุด & 0.5 เมตร & ปฏิเสธกรณีที่ช่วงมุมคร่อมทั้งวัตถุและฉากหลัง \\
\hline
มุมเผื่อของช่วงมุม & 0.035 เรเดียน & ชดเชยระยะเยื้อง 3 เซนติเมตรระหว่างกล้องกับไลดาร์ \\
\hline
ความเร็วเชิงมุมสูงสุดที่ยอมผสานข้อมูล & 0.3 เรเดียนต่อวินาที & ไม่ผสานข้อมูลขณะหุ่นยนต์หมุน เพราะการสแกนหนึ่งรอบกินเวลา 86 มิลลิวินาที \\
\hline
อายุสูงสุดของข้อมูลโอดอมิเตอร์ & 0.5 วินาที & หากข้อมูลเก่ากว่านี้ ระบบปิดประตูการผสานข้อมูลทันที \\
\hline
ระยะรวมวัตถุเป็นชิ้นเดียวกัน & 0.5 เมตร & ป้องกันการบันทึกวัตถุเดียวกันซ้ำ \\
\hline
ค่าคงที่ของค่าเฉลี่ยเคลื่อนที่ & 0.3 & น้ำหนักของการสังเกตใหม่ต่อตำแหน่งที่บันทึกไว้ \\
\hline
จำนวนครั้งที่ต้องพบก่อนเผยแพร่ & 2 ครั้ง & ยืนยันวัตถุก่อนแสดงผล ลดผลบวกลวงชั่วคราว \\
\hline
ช่วงเวลาที่ยอมให้ข้อมูลสองแหล่งต่างกัน & 0.1 วินาที & จับคู่ข้อมูลการสแกนกับผลการตรวจจับที่ประทับเวลาใกล้กันที่สุด \\
\hline
\end{tabularx}
\end{table}

ประเด็นสำคัญเชิงเรขาคณิตที่ต้องระบุไว้คือ
\textbf{ระบบตัดเส้นสายตาจากกล้องกับวงกลมระยะทางของไลดาร์}
ไม่ใช่การวางจุดวัตถุที่ระยะทางนั้นนับจากตำแหน่งกล้อง
เพราะกล้องและไลดาร์ติดตั้งอยู่คนละตำแหน่งบนตัวหุ่นยนต์
การละเลยระยะเยื้องนี้ทำให้เกิดความคลาดเคลื่อนคงที่ประมาณ 9 เซนติเมตร
และระยะเยื้องด้านข้าง 3 เซนติเมตรคิดเป็นมุม 1.7$^\circ$ ที่ระยะ 1 เมตร
ซึ่งมากกว่าความละเอียดเชิงมุมของไลดาร์หนึ่งลำแสง
การตัดกันแบบตรงไปตรงมาให้ผลที่ถูกต้องสำหรับทั้งสองกรณีโดยใช้การคำนวณเพิ่มเพียงรากที่สองหนึ่งครั้ง

เงื่อนไขการกระจายตัวของระยะทางในขั้นตอนที่ 4 มีผลโดยตรงต่อระยะที่เหมาะสมในการวางวัตถุ
ผลการวัดสัดส่วนของการตรวจจับที่ผ่านเงื่อนไขและถูกนำไปบันทึกเป็นวัตถุบนแผนที่
ที่ระยะต่าง ๆ ของวัตถุชิ้นเดียวกัน (เก้าอี้) แสดงดังตารางที่ \ref{tab:fuserange}

\begin{table}[htbp]
\centering
\caption[สัดส่วนการผสานข้อมูลสำเร็จที่ระยะต่าง ๆ ของวัตถุ (บทที่ 4)]{สัดส่วนการผสานข้อมูลสำเร็จที่ระยะต่าง ๆ ของวัตถุ}
\label{tab:fuserange}
\small
\begin{tabularx}{\textwidth}{|c|c|X|}
\hline
\textbf{ระยะของวัตถุ} & \textbf{สัดส่วนที่ผสานสำเร็จ} & \textbf{หมายเหตุ} \\
\hline
ประมาณ 0.7 เมตร & ร้อยละ 0 & ถูกปฏิเสธทั้ง 57 จาก 57 เฟรม
  การกระจายตัวของระยะทางที่วัดได้คือ 0.62 เมตร เทียบกับเกณฑ์ 0.5 เมตร \\
\hline
ระยะขณะขับผ่าน & ประมาณร้อยละ 45 & วัดจากสี่ช่วงเวลา ช่วงละ 5 วินาที
  ขณะหุ่นยนต์อยู่กับที่เป็นส่วนใหญ่ การปฏิเสธทุกครั้งเกิดจากเงื่อนไขการกระจายตัว \\
\hline
1.64 เมตร & ร้อยละ 94.3 & การทดสอบที่จุดเดียวขณะหุ่นยนต์หยุดนิ่ง วัดระยะด้วยตลับเมตร \\
\hline
\end{tabularx}
\end{table}

ที่ระยะประมาณ 0.7 เมตร เก้าอี้กินมุมมองภาพประมาณ 39 องศา
ช่วงมุมที่ใช้เลือกลำแสงไลดาร์จึงคร่อมทั้งตัวเก้าอี้และฉากหลังอย่างหลีกเลี่ยงไม่ได้
\textbf{การที่ระบบปฏิเสธการตรวจจับทั้งหมดในกรณีนี้คือการทำงานที่ถูกต้องตามที่ออกแบบไว้
ไม่ใช่ข้อบกพร่อง} และไม่ใช่เหตุผลที่จะลดเกณฑ์การกระจายตัวลง
วิธีแก้คือการวางวัตถุให้ไกลขึ้นหรือถอยหุ่นยนต์ออกมา

นอกจากนี้ระบบอ่านตำแหน่งและมุมติดตั้งของกล้องจากระบบพิกัดที่ประกาศไว้ในไฟล์ URDF
แทนการกำหนดเป็นค่าคงที่ในไฟล์พารามิเตอร์
เพื่อให้มีแหล่งข้อมูลเพียงแหล่งเดียวและป้องกันไม่ให้ค่าทั้งสองที่คลาดเคลื่อนจากกัน


\subsection{ส่วนหลังและการจัดเก็บข้อมูล}

ส่วนหลังพัฒนาด้วย FastAPI ทำหน้าที่เป็นสะพานเชื่อมระหว่างระบบ ROS 2 กับเบราว์เซอร์
โดยเปิดปลายทาง REST API สำหรับการเรียกข้อมูลแบบครั้งต่อครั้ง ได้แก่
การตรวจสอบสถานะระบบ การเรียกรายการวัตถุที่พบทั้งหมด การเรียกแผนที่ปัจจุบัน
และการสั่งล้างรายการวัตถุ และเปิดช่องทาง WebSocket สองช่อง
สำหรับส่งสถานะของระบบกับตำแหน่งหุ่นยนต์ และสำหรับส่งภาพจากกล้องแบบต่อเนื่อง

ข้อมูลวัตถุที่ระบบยืนยันแล้วถูกบันทึกลงไฟล์ JSON ไฟล์เดียว
ซึ่งเก็บชื่อคลาสของวัตถุ พิกัดบนแผนที่ ค่าความเชื่อมั่น จำนวนครั้งที่พบ
และเวลาที่พบล่าสุด \textbf{ระบบไม่ได้ใช้ฐานข้อมูล}
เพราะจำนวนวัตถุที่พบในการสำรวจหนึ่งรอบอยู่ในระดับหลักสิบ
ซึ่งไม่คุ้มกับความซับซ้อนที่ระบบฐานข้อมูลเพิ่มเข้ามา

ข้อจำกัดที่ต้องระบุไว้คือไฟล์นี้ถูกเขียนทับใหม่ในแต่ละรอบการทำงาน และ
\textbf{ไม่ถูกอ่านกลับขึ้นมาใช้เมื่อเปิดระบบครั้งถัดไป}
เนื่องจากระบบสร้างแผนที่ใหม่ทุกครั้งที่เริ่มทำงาน
ตำแหน่งที่บันทึกไว้จากรอบก่อนจึงอ้างอิงกับเฟรมแผนที่คนละชุดกันและไม่สามารถนำมาใช้ต่อได้

โครงสร้างของข้อมูลที่บันทึกลงไฟล์แสดงดังรูปที่ \ref{fig:json-schema}

\begin{figure}[htbp]
\centering
\includegraphics[width=0.85\textwidth]{images/landmark_json_schema.png}
\caption[โครงสร้างข้อมูลของไฟล์ JSON ที่ใช้บันทึกวัตถุที่พบ (บทที่ 4)]{โครงสร้างข้อมูลของไฟล์ JSON ที่ใช้บันทึกวัตถุที่พบ}
\label{fig:json-schema}
\end{figure}

\subsection{ส่วนติดต่อผู้ใช้}

ส่วนติดต่อผู้ใช้พัฒนาด้วย React.js เป็นเว็บแอปพลิเคชันหน้าเดียว
หน้าจอหลักแสดงแผนที่ที่สร้างจาก SLAM พร้อมสัญลักษณ์ตำแหน่งและทิศทางของหุ่นยนต์
และหมุดตำแหน่งของวัตถุที่ตรวจพบซึ่งแยกสีตามชนิดของวัตถุ
แผงด้านข้างแสดงภาพจากกล้องพร้อมกรอบผลการตรวจจับ ตัวบ่งชี้สถานะการเชื่อมต่อ
ตารางรายการวัตถุที่พบ และปุ่มล้างรายการวัตถุ
\textbf{ขอบเขตที่ต้องระบุไว้คือเว็บแอปพลิเคชันทำหน้าที่แสดงผลและล้างรายการวัตถุเท่านั้น}
ไม่ได้ทำหน้าที่สั่งเริ่มหรือหยุดการสำรวจ
การสั่งเคลื่อนที่ทำผ่านคำสั่งจากแป้นพิมพ์ การกำหนดเป้าหมายให้ระบบนำทาง
หรือคำสั่ง \texttt{make explore} สำหรับการสำรวจอัตโนมัติ
และการบันทึกแผนที่ทำผ่านบริการของ slam\_toolbox โดยตรง

หน้าจอของเว็บแอปพลิเคชันขณะระบบทำงานจริงแสดงดังรูปที่ \ref{fig:ui-dashboard}
โดยแผงด้านซ้ายคือแผนที่ที่สร้างจาก SLAM พร้อมสัญลักษณ์ตำแหน่งและทิศทางของหุ่นยนต์
ขอบเขตการมองเห็นของกล้อง และหมุดตำแหน่งวัตถุที่แยกสีตามชนิด
ส่วนแผงด้านขวาเรียงจากบนลงล่างคือ ภาพจากกล้องพร้อมกรอบผลการตรวจจับซึ่งกำกับด้วยหมายเลขติดตาม
ชื่อคลาส และค่าความเชื่อมั่น ถัดมาคือแถบสถานะการเชื่อมต่อและจำนวนวัตถุที่ยืนยันแล้ว
และล่างสุดคือรายการวัตถุที่พบซึ่งแสดงชนิด พิกัดบนแผนที่ ค่าความเชื่อมั่น และจำนวนครั้งที่พบ

\begin{figure}[htbp]
\centering
\includegraphics[width=\textwidth]{images/ui.jpg}
\caption[หน้าจอของเว็บแอปพลิเคชันขณะระบบทำงาน (บทที่ 4)]{หน้าจอของเว็บแอปพลิเคชันขณะระบบทำงาน}
\label{fig:ui-dashboard}
\end{figure}

การเชื่อมต่อระหว่างเบราว์เซอร์กับหุ่นยนต์ใช้เครือข่ายไร้สายภายในพื้นที่
ซึ่งต้องระบุให้ชัดว่าเป็นการเชื่อมต่อเพื่อการแสดงผลเท่านั้น
การประมวลผล SLAM และการอนุมานแบบจำลองทั้งหมดยังคงเกิดขึ้นภายในตัวอุปกรณ์
ระบบจึงยังทำงานได้ครบถ้วนแม้ไม่มีการเชื่อมต่ออินเทอร์เน็ต

% ---------------------------------------------------------------------
\chapter{สรุปผลการดำเนินงาน}

\section{สรุปผลการทำโครงงาน}

โครงงานนี้ได้พัฒนาหุ่นยนต์เคลื่อนที่อัตโนมัติแบบขับเคลื่อนสองล้ออิสระ
ที่สามารถสำรวจพื้นที่ปิดภายในอาคาร สร้างแผนที่สองมิติด้วยไลดาร์
ตรวจจับวัตถุด้วยแบบจำลองการเรียนรู้เชิงลึกที่ทำงานบนอุปกรณ์เอดจ์เอไอ
และระบุพิกัดของวัตถุที่ตรวจพบลงบนแผนที่
โดยการประมวลผลทั้งหมดเกิดขึ้นภายในตัวอุปกรณ์

ผลการดำเนินงานแยกตามส่วนของระบบสรุปได้ดังนี้

\begin{itemize}[leftmargin=1.25cm, itemsep=4pt, topsep=6pt]
    \item \textbf{ชั้นควบคุมการเคลื่อนที่} เฟิร์มแวร์บน ESP32 ร่วมกับ Hardware Interface
    ของ ros2\_control สามารถรับคำสั่งความเร็วและอ่านค่าเอนโคเดอร์ได้ตามที่ออกแบบ
    คอนโทรลเลอร์ \texttt{diff\_cont} คำนวณและเผยแพร่โอดอมิเตอร์
    พร้อมกระจายการแปลงพิกัดได้ที่อัตรา 30 ครั้งต่อวินาที
    ระยะห่างระหว่างล้อที่ได้จากการสอบเทียบ (0.25168 เมตร)
    ได้รับการยืนยันซ้ำด้วยข้อมูลเอนโคเดอร์ขณะหมุนอยู่กับที่ โดยต่างกันร้อยละ 0.003

    \item \textbf{ชั้นสร้างแผนที่และนำทาง} ระบบสร้างแผนที่สองมิติของพื้นที่ทดสอบได้สำเร็จ
    ด้วย slam\_toolbox โหมดออนไลน์แบบอะซิงโครนัส
    แผนที่อ้างอิงจากการทดสอบวนกลับมาที่เดิม (15 ก.ย.) มีขนาด 271$\times$488 ช่อง
    ที่ความละเอียด 0.05 เมตรต่อช่อง คิดเป็นพื้นที่ 13.6$\times$24.4 เมตร
    และแผนที่จากรอบการสาธิต (22 ก.ย.) ตามรูปที่ \ref{fig:map-result} มีขนาด 666$\times$361 ช่อง
    คิดเป็นกรอบขนาด 33.3$\times$18.1 เมตร
    ระบบนำทางสามารถวางแผนเส้นทางและเคลื่อนที่ไปยังเป้าหมายที่กำหนดบนแผนที่ที่สร้างขึ้นได้สำเร็จ
    รวมถึงพฤติกรรมกู้คืนสถานการณ์ด้วยการหมุนตัวที่ทำงานได้ตามที่ออกแบบ
    และสามารถสำรวจพื้นที่ได้เองด้วยการเลือกเป้าหมายที่ขอบเขตของพื้นที่ที่ยังไม่ทราบข้อมูล
    (หัวข้อ \ref{sec:explore})
    ทั้งนี้ความคลาดเคลื่อนเชิงตำแหน่งเทียบกับจุดอ้างอิงที่ทราบพิกัดจริงสามจุด
    มีค่าสูงสุด \textbf{6.4 เซนติเมตร} ผ่านเกณฑ์ 10 เซนติเมตร
    (รายละเอียดในตารางที่ \ref{tab:obj1})

    \item \textbf{ชั้นตรวจจับวัตถุ} แบบจำลองรุ่นแรก YOLO26s รูปแบบ ONNX ความละเอียด 16 บิต
    ทำงานบนหน่วยประมวลผลกราฟิกของอุปกรณ์ได้ที่อัตรา 15.13 เฟรมต่อวินาที
    ตลอดการทดสอบต่อเนื่อง 301 วินาทีเมื่อทำงานลำพัง โดยประมวลผลได้ครบทุกเฟรมที่กล้องส่งมา
    และได้ 13.05 เฟรมต่อวินาทีขณะทำงานพร้อมกับ SLAM และส่วนอื่นของระบบทั้งหมด
    ซึ่งสูงกว่าเกณฑ์ 5 เฟรมต่อวินาทีที่กำหนดไว้ประมาณ 2.6 เท่า
    และระบบถูกจำกัดด้วยอัตราเฟรมของกล้อง ไม่ใช่ด้วยความเร็วของการอนุมาน
    ด้านความแม่นยำ ประเมินบนภาพจากพื้นที่ใช้งานจริง 143 ภาพที่กำกับข้อมูลด้วยมือ
    YOLO26s ที่ติดตั้งใช้งานเดิมได้ค่า F1 เฉลี่ยร้อยละ 64.4
    โดยความแม่นยำสูง (ร้อยละ 96.5) แต่ตรวจไม่พบวัตถุระยะไกลเป็นส่วนใหญ่
    หลังปรับปรุงเป็น YOLO26l บน TensorRT ที่รับภาพขนาด 480$\times$640 ตรงกับกล้อง
    และลดค่าความเชื่อมั่นขั้นต่ำเป็น 0.4 ได้ค่า F1 เฉลี่ยร้อยละ \textbf{82.2} ผ่านเกณฑ์ร้อยละ 80
    (รายละเอียดในหัวข้อ \ref{sec:eval-det})
    หลังปรับปรุง YOLO26l บน TensorRT ยังประมวลผลได้ 15.16 เฟรมต่อวินาที
    ขณะทำงานพร้อมกับ SLAM และชั้นผสานข้อมูล ซึ่งเท่ากับอัตราเฟรมของกล้อง

    \item \textbf{ชั้นผสานข้อมูล} การวัดตำแหน่งเก้าอี้จากสี่ทิศทางรอบวัตถุ
    (หน้า ขวา หลัง ซ้าย) เทียบกับระยะที่วัดด้วยตลับเมตร ที่ระยะ 1.46--1.70 เมตร
    ให้ความคลาดเคลื่อนเฉลี่ย 0.19 เมตร และสูงสุด 0.34 เมตร (ด้านหลัง)
    ผ่านเกณฑ์ 50 เซนติเมตรที่กำหนดไว้ในวัตถุประสงค์ทั้งสี่ทิศทาง
    และเก้าอี้ถูกบันทึกเป็นวัตถุเพียงรายการเดียวตลอดการวัดทั้งสี่ทิศทาง
    ทั้งนี้ \textbf{หุ่นยนต์หยุดนิ่งขณะวัดในแต่ละทิศทาง}
    ไม่ได้วัดขณะหุ่นยนต์เคลื่อนที่
    ส่วนการทดสอบก่อนหน้าที่จุดเดียว ระยะ 1.64 เมตร ทำมุม 18.4 องศาจากแกนกลางของกล้อง
    ซึ่งใช้ตรวจโซ่การคำนวณเชิงเรขาคณิต ให้ความคลาดเคลื่อน 0.08 เมตร
    และมีสัดส่วนการตรวจจับที่ผสานข้อมูลสำเร็จร้อยละ 94.3

    \item \textbf{ชั้นแสดงผล} เว็บแอปพลิเคชันแสดงแผนที่ ตำแหน่งหุ่นยนต์
    และตำแหน่งวัตถุที่พบได้แบบเรียลไทม์ผ่านเบราว์เซอร์
    โดยผู้ใช้ไม่ต้องติดตั้งโปรแกรมเพิ่มเติม

    \item \textbf{การใช้ทรัพยากรประมวลผล} ขณะเปิดทั้งระบบและขับหุ่นยนต์
    อัตราการใช้งานหน่วยประมวลผลกลางเฉลี่ย 6 แกนเดิมอยู่ที่ร้อยละ 81.5 ซึ่งไม่ผ่านเกณฑ์
    สาเหตุหลักคือเธรดของไลบรารี OpenBLAS ที่วนรอระหว่างเฟรมในขั้นตอนติดตามวัตถุ
    หลังกำหนดให้ใช้เธรดเดียวค่าเฉลี่ยลดลงเหลือร้อยละ \textbf{55.8} ผ่านเกณฑ์ร้อยละ 80
    โดยหน่วยความจำสูงสุด 2,915 MB และอุณหภูมิชิปสูงสุด 53.1 องศาเซลเซียส
    (รายละเอียดในตารางที่ \ref{tab:obj5})
\end{itemize}

พื้นที่ที่ใช้ทดสอบระบบและการวางวัตถุทดสอบแสดงดังรูปที่ \ref{fig:test-area}

\begin{figure}[htbp]
\centering
\includegraphics[width=0.85\textwidth]{images/test_area.jpg}
\caption[พื้นที่ทดสอบและการวางวัตถุทดสอบ (บทที่ 5)]{พื้นที่ทดสอบภายในอาคาร และวัตถุทดสอบ (เก้าอี้) ที่วางไว้กลางพื้นที่}
\label{fig:test-area}
\end{figure}

\section{การประเมินคุณภาพของระบบงาน}
\label{sec:eval}

การประเมินคุณภาพของระบบดำเนินการโดยเทียบผลการทดสอบกับเกณฑ์เชิงตัวเลข
ที่กำหนดไว้ในวัตถุประสงค์ของโครงงาน สรุปดังตารางที่ \ref{tab:eval}

\begin{table}[htbp]
\centering
\caption[เกณฑ์และวิธีการประเมินคุณภาพของระบบ (บทที่ 5)]{เกณฑ์และวิธีการประเมินคุณภาพของระบบ}
\label{tab:eval}
\small
\begin{tabularx}{\textwidth}{|c|X|>{\raggedright\arraybackslash}p{3.3cm}|>{\raggedright\arraybackslash}p{2.7cm}|}
\hline
\textbf{ข้อ} & \textbf{เกณฑ์ที่กำหนด} & \textbf{วิธีการวัด} & \textbf{ผล} \\
\hline
1 & ความคลาดเคลื่อนเชิงตำแหน่งของ SLAM ไม่เกิน 10 เซนติเมตร
  & ทำเครื่องหมายจุดอ้างอิงที่ทราบพิกัดจริงในพื้นที่ทดสอบ
    ขับหุ่นยนต์กลับมายังจุดเดิม แล้วเทียบพิกัดที่ระบบรายงานกับพิกัดจริง
  & \textbf{ผ่าน} \newline สูงสุด 6.4 ซม. ที่ 3 จุด \newline (ไม่มีขากลับจุดเริ่มต้น) \\
\hline
2 & ความแม่นยำของแบบจำลองไม่น้อยกว่าร้อยละ 80
  & กำกับข้อมูลด้วยมือบนภาพ 143 ภาพจากพื้นที่ใช้งานจริง
    เทียบผลการตรวจจับที่ IoU 0.5 และรายงานค่า F1 เฉลี่ยของ 4 คลาสเป้าหมาย
  & \textbf{ผ่าน (หลังปรับปรุง)} \newline F1 ร้อยละ 82.2 \newline (เดิมร้อยละ 64.4) \\
\hline
3 & ความเร็วการประมวลผลไม่น้อยกว่า 5 เฟรมต่อวินาที
  & วัดอัตราเฟรมเฉลี่ยขณะรันจริงบนอุปกรณ์เป้าหมาย
    พร้อมกับที่ SLAM ทำงานอยู่
  & \textbf{ผ่าน} \newline YOLO26l: 15.16 เฟรมต่อวินาที \newline (YOLO26s: 13.05) \\
\hline
4 & ความคลาดเคลื่อนของตำแหน่งวัตถุไม่เกิน 50 เซนติเมตร
  & วางวัตถุทดสอบที่ทราบพิกัดจริง ให้หุ่นยนต์ตรวจจับจากหลายมุมมอง
    แล้วเทียบพิกัดที่ระบบบันทึกกับพิกัดจริง
  & \textbf{ผ่าน} \newline สูงสุด 0.34 เมตร \newline เฉลี่ย 0.19 เมตร \newline จาก 4 ทิศทาง \\
\hline
5 & อัตราการใช้หน่วยประมวลผลเฉลี่ยไม่เกินร้อยละ 80
  & บันทึกอัตราการใช้งานหน่วยประมวลผลตลอดรอบการสำรวจ
    ขณะที่ SLAM และการตรวจจับวัตถุทำงานพร้อมกัน
  & \textbf{ผ่าน (หลังปรับปรุง)} \newline เฉลี่ยร้อยละ 55.8 \newline (เดิมร้อยละ 81.5) \\
\hline
\end{tabularx}
\end{table}

รายละเอียดของการวัดตามเกณฑ์ข้อที่ 4 แสดงดังตารางที่ \ref{tab:obj4}
หุ่นยนต์จอดหันเข้าหาเก้าอี้ทดสอบจากสี่ทิศทาง ครั้งละ 30 วินาทีขณะหยุดนิ่ง
ระยะจริงวัดด้วยตลับเมตรจากจุดกึ่งกลางเพลาล้อในกรอบพิกัดของหุ่นยนต์
(ระยะตามแนวแกนกลาง และระยะตั้งฉากออกไปทางซ้าย)
ส่วนตำแหน่งที่ระบบรายงานคือตำแหน่งของวัตถุบนแผนที่ที่แปลงกลับมาเป็นกรอบพิกัดของหุ่นยนต์
ผ่านตำแหน่งของหุ่นยนต์บนแผนที่ ณ ขณะนั้น
ทำให้ความคลาดเคลื่อนของจุดกำเนิดแผนที่ไม่มีผลต่อการวัด

\begin{table}[htbp]
\centering
\caption[ผลการวัดตำแหน่งวัตถุจากสี่ทิศทาง (บทที่ 5)]{ผลการวัดตำแหน่งวัตถุ (เก้าอี้) จากสี่ทิศทาง เทียบกับระยะที่วัดด้วยตลับเมตร
(หน่วยเมตร ค่าบวกของระยะด้านข้างคือทางซ้ายของหุ่นยนต์)}
\label{tab:obj4}
\small
\begin{tabular}{|l|c|c|c|c|}
\hline
\textbf{ทิศทาง} & \textbf{ตลับเมตร (หน้า, ข้าง)} & \textbf{ระบบ (หน้า, ข้าง)} & \textbf{ความคลาดเคลื่อน} & \textbf{การกระจายตัวในรอบ} \\
\hline
หน้า & (1.60, 0.00) & (1.67, $-$0.02) & 0.07 & 0.16 \\
\hline
ขวา & (1.60, 0.00) & (1.50, $-$0.22) & 0.25 & 0.14 \\
\hline
หลัง & (1.46, 0.00) & (1.35, $-$0.32) & 0.34 & 0.11 \\
\hline
ซ้าย & (1.70, 0.12) & (1.78, 0.08) & 0.09 & 0.10 \\
\hline
\multicolumn{3}{|l|}{\textbf{เฉลี่ย / สูงสุด} (เกณฑ์ไม่เกิน 0.50 เมตร)} & \textbf{0.19 / 0.34} & --- \\
\hline
\multicolumn{3}{|l|}{การกระจายตัวของตำแหน่งบนแผนที่ระหว่างสี่ทิศทาง} & 0.28 & --- \\
\hline
\multicolumn{3}{|l|}{จำนวนรายการที่บันทึกสำหรับวัตถุจริงหนึ่งชิ้น} & 1 & --- \\
\hline
\end{tabular}
\end{table}

ทั้งสี่ทิศทางผ่านเกณฑ์ และเก้าอี้คงเป็นรายการเดียวกันตลอดการวัดทั้งสี่ทิศทาง
ความคลาดเคลื่อนส่วนใหญ่อยู่ในแนวด้านข้างมากกว่าแนวระยะ
(ด้านหลังคลาดไปทางขวา 0.32 เมตร ด้านขวาคลาด 0.22 เมตร)
ซึ่งยังไม่ได้วิเคราะห์หาสาเหตุ
ผลนี้ต่อยอดจากการทดสอบก่อนหน้าที่จุดเดียว ซึ่งออกแบบมาเพื่อตรวจโซ่การคำนวณเชิงเรขาคณิต
แสดงดังตารางที่ \ref{tab:benchcheck}

\begin{table}[htbp]
\centering
\caption[ผลการทดสอบตำแหน่งวัตถุที่จุดเดียว (บทที่ 5)]{ผลการทดสอบตำแหน่งวัตถุบนแผนที่ที่จุดเดียว เพื่อตรวจโซ่การคำนวณเชิงเรขาคณิต}
\label{tab:benchcheck}
\small
\begin{tabularx}{\textwidth}{|X|c|c|}
\hline
\textbf{รายการ} & \textbf{ค่าที่วัดได้} & \textbf{เกณฑ์} \\
\hline
ระยะจากกล้องถึงวัตถุ (วัดด้วยตลับเมตร) & 1.50 เมตร & --- \\
\hline
ระยะเยื้องด้านข้างของวัตถุ & 0.50 เมตร ทางขวา & --- \\
\hline
มุมของวัตถุเทียบกับแกนกลางของกล้อง & 18.4$^\circ$ & --- \\
\hline
ความคลาดเคลื่อนสัมบูรณ์ของตำแหน่งบนแผนที่ & \textbf{0.08 เมตร} & ไม่เกิน 0.50 เมตร \\
\hline
การกระจายตัวของตำแหน่งที่บันทึกได้ภายในรอบเดียวกัน & 0.04 เมตร & --- \\
\hline
จำนวนวัตถุที่ถูกบันทึกสำหรับวัตถุจริงหนึ่งชิ้น & 1 ชิ้น & 1 ชิ้น \\
\hline
สัดส่วนการตรวจจับที่ผสานข้อมูลสำเร็จ & 132/140 (ร้อยละ 94.3) & --- \\
\hline
\end{tabularx}
\end{table}

การที่ค่าความคลาดเคลื่อนอยู่ที่ 0.08 เมตรมีความหมายมากกว่าการผ่านเกณฑ์
เพราะการวางวัตถุให้ทำมุม 18.4 องศาจากแกนกลางของกล้องเป็นการออกแบบการทดลองโดยตั้งใจ
หากช่วงมุมที่ใช้ค้นหาลำแสงไลดาร์กลับด้านซ้ายขวา
ตำแหน่งที่ได้จะคลาดเคลื่อนไปประมาณ 1.06 เมตร
และหากใช้ตำแหน่งกล้องเป็นจุดเริ่มของระยะทางแทนตำแหน่งไลดาร์
จะเกิดความคลาดเคลื่อนคงที่ประมาณ 0.09 เมตร ซึ่งจะกินงบประมาณความคลาดเคลื่อนทั้งหมดไปเอง
ผลที่ได้จึงยืนยันว่าโซ่การคำนวณเชิงเรขาคณิตทั้งเส้นถูกต้อง
ตั้งแต่พารามิเตอร์ภายในของกล้อง ทิศทางของช่วงมุม ด้านที่ติดตั้งกล้อง
จุดเริ่มของระยะทาง ไปจนถึงการสอบถามการแปลงพิกัดที่ประทับเวลาของการถ่ายภาพ

รายละเอียดของการวัดตามเกณฑ์ข้อที่ 1 แสดงดังตารางที่ \ref{tab:obj1}
ทำเครื่องหมายจุดอ้างอิงสี่จุดบนพื้น ได้แก่ จุดเริ่มต้น (HOME) และจุด A, B, C
วัดพิกัดจริงด้วยตลับเมตรจากจุดเริ่มต้น
โดยแกน $x$ อยู่ตามทิศหัวของหุ่นยนต์ขณะเริ่มสร้างแผนที่ และแกน $y$ ไปทางซ้าย
ซึ่งตรงกับกรอบพิกัดของแผนที่ที่ slam\_toolbox สร้าง
ขับหุ่นยนต์ด้วยมือที่ความเร็ว 0.10 เมตรต่อวินาทีไปจอดที่แต่ละจุดตามลำดับ
โดยใช้โอดอมิเตอร์จากล้อเพียงอย่างเดียว (ไม่เปิดการผสานข้อมูลไจโรสโคป)
แล้วเฉลี่ยพิกัดที่ระบบรายงานขณะหยุดนิ่ง 5 วินาที
พิกัดจริงของแต่ละจุดคือตำแหน่งที่หุ่นยนต์จอดอยู่จริงตามที่วัดด้วยตลับเมตร
ที่จุด C หุ่นยนต์จอดเยื้องจากเครื่องหมาย 2 เซนติเมตร จึงใช้ค่าที่วัดได้ (0.00, 3.52) เมตร
จุดเริ่มต้นเป็นจุดกำเนิดของแผนที่ จึงมีความคลาดเคลื่อนเป็นศูนย์โดยนิยามและไม่นับเป็นจุดวัด

\begin{table}[htbp]
\centering
\caption[ผลการวัดความคลาดเคลื่อนเชิงตำแหน่งของ SLAM (บทที่ 5)]{ผลการวัดความคลาดเคลื่อนเชิงตำแหน่งของ SLAM
ที่จุดอ้างอิงสามจุด เทียบกับพิกัดที่วัดด้วยตลับเมตร (พิกัดหน่วยเมตร)}
\label{tab:obj1}
\small
\begin{tabular}{|c|c|c|c|}
\hline
\textbf{จุด} & \textbf{ตลับเมตร ($x$, $y$)} & \textbf{ระบบ ($x$, $y$)} & \textbf{ความคลาดเคลื่อน} \\
\hline
A & (4.85, 0.00) & (4.91, 0.03) & 6.4 ซม. \\
\hline
B & (4.85, 3.50) & (4.89, 3.50) & 4.1 ซม. \\
\hline
C & (0.00, 3.52) & (0.00, 3.47) & 5.3 ซม. \\
\hline
\multicolumn{3}{|l|}{\textbf{สูงสุด} (เกณฑ์ไม่เกิน 10 ซม.)} & \textbf{6.4 ซม.} \\
\hline
\multicolumn{3}{|l|}{สูงสุดหลังหักการหมุนระหว่างแกนตลับเมตรกับแกนแผนที่ ($-0.02^\circ$)} & 6.5 ซม. \\
\hline
\end{tabular}
\end{table}

ทั้งสามจุดผ่านเกณฑ์
ความคลาดเคลื่อนของระยะระหว่างจุดทุกคู่อยู่ระหว่าง $-5.3$ ถึง $+5.7$ เซนติเมตร
และมีทั้งค่าบวกและค่าลบ จึงไม่พบความคลาดเคลื่อนเชิงสเกลจากโอดอมิเตอร์
ค่าที่วัดได้รวมความคลาดเคลื่อนจากการจอดหุ่นยนต์ด้วยมือ (ประมาณ 2 เซนติเมตร)
และจากตลับเมตร (ประมาณ 1 เซนติเมตร) ไว้ด้วย จึงเป็นขอบเขตบนของความคลาดเคลื่อนของ SLAM

ทั้งนี้แบตเตอรี่หมดหลังจอดที่จุด C
จึงไม่ได้บันทึกขากลับจุดเริ่มต้น ซึ่งเป็นขาที่สะท้อนความคลาดเคลื่อนสะสมทั้งรอบ
นอกจากนี้ ในรอบทดสอบก่อนหน้าที่ใช้ปรับขั้นตอนการวัด
พบว่าความคลาดเคลื่อนเกินเกณฑ์เมื่อทิศหัวของหุ่นยนต์ขณะเริ่มสร้างแผนที่ไม่ตรงกับแกนตลับเมตร
และมีรอบหนึ่งที่ตำแหน่งกระโดดไปประมาณ 46 เซนติเมตรระหว่างจุด B กับ C ในรอบที่สอง

ผลการวัดตามเกณฑ์ข้อที่ 1 แสดงดังรูปที่ \ref{fig:slam-error}

\begin{figure}[htbp]
\centering
\includegraphics[width=\textwidth]{images/slam_error_result.png}
\caption[ผลการวัดความคลาดเคลื่อนเชิงตำแหน่งของ SLAM (บทที่ 5)]{ผลการวัดความคลาดเคลื่อนเชิงตำแหน่งของ SLAM
(ก) พิกัดจริงและพิกัดที่ระบบรายงานของจุดอ้างอิงบนแผนที่ วงกลมเส้นประคือรัศมีเกณฑ์ 10 เซนติเมตร
(ข) ความคลาดเคลื่อนของแต่ละจุดเทียบกับเกณฑ์ ทั้งค่าดิบและค่าหลังหักการหมุนของแกน}
\label{fig:slam-error}
\end{figure}

ผลการวัดตามเกณฑ์ข้อที่ 2 อธิบายแยกไว้ในหัวข้อ \ref{sec:eval-det}
พร้อมวิธีการวัดและขั้นตอนการปรับปรุงแบบจำลอง
ผลการวัดตามเกณฑ์ข้อที่ 3 แสดงดังรูปที่ \ref{fig:model-eval}

\begin{figure}[htbp]
\centering
\includegraphics[width=0.85\textwidth]{images/detection_fps.png}
\caption[ผลการวัดความเร็วของการตรวจจับวัตถุเทียบกับเกณฑ์ (บทที่ 5)]{อัตราการตรวจจับวัตถุที่วัดได้เทียบกับเกณฑ์ 5 เฟรมต่อวินาที ได้แก่ YOLO26s รูปแบบ ONNX เมื่อทำงานลำพังและเมื่อเปิดทั้งระบบพร้อมบันทึกข้อมูล และ YOLO26l บน TensorRT ขณะทำงานพร้อมกับ SLAM และชั้นผสานข้อมูล}
\label{fig:model-eval}
\end{figure}

ผลการวัดตามเกณฑ์ข้อที่ 4 แสดงดังรูปที่ \ref{fig:object-error}

\begin{figure}[htbp]
\centering
\includegraphics[width=\textwidth]{images/object_position_error.png}
\caption[ผลการวัดความคลาดเคลื่อนของตำแหน่งวัตถุบนแผนที่ (บทที่ 5)]{ผลการวัดความคลาดเคลื่อนของตำแหน่งวัตถุ (เก้าอี้) จากสี่ทิศทาง
(ก) ตำแหน่งที่ระบบรายงานเทียบกับตำแหน่งจริงที่วัดด้วยตลับเมตร ในกรอบพิกัดของหุ่นยนต์
วงกลมเส้นประคือรัศมีเกณฑ์ 50 เซนติเมตร
(ข) ความคลาดเคลื่อนของแต่ละทิศทางเทียบกับเกณฑ์ และค่าเฉลี่ยของทั้งสี่ทิศทาง}
\label{fig:object-error}
\end{figure}

ผลการวัดตามเกณฑ์ข้อที่ 5 สรุปดังตารางที่ \ref{tab:obj5} และรูปที่ \ref{fig:resource}
การวัดอ่านค่าจาก \texttt{tegrastats} ทุก 1 วินาทีตลอดช่วงการวัด
ตัวเลขหลักคืออัตราการใช้งานเฉลี่ยของหน่วยประมวลผลกลางทั้ง 6 แกน
โดยช่วง ``ทั้งระบบ'' เปิดฐานหุ่นยนต์ SLAM ระบบนำทาง การตรวจจับวัตถุ และชั้นผสานข้อมูลพร้อมกัน
ขณะขับหุ่นยนต์ในพื้นที่ทดสอบ ตารางแสดงทุกช่วงที่บันทึกไว้
ยกเว้นช่วงหนึ่งที่ถูกหยุดหลังเริ่มได้ 5 วินาที

\begin{table}[htbp]
\centering
\caption[อัตราการใช้งานทรัพยากรประมวลผลแยกตามสถานการณ์ (บทที่ 5)]{อัตราการใช้งานทรัพยากรประมวลผลแยกตามสถานการณ์ ก่อนและหลังการแก้ไขการตั้งค่าไลบรารี OpenBLAS}
\label{tab:obj5}
\small
\begin{tabularx}{\textwidth}{|X|r|r|r|r|r|}
\hline
\textbf{สถานการณ์} & \textbf{เวลา (วินาที)} & \textbf{CPU (\%)} & \textbf{GPU (\%)} & \textbf{RAM สูงสุด (MB)} & \textbf{กำลังไฟฟ้า (W)} \\
\hline
ขณะว่าง ยังไม่เปิดโหนดใด & 120 & 0.3 & 0 & 1,328 & 4.4 \\
\hline
ฐานหุ่นยนต์ + SLAM & 120 & 0.8 & 0 & 1,557 & 4.5 \\
\hline
ทั้งระบบ ก่อนแก้ไข รอบที่ 1 (YOLO26s ONNX) & 356 & 76.3 & 46 & 3,762 & 8.4 \\
\hline
ทั้งระบบ ก่อนแก้ไข รอบที่ 2 (YOLO26s ONNX) & 596 & \textbf{81.5} & 52 & 3,852 & 9.0 \\
\hline
\textbf{ทั้งระบบ หลังแก้ไข (YOLO26l TensorRT)} & 472 & \textbf{55.8} & 49 & 2,915 & 8.3 \\
\hline
\end{tabularx}
\end{table}

\textbf{ก่อนแก้ไข ช่วงที่ยาวที่สุด (596 วินาที) เฉลี่ยร้อยละ 81.5 ไม่ผ่านเกณฑ์}
ภาระกระจายเท่ากันทุกแกน (ร้อยละ 80.6--82.9) จึงไม่ใช่คอขวดที่แกนใดแกนหนึ่ง
เมื่อตรวจการใช้งานรายเธรด พบว่าโหนดตรวจจับวัตถุใช้ร้อยละ 371 ของหนึ่งแกน
โดยมี 5 เธรดที่ใช้ราวร้อยละ 40 ต่อเนื่องตลอดเวลา
สาเหตุคือกลุ่มเธรดของไลบรารี OpenBLAS ที่ numpy ใช้คำนวณเมทริกซ์ในขั้นตอนติดตามวัตถุ (ByteTrack)
เธรดเหล่านี้วนรอคำสั่งถัดไปแทนการหยุดพัก
และเนื่องจากเฟรมใหม่มาถึงทุก 66 มิลลิวินาที เธรดจึงไม่เคยได้หยุดพัก
ใช้หน่วยประมวลผลรวมประมาณ 2 แกนโดยไม่เกิดงานจริง
ทั้งที่เมทริกซ์ของขั้นตอนติดตามวัตถุมีขนาดเล็ก และไม่ได้ประโยชน์จากการแบ่งงานหลายเธรด
การทดสอบแยกด้วยภาพเดียวกันยืนยันสาเหตุนี้
การจำกัดเธรดของ ONNX Runtime ไม่เปลี่ยนผล (ร้อยละ 263 เป็น 258 ของหนึ่งแกน)
ส่วนการกำหนด \texttt{OPENBLAS\_NUM\_THREADS=1} ลดเหลือร้อยละ 38 ที่อัตราเฟรมเท่าเดิม
การแก้ไขจึงเป็นการกำหนดตัวแปรนี้ในไฟล์เปิดใช้งานของโหนดตรวจจับ
โดยไม่ลดอัตราเฟรม ไม่เปลี่ยนความละเอียดภาพ และไม่ต้องแก้โค้ด

\textbf{หลังแก้ไขเฉลี่ยร้อยละ 55.8 ผ่านเกณฑ์ โดยเหลือส่วนเผื่อ 24 จุด}
ช่วงการวัด 596 วินาทีนี้มีสองช่วงที่ระบบไม่ได้เปิดครบ
คือช่วงวินาทีที่ 129--175 ที่โหนดตรวจจับถูกหยุดและเปิดใหม่
และตั้งแต่วินาทีที่ 521 ที่ปิดการตรวจจับและชั้นผสานข้อมูลก่อนหมดเวลาบันทึก
(แรเงาสีเทาในรูปที่ \ref{fig:resource})
ค่าร้อยละ 55.8 คิดจาก 472 วินาทีที่เหลือซึ่งเปิดครบทั้งระบบ
การตัดสองช่วงนี้ทำให้ค่าเฉลี่ย\emph{สูงขึ้น} จากร้อยละ 48.4 ของทั้งไฟล์
จึงเป็นการรายงานแบบระมัดระวัง
ในช่วงที่เปิดครบ ค่าที่ร้อยละ 95 ของตัวอย่างอยู่ที่ร้อยละ 68.8
มีตัวอย่างที่เกินร้อยละ 80 เพียงร้อยละ 1.5 ของเวลา
และแกนที่ใช้งานมากที่สุดเฉลี่ยร้อยละ 59.8 จึงไม่มีคอขวดที่แกนเดียว
หน่วยความจำสูงสุดลดลง 937 MB เหลือ 2,915 MB จาก 7,620 MB
และอุณหภูมิชิปสูงสุด 53.1 องศาเซลเซียส ต่ำกว่าจุดลดความเร็วสัญญาณนาฬิกา (ประมาณ 97 องศาเซลเซียส) มาก

ระหว่างการวัดก่อนและหลังแก้ไขมีการเปลี่ยนแบบจำลองจาก YOLO26s ONNX เป็น YOLO26l TensorRT ด้วย
(หัวข้อ \ref{sec:eval-det}) แต่เมื่อวัดด้วยการตั้งค่า OpenBLAS เดียวกัน
แบบจำลองทั้งสองใช้หน่วยประมวลผลกลางใกล้เคียงกัน คือร้อยละ 72 และ 67 ของหนึ่งแกน (ตารางที่ \ref{tab:det-cost})
ต่างกันไม่ถึง 1 จุดเมื่อคิดเป็นค่าเฉลี่ย 6 แกน
การลดลงประมาณ 26 จุดจึงมาจากการแก้ไข OpenBLAS เป็นหลัก
ทั้งนี้เธรดที่วนรอใช้ประมาณ 2 แกน แต่การลดลงที่วัดได้เทียบเท่าประมาณ 1.5 แกน
ส่วนต่างนี้ไม่ได้แยกสาเหตุ
ส่วนหนึ่งน่าจะมาจากการวนรอเกิดเฉพาะเมื่อมีวัตถุให้ติดตาม
ภาระจึงขึ้นกับฉากที่กล้องเห็นในแต่ละรอบ
ซึ่งเห็นได้จากรอบก่อนแก้ไขเองที่ลดลงจากประมาณร้อยละ 92 ในช่วงแรกเหลือประมาณร้อยละ 60 ในช่วงท้าย

\begin{figure}[htbp]
\centering
\includegraphics[width=\textwidth]{images/resource_usage.png}
\caption[อัตราการใช้งานทรัพยากรประมวลผล (บทที่ 5)]{อัตราการใช้งานทรัพยากรประมวลผล
(ก) ตลอดช่วงการวัดหลังแก้ไข ขณะที่ SLAM การตรวจจับวัตถุ และชั้นผสานข้อมูลทำงานพร้อมกัน
ช่วงแรเงาคือช่วงที่ระบบไม่ได้เปิดครบและไม่นำมาคิดค่าเฉลี่ย
(ข) หน่วยประมวลผลกลางเฉลี่ยของแต่ละสถานการณ์ในตารางที่ \ref{tab:obj5} เทียบกับเกณฑ์ร้อยละ 80}
\label{fig:resource}
\end{figure}

\subsection{วิธีการและผลการประเมินความแม่นยำของการตรวจจับวัตถุ (เกณฑ์ข้อที่ 2)}
\label{sec:eval-det}

เกณฑ์ข้อที่ 2 ต้องระบุวิธีวัดให้ชัดเจนกว่าข้ออื่น
เนื่องจากแบบจำลองที่ใช้เป็นแบบจำลองที่ผ่านการฝึกมาแล้วบนชุดข้อมูลสาธารณะ
โครงงานนี้ไม่ได้ฝึกแบบจำลองขึ้นเอง จึงไม่มีชุดข้อมูลทดสอบที่แยกไว้ของตนเอง
การรายงานค่าความแม่นยำที่ผู้พัฒนาแบบจำลองประกาศไว้ก็ไม่ได้บอกอะไรเกี่ยวกับระบบนี้
เพราะวัดจากชุดข้อมูลคนละชุดและสภาพแวดล้อมคนละแบบ

\textbf{สิ่งที่โครงงานนี้วัดจึงเป็นความแม่นยำของตัวตรวจจับที่ติดตั้งใช้งานจริง
ในสภาพแวดล้อมเป้าหมาย} ซึ่งเป็นคำถามที่ตรงกับการใช้งานมากกว่า
ภาพในชุดทดสอบนี้ใช้เพื่อการประเมินเท่านั้น ไม่ได้เผยแพร่ และได้รับความยินยอมจากบุคคลที่ปรากฏในภาพ

\subsubsection{ชุดภาพและการกำกับข้อมูล}

บันทึกภาพจากกล้องของหุ่นยนต์ (ภาพดิบก่อนวาดกรอบ) พร้อมผลการตรวจจับของระบบ
เป็นเวลา 290 วินาที ขณะเปิดเฉพาะตัวตรวจจับ (ไม่ได้เปิด SLAM)
ผลการตรวจจับแต่ละชุดมีเวลาประทับตรงกับภาพต้นทาง จึงจับคู่กันได้แบบตรงตัว
เลือกเฉพาะภาพที่ตัวตรวจจับประมวลผลจริง (3,283 จาก 3,962 ภาพ
เพราะตัวตรวจจับทำงานช้ากว่ากล้องเล็กน้อย) แล้วสุ่มอย่างเป็นระบบทุก ๆ 23 ภาพ
ได้ชุดภาพ 143 ภาพ
\textbf{ภาพทั้งหมดถ่ายจากมุมมองเดียวขณะหุ่นยนต์จอดนิ่ง} โดยมีบุคคลเคลื่อนที่ในฉากตลอดการบันทึก
ทั้งเดินผ่านในระยะใกล้ ยืน และนั่งบนเก้าอี้

กรอบวัตถุเบื้องต้นร่างด้วยแบบจำลอง YOLO26x ซึ่งเป็นแบบจำลองคนละตัวกับที่ประเมิน
แล้วตรวจแก้ด้วยมือทุกภาพ ได้แก่ ลบกรอบที่ผิด เพิ่มกรอบที่ขาด และปรับขอบกรอบ
โดยกำกับวัตถุเป้าหมายทุกชิ้นที่มองเห็นในภาพ รวมถึงเก้าอี้ที่อยู่ไกลและถูกบังบางส่วน
ระหว่างกำกับข้อมูลไม่ได้แสดงผลของแบบจำลองที่กำลังประเมิน
เพื่อไม่ให้ข้อมูลอ้างอิงเอนเอียงเข้าหาคำตอบของแบบจำลองนั้น

\textbf{คลาสเป้าหมายกำหนดไว้ล่วงหน้าก่อนเริ่มกำกับข้อมูล} ได้แก่
\texttt{person}, \texttt{chair}, \texttt{backpack} และ \texttt{laptop}
เกณฑ์การเลือกคือ \emph{คลาสที่ระบบสามารถบันทึกลงแผนที่เป็นวัตถุอ้างอิงได้จริง
ในรอบการทดสอบ} ไม่ใช่การเลือกตามความเหมาะสมของผลลัพธ์
คลาสจะนำมาคิดค่าเฉลี่ยก็ต่อเมื่อพบไม่น้อยกว่า 50 ครั้งในชุดทดสอบ
ซึ่งทั้ง 4 คลาสผ่านเงื่อนไขนี้ (person 292, chair 441, backpack 143 และ laptop 116 ชิ้น)

\subsubsection{การวัด}

การจับคู่ผลการตรวจจับกับข้อมูลที่กำกับด้วยมือใช้เกณฑ์การซ้อนทับของกรอบ (IoU)
ไม่น้อยกว่า 0.5 ต้องเป็นคลาสเดียวกัน และพิจารณาเฉพาะผลที่มีค่าความเชื่อมั่นไม่ต่ำกว่าค่าที่ระบบใช้งาน
ค่าที่รายงานประกอบด้วยความแม่นยำ (Precision) ความครบถ้วน (Recall)
และค่าเฉลี่ยฮาร์มอนิกของทั้งสอง (F1) แยกตามคลาส พร้อมค่าเฉลี่ยข้ามคลาส

\textbf{เกณฑ์ร้อยละ 80 ตรวจสอบกับค่า F1 เฉลี่ย ไม่ใช่กับความแม่นยำเพียงอย่างเดียว}
เหตุผลคือความแม่นยำสูงขึ้นได้ด้วยการยกระดับค่าความเชื่อมั่นขั้นต่ำ
และความครบถ้วนสูงขึ้นได้ด้วยการลดค่าดังกล่าว
การเลือกรายงานค่าใดค่าหนึ่งเพียงลำพังจึงทำให้ตัวเลขดูดีได้โดยที่ระบบไม่ได้ดีขึ้น
ขณะที่ค่า F1 ไม่สามารถทำเช่นนั้นได้ ทั้งสามค่าจึงถูกรายงานคู่กันเสมอ

\subsubsection{ผลของตัวตรวจจับเดิม}

ตัวตรวจจับที่ติดตั้งใช้งานเดิม (YOLO26s รูปแบบ ONNX ความละเอียด 16 บิต ค่าความเชื่อมั่น 0.5)
ได้ผลดังตารางที่ \ref{tab:det-baseline}
ค่า F1 เฉลี่ยร้อยละ 64.4 \textbf{ไม่ผ่านเกณฑ์} โดยปัญหาอยู่ที่ความครบถ้วน ไม่ใช่ความแม่นยำ
เมื่อแยกตามขนาดของวัตถุในภาพ วัตถุขนาด 32--96 พิกเซลซึ่งเป็นเก้าอี้ใต้โต๊ะและบุคคลที่อยู่ไกล
ตรวจพบเพียง 4 จาก 231 ชิ้น (ร้อยละ 1.7) ขณะที่วัตถุขนาดใหญ่กว่า 96 พิกเซลตรวจพบร้อยละ 51.1

\begin{table}[htbp]
\centering
\caption[ผลการประเมินตัวตรวจจับเดิม YOLO26s (บทที่ 5)]{ผลการประเมินตัวตรวจจับเดิม YOLO26s ONNX 16 บิต ค่าความเชื่อมั่น 0.5 (143 ภาพ)}
\label{tab:det-baseline}
\small
\begin{tabular}{|l|r|r|r|r|r|r|}
\hline
\textbf{คลาส} & \textbf{จำนวนจริง} & \textbf{TP} & \textbf{FP} & \textbf{Precision} & \textbf{Recall} & \textbf{F1} \\
\hline
person & 292 & 134 & 9 & 93.7 & 45.9 & 61.6 \\
chair & 441 & 69 & 5 & 93.2 & 15.7 & 26.8 \\
backpack & 143 & 90 & 1 & 98.9 & 62.9 & 76.9 \\
laptop & 116 & 99 & 0 & 100.0 & 85.3 & 92.1 \\
\hline
\textbf{เฉลี่ย} & & & & \textbf{96.5} & \textbf{52.5} & \textbf{64.4} \\
\hline
\end{tabular}
\end{table}

\subsubsection{การปรับปรุงแบบจำลอง}

การปรับปรุงทำทีละขั้นบนชุดภาพเดียวกัน และบันทึกผลทุกขั้นดังตารางที่ \ref{tab:det-improve}
พบข้อสังเกตสำคัญสามประการ
ประการแรก แบบจำลองขนาดใหญ่ขึ้นช่วยเรื่องวัตถุระยะไกล โดย YOLO26l เพิ่มค่า F1 ของ chair
จากร้อยละ 22 เป็น 48 และของ person จากร้อยละ 64 เป็น 72
ประการที่สอง \textbf{การส่งออกเป็น TensorRT ความละเอียด 16 บิตไม่ทำให้ความแม่นยำลดลง}
แต่การกำหนดขนาดภาพนำเข้าเป็นสี่เหลี่ยมจัตุรัส 640$\times$640
ทำให้ภาพขนาด 640$\times$480 ของกล้องถูกเติมขอบว่าง 160 แถว และค่า F1 ลดลงประมาณ 2 จุด
การส่งออกที่ขนาด 480$\times$640 ตรงกับกล้องให้ผลเท่ากับแบบจำลองต้นฉบับ
ประการที่สาม การลดค่าความเชื่อมั่นขั้นต่ำเพิ่มความครบถ้วนโดยแลกกับความแม่นยำเพียงเล็กน้อย
ทั้งนี้ไม่ได้ใช้ YOLO26x เป็นตัวเลือก เนื่องจากเป็นแบบจำลองที่ใช้ร่างข้อมูลอ้างอิง

\begin{table}[htbp]
\centering
\caption[ขั้นตอนการปรับปรุงแบบจำลองตรวจจับวัตถุ (บทที่ 5)]{ขั้นตอนการปรับปรุงแบบจำลองตรวจจับวัตถุ ประเมินบนชุดภาพเดียวกัน 143 ภาพ}
\label{tab:det-improve}
\small
\begin{tabularx}{\textwidth}{|X|c|r|r|r|}
\hline
\textbf{แบบจำลองและรูปแบบ} & \textbf{ค่าความเชื่อมั่น} & \textbf{Precision} & \textbf{Recall} & \textbf{F1} \\
\hline
YOLO26s ONNX 640$\times$640 (ตัวตรวจจับเดิม) & 0.5 & 96.5 & 52.5 & 64.4 \\
YOLO26s ต้นฉบับ & 0.5 & 93.0 & 56.3 & 66.3 \\
YOLO26m ต้นฉบับ & 0.5 & 96.8 & 54.7 & 67.4 \\
YOLO26m TensorRT 16 บิต 640$\times$640 & 0.5 & 95.3 & 52.7 & 65.6 \\
YOLO26l ต้นฉบับ & 0.5 & 97.1 & 71.0 & 79.0 \\
YOLO26l TensorRT 16 บิต 480$\times$640 & 0.5 & 97.1 & 71.1 & 79.2 \\
YOLO26l TensorRT 16 บิต 480$\times$640 & 0.45 & 97.0 & 73.6 & 81.0 \\
\textbf{YOLO26l TensorRT 16 บิต 480$\times$640} & \textbf{0.4} & \textbf{96.2} & \textbf{75.1} & \textbf{82.2} \\
\hline
\end{tabularx}
\end{table}

ผลของตัวตรวจจับหลังปรับปรุงแยกตามคลาสแสดงดังตารางที่ \ref{tab:det-final}
และรูปที่ \ref{fig:det-before-after}
ค่า F1 เฉลี่ยร้อยละ 82.2 \textbf{ผ่านเกณฑ์ร้อยละ 80}
โดยความแม่นยำแทบไม่เปลี่ยน (ร้อยละ 96.5 เป็น 96.2) ขณะที่ความครบถ้วนเพิ่มจากร้อยละ 52.5 เป็น 75.1
เมื่อทดสอบผ่านโหนดตรวจจับจริงพร้อมตัวติดตามวัตถุ โดยป้อนภาพชุดเดิมย้อนหลัง
ที่ค่าความเชื่อมั่น 0.5 ได้ค่า F1 ร้อยละ 78.0 เทียบกับ 63.7 ของตัวตรวจจับเดิมบนภาพชุดเดียวกัน
ซึ่งสอดคล้องกับการประเมินแบบออฟไลน์

\begin{table}[htbp]
\centering
\caption[ผลการประเมินตัวตรวจจับหลังปรับปรุง YOLO26l (บทที่ 5)]{ผลการประเมินตัวตรวจจับหลังปรับปรุง YOLO26l TensorRT 16 บิต ค่าความเชื่อมั่น 0.4 (143 ภาพ)}
\label{tab:det-final}
\small
\begin{tabular}{|l|r|r|r|r|r|r|}
\hline
\textbf{คลาส} & \textbf{จำนวนจริง} & \textbf{TP} & \textbf{FP} & \textbf{Precision} & \textbf{Recall} & \textbf{F1} \\
\hline
person & 292 & 190 & 19 & 90.9 & 65.1 & 75.8 \\
chair & 441 & 172 & 8 & 95.6 & 39.0 & 55.4 \\
backpack & 143 & 138 & 0 & 100.0 & 96.5 & 98.2 \\
laptop & 116 & 116 & 2 & 98.3 & 100.0 & 99.2 \\
\hline
\textbf{เฉลี่ย} & & & & \textbf{96.2} & \textbf{75.1} & \textbf{82.2} \\
\hline
\end{tabular}
\end{table}

\begin{figure}[htbp]
\centering
\includegraphics[width=0.85\textwidth]{images/accuracy_before_after.png}
\caption[ความแม่นยำของการตรวจจับก่อนและหลังปรับปรุง (บทที่ 5)]{ค่าเฉลี่ยข้ามคลาสของความแม่นยำ ความครบถ้วน และ F1 ของตัวตรวจจับเดิม (YOLO26s ONNX ค่าความเชื่อมั่น 0.5) และหลังปรับปรุง (YOLO26l TensorRT ค่าความเชื่อมั่น 0.4) เทียบกับเกณฑ์ร้อยละ 80}
\label{fig:det-before-after}
\end{figure}

การใช้ TensorRT ทำให้แบบจำลองที่ใหญ่ขึ้นใช้ทรัพยากรไม่มากกว่าเดิม ดังตารางที่ \ref{tab:det-cost}
YOLO26l บน TensorRT ใช้เวลาอนุมานลดลงร้อยละ 51 เมื่อเทียบกับรูปแบบต้นฉบับ
และใช้หน่วยประมวลผลกราฟิกและหน่วยความจำน้อยกว่า YOLO26s ONNX ที่ใช้เดิม
(วัดบนอุปกรณ์จริงที่โหมดพลังงาน 15 วัตต์ ขณะไม่มีงานอื่นใช้หน่วยประมวลผลกราฟิก)

\begin{table}[htbp]
\centering
\caption[ทรัพยากรที่ใช้ของแบบจำลองแต่ละรูปแบบ (บทที่ 5)]{ทรัพยากรที่ใช้ของแบบจำลองแต่ละรูปแบบบน Jetson Orin NX 8 GB}
\label{tab:det-cost}
\small
\begin{tabularx}{\textwidth}{|X|r|r|r|r|}
\hline
\textbf{แบบจำลองและรูปแบบ} & \textbf{เวลาอนุมาน (ms)} & \textbf{GPU (\%)} & \textbf{CPU (\% ของ 1 แกน)} & \textbf{หน่วยความจำ (MB)} \\
\hline
YOLO26s ONNX 16 บิต (เดิม) & 36.7 & 50.7 & 72 & 1,091 \\
YOLO26l ต้นฉบับ & 74.4 & 68.9 & 78 & 946 \\
\textbf{YOLO26l TensorRT 16 บิต} & \textbf{36.4} & \textbf{44.8} & \textbf{67} & \textbf{766} \\
\hline
\end{tabularx}
\end{table}

\subsubsection{ข้อจำกัดของผลการประเมิน}

\begin{itemize}[leftmargin=1.25cm, itemsep=6pt, topsep=6pt]
    \item \textbf{มุมมองเดียว} ภาพทั้งหมดถ่ายขณะหุ่นยนต์จอดนิ่ง
    คลาส laptop และ backpack มีวัตถุอย่างละหนึ่งชิ้นในตำแหน่งคงที่
    ผลของสองคลาสนี้ (ร้อยละ 98--99) จึงสะท้อนวัตถุชิ้นเดียว
    และไม่ได้ทดสอบความหลากหลายของระยะและมุมมอง
    ขณะที่คลาส chair ซึ่งมีวัตถุหลายชิ้นหลายระยะยังได้ค่า F1 เพียงร้อยละ 55.4
    \item \textbf{ข้อมูลอ้างอิงเริ่มจากร่างของ YOLO26x} แม้จะตรวจแก้ด้วยมือทุกภาพ
    แบบจำลองในตระกูลเดียวกันที่ขนาดใกล้เคียงอาจได้ประโยชน์จากข้อผิดพลาดที่ตรวจไม่พบในร่าง
    \item \textbf{อัตราเฟรมของ YOLO26l วัดขณะหุ่นยนต์จอดนิ่ง}
    ได้ 15.16 เฟรมต่อวินาทีตลอด 120 วินาที ขณะ SLAM และชั้นผสานข้อมูลทำงาน
    (ไม่ได้เปิดระบบนำทางและไม่ได้บันทึกข้อมูล) ส่วนค่า 13.05 เฟรมต่อวินาทีของ YOLO26s
    วัดขณะเปิดทั้งระบบพร้อมการบันทึกข้อมูลระหว่างขับหุ่นยนต์ เงื่อนไขจึงไม่เหมือนกัน
    ชั้นผสานข้อมูลยังคัดผลการตรวจจับที่ค่าความเชื่อมั่นต่ำกว่า 0.5 ออกเช่นเดิม
    การลดค่าความเชื่อมั่นของตัวตรวจจับจึงไม่เพิ่มวัตถุที่ไม่น่าเชื่อถือลงบนแผนที่
\end{itemize}

\section{ปัญหาและอุปสรรคในการทำโครงงาน}

\subsection{ปัญหาด้านฮาร์ดแวร์และเฟิร์มแวร์}

\begin{itemize}[leftmargin=1.25cm, itemsep=6pt, topsep=6pt]
    \item \textbf{การใช้ขาสัญญาณของ ESP32 ที่เป็นขากำหนดสถานะขณะบูต}
    ขาสัญญาณควบคุมทิศทางมอเตอร์ข้างขวาถูกต่อเข้ากับขา GPIO12
    ซึ่งเป็นขาที่ชิปอ่านค่าขณะบูตเพื่อกำหนดแรงดันของหน่วยความจำแฟลช
    หากขานี้ถูกดึงขึ้นสู่แรงดันสูงในจังหวะรีเซ็ต บอร์ดอาจไม่บูตหรือทำงานผิดปกติ
    เนื่องจากหุ่นยนต์ถูกเดินสายไว้เช่นนี้ตั้งแต่ต้นและการย้ายขาหมายถึงการรื้อสายใหม่
    ระบบจึงเลือกวิธีตรวจสอบความน่าเชื่อถือของการบูตซ้ำหลายรอบแทนการย้ายขา
    และยืนยันว่าวงจรขับมอเตอร์ไม่ได้ดึงขานี้ขึ้นในจังหวะรีเซ็ต
    \textbf{ข้อจำกัดที่ต้องระบุไว้คือขานี้ยังคงเป็นขากำหนดสถานะขณะบูต}
    หากเปลี่ยนโมดูลขับมอเตอร์หรือเพิ่มตัวต้านทานดึงขึ้นที่ขานี้ในอนาคต
    ปัญหาการบูตจะกลับมาอีก

    \item \textbf{ขาเอนโคเดอร์ของล้อซ้ายเป็นขาที่รับสัญญาณเข้าได้อย่างเดียว}
    ขา GPIO34 และ GPIO35 ไม่มีตัวต้านทานดึงขึ้นภายใน
    เอนโคเดอร์จึงต้องขับสัญญาณออกมาเองอย่างชัดเจน
    และทำให้ทุกครั้งที่บูตจะมีข้อความแจ้งเตือนเรื่องขาที่ไม่มีตัวต้านทานดึงขึ้นปรากฏสี่บรรทัดเสมอ
    ซึ่งเป็นพฤติกรรมปกติของระบบ ไม่ใช่ความผิดพลาด

    \item \textbf{จัมเปอร์บนโมดูลขับมอเตอร์} โมดูลขับมอเตอร์ที่ใช้มีจัมเปอร์
    สำหรับตรึงขาเปิดใช้งานช่องสัญญาณไว้ที่แรงดันสูงตลอดเวลา
    เมื่อไม่ถอดจัมเปอร์ออก สัญญาณควบคุมความเร็วจะไม่มีผล
    ทำให้มอเตอร์หมุนด้วยกำลังเต็มที่ตลอดเวลา

    \item \textbf{การสร้างสัญญาณควบคุมความเร็วบน ESP32} ฟังก์ชันสร้างสัญญาณ
    ที่ใช้กันทั่วไปบนไมโครคอนโทรลเลอร์ตระกูลอื่นไม่สามารถใช้ได้โดยตรงบน ESP32
    ต้องเปลี่ยนไปใช้ชุดคำสั่งของตัวควบคุมสัญญาณที่ชิปมีมาให้โดยเฉพาะ

    \item \textbf{ข้อจำกัดจำนวนพอร์ตยูเอสบี} หน่วยประมวลผลหลักมีพอร์ตยูเอสบีจำกัด
    ซึ่งถูกใช้งานครบโดยไลดาร์ กล้อง และไมโครคอนโทรลเลอร์
    การเพิ่มอุปกรณ์ใด ๆ จึงต้องใช้ฮับที่มีแหล่งจ่ายไฟภายนอก
    เพราะการแบ่งกระแสจากพอร์ตเดิมทำให้มอเตอร์ของไลดาร์หมุนไม่สม่ำเสมอ

    \item \textbf{รุ่นของไลดาร์ถูกระบุผิดและไม่มีทางตรวจสอบจากซอฟต์แวร์}
    เอกสารและไฟล์กำหนดค่าของโครงงานระบุรุ่นของไลดาร์ผิดมาโดยตลอดจนถึงช่วงท้ายของการพัฒนา
    สาเหตุคือชุดพัฒนาของผู้ผลิตไม่มีรายการรุ่นนี้อยู่ในตัวเลย
    และไลดาร์แบบช่องสัญญาณเดียวไม่ตอบคำสั่งสอบถามข้อมูลอุปกรณ์
    ทำให้เกิดข้อความแจ้งเตือนตอนเริ่มต้นที่ถูกเข้าใจว่าเป็นข้อความปกติ
    ผลกระทบที่แท้จริงไม่ได้อยู่ที่ชื่อรุ่น แต่อยู่ที่ระยะวัดสูงสุดที่ประกาศไว้ในไฟล์กำหนดค่า
    ซึ่งเดิมตั้งไว้ที่ 12 เมตรสำหรับเซนเซอร์ที่มีระยะวัดจริง 8 เมตร
    เท่ากับบอกให้ทั้ง slam\_toolbox และระบบนำทางเชื่อถือระยะทางที่ฮาร์ดแวร์ให้ไม่ได้
    และนับทิศทางที่ไม่มีวัตถุอยู่ในระยะเป็นลำแสงที่ไม่ได้ค่า
    การแก้ไขคือการอ่านฉลากบนตัวอุปกรณ์แล้วแก้ค่าระยะวัดในไฟล์กำหนดค่าทั้งสี่ไฟล์ให้ตรงกัน
    \textbf{บทเรียนคือคุณลักษณะเชิงกายภาพของเซนเซอร์ต้องยืนยันจากตัวอุปกรณ์จริง
    ไม่ใช่จากเอกสารที่ส่งต่อกันมา}

    \item \textbf{ขั้วต่อแบตเตอรี่ที่หลวมทำให้วินิจฉัยสาเหตุผิด}
    การทดสอบพฤติกรรมกู้คืนสถานการณ์ด้วยการหมุนตัวล้มเหลวสองครั้งติดกัน
    โดยหุ่นยนต์ไม่ขยับเลย ข้อมูลทุกอย่างที่ตรวจสอบได้ในขณะนั้นชี้ไปที่แรงเสียดทานสถิต
    คือคำสั่งความเร็วถูกส่งถึงคอนโทรลเลอร์ครบถ้วน แต่ค่าเอนโคเดอร์ไม่เปลี่ยนแปลงเลย
    สาเหตุที่แท้จริงคือขั้วต่อแบตเตอรี่ของภาคขับมอเตอร์หลวม
    และเนื่องจากไมโครคอนโทรลเลอร์รับไฟจากพอร์ตยูเอสบี
    การสื่อสารอนุกรมและการอ่านค่าเอนโคเดอร์จึงทำงานเป็นปกติทุกประการขณะที่มอเตอร์ไม่มีไฟเลี้ยง
    เมื่อเสียบขั้วต่อใหม่ การทดสอบผ่านทันทีด้วยค่าพารามิเตอร์ชุดเดิม
    \textbf{บทเรียนคือต้องตรวจแหล่งจ่ายไฟของภาคขับมอเตอร์ก่อนสงสัยเรื่องแรงเสียดทานเสมอ}
\end{itemize}

\subsection{ปัญหาด้านซอฟต์แวร์และสภาพแวดล้อมการพัฒนา}

\begin{itemize}[leftmargin=1.25cm, itemsep=6pt, topsep=6pt]
    \item \textbf{ไลบรารีเรียนรู้เชิงลึกรุ่นที่ไม่รองรับหน่วยประมวลผลกราฟิก}
    การติดตั้งไลบรารีจากแหล่งมาตรฐานจะได้รุ่นที่รองรับเฉพาะหน่วยประมวลผลกลาง
    เมื่อนำมาใช้บนอุปกรณ์เป้าหมาย แบบจำลองจึงไม่ได้ใช้หน่วยประมวลผลกราฟิกเลย
    และทำงานช้าลงหลายเท่า \textbf{โดยไม่มีข้อความแจ้งเตือนใด ๆ}
    อาการที่สังเกตได้มีเพียงอัตราเฟรมที่ตกลง
    การแก้ไขคือติดตั้งไลบรารีรุ่นที่จัดทำมาสำหรับอุปกรณ์ตระกูลนี้โดยเฉพาะ
    ระบุที่อยู่ของไฟล์ติดตั้งโดยตรงเพื่อไม่ให้ถูกแทนที่ด้วยรุ่นมาตรฐาน
    และตรวจสอบทุกครั้งว่าระบบเลือกใช้หน่วยประมวลผลกราฟิกจริง

    \item \textbf{ลำดับการติดตั้งไลบรารีมีผลต่อผลลัพธ์} การติดตั้งไลบรารีตรวจจับวัตถุก่อน
    จะทำให้ตัวจัดการแพ็กเกจดึงไลบรารีเรียนรู้เชิงลึกรุ่นมาตรฐานเข้ามาแทนที่รุ่นที่ถูกต้อง
    จึงต้องติดตั้งไลบรารีเรียนรู้เชิงลึกให้เสร็จก่อนเสมอ
    แล้วจึงติดตั้งไลบรารีตรวจจับวัตถุโดยระบุไม่ให้ติดตั้งไลบรารีที่ต้องพึ่งพาตามมาโดยอัตโนมัติ

    \item \textbf{แบบจำลองที่ผูกกับรุ่นของไลบรารีอนุมาน} เครื่องยนต์อนุมานของ TensorRT
    ที่สร้างไว้ก่อนหน้าไม่สามารถโหลดได้อีกเมื่อเปลี่ยนรุ่นชุดซอฟต์แวร์ของอุปกรณ์
    เป็นเหตุให้ระบบรุ่นแรกเปลี่ยนมาใช้รูปแบบ ONNX ซึ่งไม่ผูกกับรุ่นของไลบรารี
    ตามที่อธิบายไว้ในบทที่ 2
    เมื่อเปลี่ยนเป็น YOLO26l ระบบกลับมาใช้ TensorRT โดยแก้ปัญหานี้ด้วยการเก็บไฟล์น้ำหนักต้นฉบับไว้เป็นแหล่งที่มา
    และให้คำสั่งเปิดระบบสร้างเครื่องยนต์ใหม่บนอุปกรณ์เองเมื่อไม่มีไฟล์
    การเปลี่ยนรุ่นชุดซอฟต์แวร์จึงเสียเพียงเวลาสร้างใหม่ประมาณ 13 นาที ไม่ได้สูญเสียแบบจำลอง

    \item \textbf{ความเร็วในการขับหุ่นยนต์ด้วยมือทำให้แผนที่เสียหายถาวร}
    การสแกนหนึ่งรอบของไลดาร์กินเวลา 86 มิลลิวินาที
    และประทับเวลาของข้อมูลมีอายุประมาณหนึ่งรอบการสแกนก่อนที่ผู้ใช้ข้อมูลจะได้รับ
    แต่ระบบถือว่าข้อมูลทั้งรอบเกิดขึ้น ณ ตำแหน่งเดียว
    ที่ความเร็วเริ่มต้นของโปรแกรมขับด้วยแป้นพิมพ์คือ 0.5 เมตรต่อวินาที
    หุ่นยนต์เคลื่อนที่ไปแล้วประมาณ 8.7 เซนติเมตรภายในช่วงเวลาดังกล่าว
    ทำให้ข้อมูลการสแกนบิดเบี้ยวและแผนที่เบลอ
    ซึ่งเป็นความเสียหายที่แก้ย้อนหลังไม่ได้เมื่อถูกเขียนลงกราฟไปแล้ว
    \textbf{ประเด็นที่ทำให้ปัญหานี้ตรวจพบได้ยากคือโปรแกรมขับด้วยแป้นพิมพ์ส่งคำสั่งตรงไปยังคอนโทรลเลอร์
    โดยไม่ผ่านขีดจำกัดความเร็วของระบบนำทาง} และไม่มีการแจ้งเตือนใด ๆ
    การแก้ไขคือลดความเร็วในการขับด้วยมือให้อยู่ในระดับเดียวกับที่ระบบนำทางใช้
    ซึ่งยืนยันแล้วว่าทำให้แผนที่ของพื้นที่เดียวกันคมชัดขึ้น
    ในช่วงท้ายของโครงงาน ขีดจำกัดของการขับด้วยมือถูกปรับขึ้นเป็น 0.30 เมตรต่อวินาที
    เพื่อใช้ย้ายตำแหน่งหุ่นยนต์ระหว่างการซ้อมสาธิต
    \textbf{แต่การขับเพื่อสร้างแผนที่ทั้งหมดใช้ความเร็ว 0.10 เมตรต่อวินาที}
    ซึ่งข้อมูลการสแกนเหลื่อมกันเพียงประมาณ 0.9 เซนติเมตรต่อรอบ
    ข้อควรระวังคือระบบไม่ได้บังคับความเร็วนี้โดยอัตโนมัติ
    ผู้ใช้ต้องกำหนดความเร็วเองทุกครั้งที่สร้างแผนที่


    \item \textbf{ชื่ออุปกรณ์ต่อพ่วงสลับกันเมื่อเปิดเครื่องใหม่}
    แก้ไขโดยกำหนดกฎ udev เพื่อสร้างชื่ออุปกรณ์ถาวรตามที่อธิบายไว้ในบทที่ 4

    \item \textbf{เป้าหมายการนำทางถูกประกาศในระบบพิกัดที่ผิด}
    การทดสอบระบบนำทางครั้งแรกล้มเหลวทั้งสี่ครั้งที่กำหนดเป้าหมาย
    สาเหตุคือเครื่องมือแสดงผลประทับระบบพิกัดของเป้าหมายตามเฟรมอ้างอิงที่ตั้งไว้ในตัวเครื่องมือเอง
    ซึ่งถูกเปลี่ยนจาก \texttt{map} เป็น \texttt{odom} ระหว่างการทดสอบก่อนหน้า
    อาการที่สังเกตได้คือหุ่นยนต์เริ่มเคลื่อนที่ตามปกติในช่วงแรกแล้วหยุดไปเฉย ๆ
    เพราะต้นไม้พฤติกรรมเก็บเป้าหมายเดิมพร้อมประทับเวลาเดิมไว้และวางแผนเส้นทางซ้ำทุกวินาที
    การวางแผนจึงสำเร็จอยู่ประมาณสิบวินาทีแรกแล้วล้มเหลวตลอดไป
    \textbf{ประเด็นที่ทำให้วินิจฉัยยากคือข้อความแจ้งข้อผิดพลาดชี้ไปที่การแปลงพิกัด}
    ซึ่งชวนให้ไปปรับค่าความลึกของบัฟเฟอร์การแปลงพิกัด
    ทั้งที่การวัดพบว่าบัฟเฟอร์เก็บข้อมูลได้ครบ 10 วินาทีตามปกติ
    และค่าที่ระบบรายงานคือระยะเวลาที่คำขอเก่าเกินบัฟเฟอร์ ไม่ใช่ความลึกของบัฟเฟอร์
    การแก้ไขคือการตั้งเฟรมอ้างอิงของเครื่องมือกลับเป็น \texttt{map} โดยไม่ต้องแก้โค้ดใด ๆ

    \item \textbf{พฤติกรรมกู้คืนสถานการณ์หมดเวลาทุกครั้งด้วยเหตุผลเชิงเลขคณิต}
    ต้นไม้พฤติกรรมมาตรฐานของ Nav2 กำหนดให้การหมุนตัวเพื่อกู้คืนสถานการณ์หมุน 1.57 เรเดียน
    โดยมีเวลาที่ยอมให้ใช้ตามค่าเริ่มต้นเพียง 10 วินาที
    ขณะที่ความเร็วเชิงมุมสูงสุดของพฤติกรรมกู้คืนที่ระบบตั้งไว้คือ 0.1 เรเดียนต่อวินาที
    ซึ่งต้องใช้เวลาอย่างน้อย 15.7 วินาที การหมุนจึงถูกยกเลิกที่ 10.000 วินาทีพอดีทุกครั้ง
    และเนื่องจากค่าเวลานี้เป็นพารามิเตอร์ของต้นไม้พฤติกรรม ไม่ใช่พารามิเตอร์ของ ROS
    จึงไม่สามารถแก้จากไฟล์กำหนดค่าได้
    การแก้ไขคือการนำไฟล์ต้นไม้พฤติกรรมมาไว้ในแพ็กเกจของโครงงานเองและกำหนดเวลาเป็น 25 วินาที
    หลังแก้ไขแล้ววัดได้ว่าการหมุนสำเร็จในเวลา 16.40 วินาที

    \item \textbf{ข้อบกพร่องที่พบตอนรวมระบบ ซึ่งไม่มีข้อใดอยู่ในส่วนผสานข้อมูลเลย}
    ในการทดสอบรวมระบบครั้งแรก พบข้อบกพร่องสี่จุดที่ทุกจุดแสดงอาการเหมือนหุ่นยนต์ทำงานผิด
    จุดแรกคือภาวะแข่งกันตอนเริ่มโหนดผสานข้อมูล
    ซึ่งประกาศตัวแปรเก็บสถิติหลังจากสมัครรับข้อมูลไปแล้ว
    ทำให้การเริ่มโหนดในขณะที่ข้อมูลการตรวจจับไหลอยู่แล้วเกิดข้อผิดพลาดในฟังก์ชันรับข้อมูล
    และหยุดตัวจัดการงานทั้งหมดอย่างเงียบ ๆ โดยที่กระบวนการยังทำงานอยู่และยังปรากฏในรายการโหนด
    ข้อบกพร่องนี้ขึ้นกับลำดับการเปิดระบบ และลำดับที่ใช้งานจริงคือลำดับที่ทำให้เกิดปัญหาพอดี
    จุดที่สองคือเว็บแอปพลิเคชันแสดงข้อความว่าไม่มีสัญญาณภาพตลอดเวลา
    ทั้งที่ภาพส่งมาถึงที่อัตรา 15.3 เฟรมต่อวินาที
    จุดที่สามคือหน้าเว็บเรียกแผนที่เพียงครั้งเดียวตอนเปิดหน้า แล้วไม่เรียกซ้ำอีกเลย
    ขณะที่ระบบ SLAM ยังขยายแผนที่ต่อเนื่อง
    และจุดที่สี่คือสีของชื่อคลาสวัตถุมีอัตราส่วนความต่างของความสว่างเทียบกับพื้นหลังเพียง 1.08 ต่อ 1
    ทำให้ชื่อวัตถุมองไม่เห็นในบริเวณที่หุ่นยนต์สำรวจไปแล้วพอดี
    \textbf{บทเรียนคือเมื่อระบบแสดงผลผิด สาเหตุอยู่ที่ชั้นแสดงผลได้บ่อยพอ ๆ กับที่อยู่ในหุ่นยนต์}
\end{itemize}

\subsection{ปัญหาด้านการบริหารจัดการซอร์สโค้ดและข้อมูล}

ปัญหาที่ส่งผลกระทบต่อโครงงานมากที่สุดไม่ใช่ปัญหาเชิงเทคนิคของระบบ
แต่คือ \textbf{การที่หน่วยประมวลผลหลักเสียหายและทำให้ซอร์สโค้ดของระบบสูญหายทั้งหมด}
ในขณะนั้นแพ็กเกจหลักของหุ่นยนต์ถูกพัฒนาไว้เฉพาะบนเครื่องเป้าหมาย
โดยยังไม่ได้สร้างที่เก็บซอร์สโค้ดระยะไกลและไม่มีการสำรองข้อมูล
สิ่งที่สูญหายไปพร้อมกันไม่ได้มีเพียงโค้ด แต่รวมถึงค่าที่ได้จากการสอบเทียบทั้งหมด
ซึ่งเป็นข้อมูลที่ต้องใช้เวลาวัดใหม่บนฮาร์ดแวร์จริงและไม่สามารถเขียนขึ้นใหม่จากความจำได้

แนวทางที่ใช้แก้ไขและป้องกันไม่ให้เกิดซ้ำมีสามข้อ
ข้อแรกคือ \textbf{ทุกที่เก็บซอร์สโค้ดต้องมีที่เก็บระยะไกลและมีการคอมมิตครั้งแรก
ก่อนที่จะเพิ่มความสามารถใดลงไปเป็นอย่างที่สอง}
ข้อที่สองคือค่าที่ได้จากการวัดหรือสอบเทียบทุกค่าต้องถูกบันทึกลงเอกสารบันทึกการสอบเทียบ
พร้อมวันที่และวิธีการวัดทันทีที่วัดได้ ไม่ใช่เมื่อจบวัน
เพื่อให้สามารถตรวจสอบย้อนกลับหรือวัดซ้ำได้ในภายหลัง
และข้อที่สามคือค่ากำหนดที่มีผลต่อพฤติกรรมของระบบต้องอยู่ในไฟล์กำหนดค่าที่อยู่ในที่เก็บซอร์สโค้ด
ไม่ใช่ตั้งไว้เฉพาะบนเครื่องเป้าหมาย
ประสบการณ์นี้เป็นเหตุผลที่โครงสร้างของระบบในบทที่ 4 แยกค่ากำหนดทั้งหมด
ออกจากโค้ดอย่างชัดเจน และมีการอ้างอิงแหล่งข้อมูลเพียงแหล่งเดียวสำหรับแต่ละค่า
เช่น การอ่านตำแหน่งติดตั้งของกล้องจากไฟล์ URDF แทนการกำหนดซ้ำในไฟล์พารามิเตอร์

\subsection{ข้อจำกัดเชิงขอบเขตที่พบระหว่างการพัฒนา}

ข้อจำกัดที่พบคือแม้การประมวลผลจะเกิดขึ้นภายในตัวอุปกรณ์ทั้งหมด
แต่การแสดงผลบนเว็บแอปพลิเคชันยังคงต้องอาศัยเครือข่ายไร้สายภายในพื้นที่
ซึ่งเป็นคนละเรื่องกับการพึ่งพาการประมวลผลบนคลาวด์

\section{ข้อจำกัดของระบบที่ต้องระบุไว้}

นอกเหนือจากปัญหาที่พบระหว่างการพัฒนา
ระบบยังมีข้อจำกัดเชิงออกแบบที่เกิดจากการตัดสินใจตัดขอบเขตบางส่วนออกอย่างตั้งใจ
เพื่อให้ระบบหลักทำงานได้ครบภายในเวลาที่มี ข้อจำกัดเหล่านี้มีดังนี้

\begin{enumerate}[label=\arabic*., leftmargin=1.25cm, itemsep=6pt, topsep=6pt]
    \item \textbf{การสังเกตวัตถุจากระยะไกลถูกให้น้ำหนักมากเกินจริง}
    ระบบรวมตำแหน่งวัตถุด้วยค่าเฉลี่ยเคลื่อนที่ที่ค่าคงที่ตายตัว
    ซึ่งไม่ได้คำนึงว่าความไม่แน่นอนของตำแหน่งเพิ่มขึ้นตามระยะทาง
    การสังเกตจากระยะ 4 เมตรจึงมีน้ำหนักเท่ากับการสังเกตจากระยะ 1 เมตร

    \item \textbf{วัตถุที่บันทึกแล้วไม่ถูกลบออกแม้พิสูจน์ได้ว่าไม่มีอยู่จริง}
    ระบบไม่ได้ตรวจสอบย้อนกลับว่าบริเวณที่เคยบันทึกวัตถุไว้ยังมีวัตถุอยู่หรือไม่
    ผลบวกลวงที่เกิดขึ้นครั้งเดียวและผ่านเกณฑ์การยืนยันจะคงอยู่ตลอดรอบการสำรวจนั้น

    \item \textbf{การวนกลับมาที่เดิมปรับแผนที่แต่ไม่ปรับตำแหน่งวัตถุ}
    ระบบเก็บตำแหน่งวัตถุเป็นพิกัดสัมบูรณ์บนเฟรมแผนที่
    เมื่อ SLAM ตรวจพบการวนกลับมาที่เดิมและปรับกราฟทั้งชุด แผนที่จะขยับแต่หมุดวัตถุไม่ขยับตาม
    แนวทางแก้ไขคือเก็บผลการสังเกตเป็นคู่ของตำแหน่งหุ่นยนต์กับระยะและมุม
    แล้วคำนวณพิกัดใหม่ทุกครั้งที่เผยแพร่

    \item \textbf{วัตถุที่อยู่นอกระนาบการสแกนของไลดาร์จะถูกปฏิเสธ}
    ไลดาร์สแกนในระนาบแนวนอนที่ความสูงติดตั้งเพียงระนาบเดียว
    วัตถุที่วางอยู่บนโต๊ะหรือชั้นวางที่สูงกว่าระนาบนั้นจะไม่มีลำแสงตกกระทบ
    ระบบจึงปฏิเสธผลการตรวจจับนั้นแทนที่จะประมาณระยะจากขนาดของกรอบภาพ
    ซึ่งเป็นการแลกความสามารถในการตรวจพบกับการไม่วางวัตถุผิดตำแหน่ง

    \item \textbf{แบบจำลองความลึกเป็นแบบสองมิติ}
    ระบบละเลยความสูงของวัตถุ มุมก้มของกล้อง และความบิดเบี้ยวของเลนส์
    ซึ่งรวมกันแล้วมีผลไม่เกินระดับไม่กี่เซนติเมตรที่ระยะ 3 เมตร
    แต่เป็นข้อจำกัดที่ต้องระบุไว้เมื่อเทียบกับเกณฑ์ความคลาดเคลื่อน 50 เซนติเมตร

    \item \textbf{การยืนยันว่าเป็นวัตถุเดิมพึ่งพาหมายเลขติดตามของตัวตรวจจับเป็นหลัก}
    หากวัตถุชิ้นเดียวได้รับหมายเลขติดตามใหม่ เช่น เมื่อหลุดจากเฟรมแล้วกลับเข้ามาใหม่
    หรือเมื่อมองจากมุมที่ต่างออกไปมาก ระบบจะไม่มีทางทราบว่าเป็นวัตถุเดิม
    และจะบันทึกเป็นวัตถุสองชิ้นบนแผนที่
    การจับคู่ด้วยระยะทางที่ใช้เป็นทางเลือกสำรองช่วยได้เฉพาะเมื่อวัตถุอยู่ในรัศมี 0.5 เมตร
    และเป็นคลาสเดียวกันเท่านั้น
    ผลการวัดที่บันทึกไว้ในบทที่ 4 พบหมายเลขติดตาม 144 หมายเลขในการทดสอบ 301 วินาที
    โดยมีคลาสที่ได้รับมากกว่าหนึ่งหมายเลข ซึ่งข้อมูลที่มีไม่สามารถแยกได้ว่า
    เกิดจากวัตถุหลายชิ้นจริงหรือจากการแตกของหมายเลขติดตาม

    \item \textbf{การวัดตำแหน่งวัตถุจากสี่ทิศทางไม่ได้ล้างรายการวัตถุระหว่างทิศทาง}
    ตามที่ระบุไว้ในหัวข้อ \ref{sec:eval} ทั้งสี่ทิศทางใช้รายการวัตถุรายการเดียวกันต่อเนื่องกัน
    ตำแหน่งจึงถูกปรับด้วยค่าเฉลี่ยเคลื่อนที่แบบเลขชี้กำลัง ($\alpha = 0.3$ ต่อการผสานข้อมูลหนึ่งครั้ง)
    ซึ่งลืมค่าจากทิศทางก่อนหน้าภายในเวลาประมาณ 1--2 วินาที
    แต่ละรอบ 30 วินาทีจึงสะท้อนทิศทางนั้นเป็นหลัก
    ผลที่ว่าเก้าอี้คงเป็นรายการเดียวแสดงว่าการรวมรายการซ้ำภายในรัศมี 0.5 เมตรทำงาน
    แต่ไม่ได้แยกทดสอบว่าหากล้างรายการก่อนแต่ละทิศทาง ตำแหน่งเริ่มต้นจากแต่ละทิศทางจะต่างกันเท่าใด
    นอกจากนี้ความคลาดเคลื่อนในแนวด้านข้างที่พบในด้านหลังและด้านขวายังไม่ได้วิเคราะห์หาสาเหตุ

    \item \textbf{การแก้ทิศทางเมื่อล้อลื่นไถลอาศัยการจับคู่ข้อมูลสแกนเป็นหลัก
    และมีขอบเขตการค้นหาจำกัด} เมื่อหุ่นยนต์ติดขัดและล้อหมุนฟรี
    โอดอมิเตอร์จะรายงานการหมุนที่ไม่ได้เกิดขึ้นจริง
    ซึ่ง slam\_toolbox จะดูดซับความคลาดเคลื่อนนี้เข้าไปในการแปลงพิกัดจาก
    \texttt{map} ไปยัง \texttt{odom} ตามการออกแบบ
    แต่ขอบเขตการค้นหาเชิงมุมของการจับคู่ข้อมูลสแกนถูกตั้งไว้ที่ $\pm$10 องศา
    บนสมมติฐานว่าโอดอมิเตอร์เชื่อถือได้
    ที่ความเร็วเชิงมุมสูงสุดของคอนโทรลเลอร์ การลื่นไถลเต็มความเร็วเพียงประมาณ 0.35 วินาที
    ก็เกินขอบเขตดังกล่าวแล้ว
    นอกจากนี้ประตูควบคุมการผสานข้อมูลของชั้นระบุตำแหน่งวัตถุอ่านความเร็วเชิงมุม
    จากโอดอมิเตอร์ของล้อ ขณะล้อลื่นไถลจึงเห็นการหมุนที่ไม่มีอยู่จริงและปิดการผสานข้อมูลทั้งหมด
    \textbf{กล่าวคือระบบจะหยุดบันทึกวัตถุพอดีในจังหวะที่หุ่นยนต์อยู่นิ่งที่สุด}
    พฤติกรรมทั้งสองข้อนี้ได้จากการวิเคราะห์ค่ากำหนดของระบบ ยังไม่ได้วัดบนหุ่นยนต์จริง
\end{enumerate}


\section{ข้อเสนอแนะ}

\begin{enumerate}[label=\arabic*., leftmargin=1.25cm, itemsep=6pt, topsep=6pt]
    \item \textbf{ปรับปรุงการรวมตำแหน่งวัตถุจากค่าเฉลี่ยเคลื่อนที่เป็นตัวกรองคาลมาน}
    ระบบปัจจุบันใช้ค่าเฉลี่ยเคลื่อนที่แบบถ่วงน้ำหนักเลขชี้กำลังที่ค่าคงที่
    ทำให้การสังเกตวัตถุจากระยะไกลซึ่งมีความไม่แน่นอนสูงถูกให้น้ำหนักเท่ากับการสังเกตระยะใกล้
    การเปลี่ยนไปใช้ตัวกรองคาลมานที่กำหนดความไม่แน่นอนตามระยะทาง
    ร่วมกับการกลั่นกรองด้วยระยะทางมาฮาลาโนบิส จะแก้ข้อจำกัดนี้ได้โดยตรง

    \item \textbf{ขยายไปสู่การประมวลผลสามมิติ}
    การเปลี่ยนไปใช้กล้องวัดความลึกหรือไลดาร์สามมิติ
    จะทำให้ระบุความสูงของวัตถุได้ และรองรับพื้นที่ที่มีชั้นวางสินค้าหลายระดับ

    \item \textbf{รองรับการทำงานของหุ่นยนต์หลายตัว}
    สถาปัตยกรรมแบบโมดูลของระบบเอื้อต่อการขยายให้หุ่นยนต์หลายตัว
    ร่วมกันสำรวจพื้นที่เดียวกันและรวมแผนที่เข้าด้วยกัน
    ซึ่งจะลดเวลาการสำรวจพื้นที่ขนาดใหญ่ได้อย่างมีนัยสำคัญ
\end{enumerate}

% =====================================================================
%  ท้ายเล่ม — เลขหน้ากลับไปเป็น ฌ ญ ฎ ฏ นับต่อจากส่วนนำ
% =====================================================================
\resumefrontpage

% ------------------- เอกสารอ้างอิง -------------------
\begin{thebibliography}{99}
\addcontentsline{toc}{chapter}{เอกสารอ้างอิง}

\bibitem{ros2}
S. Macenski, T. Foote, B. Gerkey, C. Lalancette, and W. Woodall,
``Robot Operating System 2: Design, architecture, and uses in the wild,''
\textit{Science Robotics}, vol. 7, no. 66, 2022.

\bibitem{ros2control}
S. Chitta, E. Marder-Eppstein, W. Meeussen, V. Pradeep, A. Rodr\'iguez Tsouroukdissian,
J. Bohren, D. Coleman, B. Magyar, G. Raiola, M. L\"udtke, and E. Fern\'andez Perdomo,
``ros\_control: A generic and simple control framework for ROS,''
\textit{Journal of Open Source Software}, vol. 2, no. 20, p. 456, 2017.

\bibitem{probrobotics}
S. Thrun, W. Burgard, and D. Fox,
\textit{Probabilistic Robotics}. Cambridge, MA: MIT Press, 2005.

\bibitem{slamtoolbox}
S. Macenski and I. Jambrecic,
``SLAM Toolbox: SLAM for the dynamic world,''
\textit{Journal of Open Source Software}, vol. 6, no. 61, p. 2783, 2021.

\bibitem{nav2}
S. Macenski, F. Mart\'in, R. White, and J. Gin\'es Clavero,
``The Marathon 2: A navigation system,''
in \textit{Proc. IEEE/RSJ International Conference on Intelligent Robots
and Systems (IROS)}, 2020.

\bibitem{yolo26}
R. Sapkota, R. H. Cheppally, A. Sharda, and M. Karkee,
``YOLO26: Key architectural enhancements and performance benchmarking
for real-time object detection,''
arXiv preprint arXiv:2509.25164, 2025.

\bibitem{ultralytics}
Ultralytics,
``Ultralytics YOLO,''
[Online]. Available: \url{https://github.com/ultralytics/ultralytics}

\bibitem{coco}
T.-Y. Lin, M. Maire, S. Belongie, J. Hays, P. Perona, D. Ramanan, P. Doll\'ar,
and C. L. Zitnick,
``Microsoft COCO: Common objects in context,''
in \textit{Proc. European Conference on Computer Vision (ECCV)}, 2014, pp. 740--755.

\bibitem{tensorrt}
NVIDIA Corporation,
``NVIDIA TensorRT,''
[Online]. Available: \url{https://developer.nvidia.com/tensorrt}

\bibitem{onnxruntime}
ONNX Runtime developers,
``ONNX Runtime: Cross-platform accelerated machine learning,''
2021. [Online]. Available: \url{https://onnxruntime.ai}

\bibitem{rep105}
W. Meeussen,
``REP 105: Coordinate frames for mobile platforms,''
ROS Enhancement Proposal, 2010.

\bibitem{zhang2000}
Z. Zhang,
``A flexible new technique for camera calibration,''
\textit{IEEE Transactions on Pattern Analysis and Machine Intelligence},
vol. 22, no. 11, pp. 1330--1334, 2000.

\bibitem{fastapi}
S. Ram\'irez,
``FastAPI,''
2018. [Online]. Available: \url{https://fastapi.tiangolo.com}

\bibitem{react}
Meta Platforms, Inc.,
``React,''
[Online]. Available: \url{https://react.dev}

\bibitem{robotlocalization}
T. Moore and D. Stouch,
``A generalized extended Kalman filter implementation for the Robot Operating System,''
in \textit{Proc. 13th International Conference on Intelligent Autonomous Systems (IAS-13)},
Springer, 2014.

\bibitem{yamauchi1997}
B. Yamauchi,
``A frontier-based approach for autonomous exploration,''
in \textit{Proc. IEEE International Symposium on Computational Intelligence
in Robotics and Automation (CIRA)}, 1997, pp. 146--151.

\bibitem{mexplore}
robo-friends,
``m-explore-ros2: ROS 2 port of the m-explore multi-robot exploration packages,''
[Online]. Available: \url{https://github.com/robo-friends/m-explore-ros2}

\bibitem{bytetrack}
Y. Zhang, P. Sun, Y. Jiang, D. Yu, F. Weng, Z. Yuan, P. Luo, W. Liu, and X. Wang,
``ByteTrack: Multi-object tracking by associating every detection box,''
in \textit{Proc. European Conference on Computer Vision (ECCV)}, 2022, pp. 1--21.

\end{thebibliography}

% ------------------- ภาคผนวก -------------------
\appendix
\renewcommand{\thechapter}{\thaialph{chapter}}

\addtocontents{toc}{\protect\vspace{6pt}}
\addtocontents{toc}{\protect\noindent\textbf{ภาคผนวก}\protect\par\protect\vspace{2pt}}
\addtocontents{toc}{\protect\renewcommand{\protect\cftchappresnum}{ภาคผนวก\protect\ }}
\addtocontents{toc}{\protect\setlength{\protect\cftchapindent}{1cm}}
\addtocontents{toc}{\protect\setlength{\protect\cftchapnumwidth}{3.4cm}}

\chapter{คู่มือการติดตั้งและใช้งานระบบ}

ภาคผนวกนี้รวบรวมขั้นตอนการติดตั้งและเปิดใช้งานระบบตั้งแต่ต้น
ตั้งแต่ข้อกำหนดของฮาร์ดแวร์ การติดตั้งซอฟต์แวร์บนหน่วยประมวลผลหลัก
การอัปโหลดเฟิร์มแวร์ ไปจนถึงลำดับคำสั่งเปิดใช้งานระบบและการแก้ปัญหาที่พบบ่อย

\section{ข้อกำหนดของฮาร์ดแวร์และการต่อสายสัญญาณ}

การต่อขาสัญญาณระหว่าง ESP32 กับอุปกรณ์ต่อพ่วงแสดงดังตารางที่ \ref{tab:pinout}
ขาที่ต้องระวังเป็นพิเศษมีสองกลุ่ม กลุ่มแรกคือ GPIO12
ซึ่งเป็นขากำหนดสถานะขณะบูตตามที่อธิบายไว้ในบทที่ 5
และกลุ่มที่สองคือ GPIO34 และ GPIO35 ซึ่งเป็นขารับสัญญาณเข้าอย่างเดียว
ไม่มีตัวต้านทานดึงขึ้นภายใน จึงต้องให้เอนโคเดอร์ขับสัญญาณออกมาเอง

\begin{table}[htbp]
\centering
\caption[การกำหนดขาสัญญาณของไมโครคอนโทรลเลอร์ (ภาคผนวก ก)]{การกำหนดขาสัญญาณของไมโครคอนโทรลเลอร์}
\label{tab:pinout}
\small
\begin{tabularx}{\textwidth}{|X|c|X|}
\hline
\textbf{หน้าที่} & \textbf{ขา} & \textbf{หมายเหตุ} \\
\hline
มอเตอร์ซ้าย เดินหน้า / ถอยหลัง & GPIO25 / GPIO26 & ป้อนสัญญาณควบคุมความเร็ว \\
\hline
มอเตอร์ซ้าย เปิดใช้งานช่องสัญญาณ & GPIO27 & ตรึงไว้ที่แรงดันสูง \\
\hline
มอเตอร์ขวา เดินหน้า / ถอยหลัง & GPIO14 / GPIO12 & GPIO12 เป็นขากำหนดสถานะขณะบูต \\
\hline
มอเตอร์ขวา เปิดใช้งานช่องสัญญาณ & GPIO13 & ตรึงไว้ที่แรงดันสูง \\
\hline
เอนโคเดอร์ซ้าย ช่อง A / B & GPIO34 / GPIO35 & ขารับสัญญาณเข้าอย่างเดียว \\
\hline
เอนโคเดอร์ขวา ช่อง A / B & GPIO23 / GPIO22 & \\
\hline
หน่วยวัดความเฉื่อย SDA / SCL & GPIO21 / GPIO19 & I2C ที่ 400 กิโลเฮิรตซ์ ที่อยู่ 0x68 \\
\hline
\end{tabularx}
\end{table}

\section{การติดตั้งระบบปฏิบัติการและ ROS 2}

หน่วยประมวลผลหลักติดตั้งชุดซอฟต์แวร์ JetPack 6.1 ซึ่งใช้ฐาน Ubuntu 22.04
จากนั้นติดตั้ง ROS 2 Humble จากที่เก็บแพ็กเกจของ ROS ดังนี้

{\small
\begin{verbatim}
sudo apt install ros-humble-desktop ros-dev-tools
sudo apt install ros-humble-ros2-control ros-humble-ros2-controllers \
                 ros-humble-slam-toolbox ros-humble-robot-localization \
                 ros-humble-image-tools ros-humble-twist-mux \
                 ros-humble-navigation2 ros-humble-nav2-bringup
sudo apt install build-essential cmake git pkg-config python3-venv
\end{verbatim}
}

ไดรเวอร์ของไลดาร์ไม่มีในที่เก็บแพ็กเกจ ต้องคอมไพล์จากซอร์สโค้ดของผู้ผลิต
(YDLidar-SDK และ \texttt{ydlidar\_ros2\_driver})
ในทำนองเดียวกัน แพ็กเกจ explore\_lite สำหรับการสำรวจอัตโนมัติไม่มีในที่เก็บแพ็กเกจของ ROS 2 Humble
ต้องคอมไพล์จากซอร์สโค้ด m-explore-ros2 ด้วยคำสั่ง \texttt{make explore-deps}
ส่วนชุดซอฟต์แวร์สำหรับหน่วยประมวลผลกราฟิกต้องติดตั้งแยกและเลือกเฉพาะแพ็กเกจที่จำเป็น
ได้แก่ CUDA 12.6 cuDNN 9.3 และ TensorRT 10.3
\textbf{ข้อควรระวังที่ทำให้เสียเวลามากที่สุดคือ} ไลบรารีเชื่อมต่อภาษาไพทอนของ TensorRT
ต้องการแพ็กเกจ \texttt{nvidia-l4t-dla-compiler} ที่รุ่นตรงกับเคอร์เนล
และต้องสั่ง \texttt{sudo ldconfig} หลังติดตั้งเสมอ มิฉะนั้นการนำเข้าไลบรารีจะล้มเหลว

\section{การอัปโหลดเฟิร์มแวร์ลง ESP32}

{\small
\begin{verbatim}
cd esp-motor-firmware
arduino-cli compile . -b esp32:esp32:esp32doit-devkit-v1 -p /dev/esp32
arduino-cli upload  . -b esp32:esp32:esp32doit-devkit-v1 -p /dev/esp32
arduino-cli monitor . -p /dev/esp32 --config 57600
\end{verbatim}
}

เมื่อเปิดหน้าต่างเฝ้าดูพอร์ตอนุกรมที่อัตราบอด 57600
บอร์ดต้องแสดงข้อความยืนยันการบูตพร้อมสถานะของวงจรนับพัลส์
และเมื่อพิมพ์คำสั่ง \texttt{e} ต้องได้ค่าเอนโคเดอร์สองค่ากลับมา
พิมพ์ \texttt{r} ต้องได้คำตอบยืนยันการรีเซ็ต
ข้อความแจ้งเตือนเรื่องขาที่ไม่มีตัวต้านทานดึงขึ้นจำนวนสี่บรรทัดเป็นพฤติกรรมปกติ
ตามที่อธิบายไว้ในบทที่ 5

\section{การกำหนดกฎ udev สำหรับชื่ออุปกรณ์ถาวร}

ติดตั้งกฎด้วยสคริปต์ \texttt{make udev} ซึ่งจะเขียนไฟล์
\texttt{/etc/udev/rules.d/99-my-bot-serial.rules}
โดยกำหนดสิทธิ์เป็นกลุ่ม \texttt{dialout} โหมด 0660
ตรวจสอบผลด้วยคำสั่ง \texttt{ls -l /dev/esp32 /dev/ydlidar}
\textbf{เนื่องจากกฎอ้างอิงจากตำแหน่งพอร์ตยูเอสบี
การย้ายสายไปเสียบพอร์ตอื่นจะทำให้ชื่ออุปกรณ์ผิดโดยไม่มีข้อความแจ้งเตือนใด ๆ}
จึงควรทำเครื่องหมายที่พอร์ตทั้งสองไว้ และสั่ง \texttt{make udev} ใหม่ทุกครั้งที่ย้ายพอร์ต

\section{การสร้างสภาพแวดล้อมเสมือนสำหรับการอนุมานแบบจำลอง}

{\small
\begin{verbatim}
uv venv --system-site-packages --python /usr/bin/python3.10 ~/yolo/venv
# 1. ติดตั้ง torch / torchvision จาก wheel ของ JetPack โดยระบุ URL โดยตรง
# 2. ตรึง numpy < 2 หลังจากติดตั้ง torch เท่านั้น
# 3. ติดตั้ง ultralytics โดยไม่ให้ดึงไลบรารีที่ต้องพึ่งพาตามมา
# 4. ติดตั้ง onnxruntime-gpu จาก wheel ของ JetPack และ lap สำหรับตัวติดตามวัตถุ
# (TensorRT ไม่ได้ติดตั้งผ่าน pip แต่ใช้ของ JetPack ผ่าน --system-site-packages)
\end{verbatim}
}

ลำดับของทั้งสี่ขั้นตอนมีความสำคัญและไม่สามารถแทนด้วยไฟล์รายการแพ็กเกจที่ตรึงรุ่นไว้ได้
เพราะไฟล์ดังกล่าวบันทึกได้เพียงรุ่น แต่บันทึกลำดับและข้อยกเว้นไม่ได้
หลังติดตั้งเสร็จต้องตรวจสอบสามอย่างเสมอ
คือ \texttt{torch.cuda.is\_available()} ต้องเป็นจริง
\texttt{import tensorrt} ต้องทำงานได้ภายในสภาพแวดล้อมเสมือน
และรายชื่อผู้ให้บริการประมวลผลของ ONNX Runtime ต้องมี CUDA อยู่จริง (สำหรับเส้นทางสำรอง)
มิฉะนั้นระบบจะทำงานบนหน่วยประมวลผลกลางหรือสร้างเครื่องยนต์ไม่ได้โดยไม่มีข้อความแจ้งเตือนที่ชัดเจน

\section{ลำดับคำสั่งเปิดใช้งานระบบ}

แต่ละคำสั่งเปิดในหน้าต่างเทอร์มินัลของตนเอง และต้องเปิดตามลำดับนี้
\textbf{ยกเว้นการเปิดครั้งแรกบนอุปกรณ์ใหม่หรือหลังเปลี่ยนรุ่น JetPack}
ซึ่งยังไม่มีเครื่องยนต์ TensorRT ให้สั่ง \texttt{make yolo} ก่อนคำสั่งอื่นทั้งหมด
เพื่อให้สร้างเครื่องยนต์ได้ (ประมาณ 13 นาที) เพราะสคริปต์จะไม่ยอมสร้างหากหน่วยความจำว่างไม่พอ
เมื่อสร้างเสร็จแล้วจึงปิดและเปิดระบบตามลำดับด้านล่าง

{\small
\begin{verbatim}
make real       # ros2_control, ไดรเวอร์ไลดาร์, กล้อง, twist_mux
make slam       # slam_toolbox โหมดออนไลน์อะซิงโครนัส
make nav        # Nav2
make explore    # (ไม่บังคับ) สำรวจพื้นที่อัตโนมัติ ต้องเปิดสามคำสั่งแรกก่อน
make yolo       # เปิดโหนดตรวจจับวัตถุ YOLO26l (TensorRT)
make semantic   # โหนดผสานข้อมูลและระบุตำแหน่งวัตถุ
make bridge     # ส่วนหลังของเว็บแอปพลิเคชัน พอร์ต 8000
make ui         # ส่วนหน้าของเว็บแอปพลิเคชัน พอร์ต 3000
\end{verbatim}
}

การตรวจสอบว่าระบบทำงานถูกต้องใช้การวัดอัตราการเผยแพร่ของหัวข้อหลักสี่หัวข้อ
ได้แก่ ข้อมูลการสแกนประมาณ 11.6 ครั้งต่อวินาที
สถานะข้อต่อ 30 ครั้งต่อวินาที ผลการตรวจจับประมาณ 15 ครั้งต่อวินาที
และการแปลงพิกัดจาก \texttt{map} ไปยัง \texttt{odom} ประมาณ 12--50 ครั้งต่อวินาที
(ต่ำลงเมื่อเปิดครบทุกส่วน ตามที่อธิบายในบทที่ 4)
\textbf{ข้อควรระวังคือการสมัครรับข้อมูลการสแกนต้องใช้คุณภาพการส่งแบบไม่รับประกันการส่งถึง}
มิฉะนั้นจะไม่ได้รับข้อความใดเลยและมีอาการเหมือนหัวข้อที่ไม่มีผู้เผยแพร่

\section{การเข้าใช้งานเว็บแอปพลิเคชันและการแก้ปัญหาที่พบบ่อย}

เปิดเบราว์เซอร์ไปที่พอร์ต 3000 ของหมายเลขไอพีของหุ่นยนต์
โดยไฟล์ตั้งค่าของส่วนหน้าต้องชี้ไปที่หมายเลขไอพีของหุ่นยนต์ ไม่ใช่ \texttt{localhost}
เพราะเบราว์เซอร์ทำงานอยู่บนเครื่องอื่น
ปัญหาที่พบบ่อยและวิธีแก้สรุปดังตารางที่ \ref{tab:troubleshoot}

\begin{table}[htbp]
\centering
\caption[อาการที่พบบ่อยและสาเหตุ (ภาคผนวก ก)]{อาการที่พบบ่อยและสาเหตุ}
\label{tab:troubleshoot}
\small
\begin{tabularx}{\textwidth}{|>{\raggedright\arraybackslash}p{5cm}|X|}
\hline
\textbf{อาการ} & \textbf{สาเหตุและวิธีแก้} \\
\hline
หุ่นยนต์ไม่ขยับ แต่ค่าเอนโคเดอร์และการสื่อสารปกติ
 & ขั้วต่อแบตเตอรี่ของภาคขับมอเตอร์หลวม ไมโครคอนโทรลเลอร์รับไฟจากยูเอสบีจึงทำงานปกติ \\
\hline
แผนที่เบลอหรือผนังหนาผิดปกติ
 & ขับหุ่นยนต์เร็วเกินไป ให้ลดความเร็วการขับด้วยมือให้เท่ากับความเร็วของระบบนำทาง \\
\hline
ไม่ได้รับข้อมูลการสแกนเลย แต่ไม่มีข้อความแจ้งข้อผิดพลาด
 & สมัครรับข้อมูลด้วยคุณภาพการส่งแบบรับประกันการส่งถึง \\
\hline
อัตราเฟรมของการตรวจจับต่ำผิดปกติ
 & ไลบรารีเรียนรู้เชิงลึกเป็นรุ่นที่ไม่รองรับหน่วยประมวลผลกราฟิก \\
\hline
ระบบนำทางเริ่มเคลื่อนที่แล้วหยุดค้าง
 & เป้าหมายถูกประกาศในเฟรม \texttt{odom} ให้ตั้งเฟรมอ้างอิงของเครื่องมือแสดงผลเป็น \texttt{map} \\
\hline
วัตถุที่เห็นชัดไม่ถูกบันทึกลงแผนที่
 & วัตถุอยู่ใกล้เกินไปจนช่วงมุมคร่อมฉากหลัง ให้ถอยหุ่นยนต์ออกมา \\
\hline
\end{tabularx}
\end{table}

แผนผังการต่อสายสัญญาณสำหรับขั้นตอนที่ 1 ได้แสดงไว้แล้วในเนื้อหาหลัก โดยภาพรวมการเชื่อมต่ออุปกรณ์ทั้งระบบแสดงดังรูปที่ \ref{fig:hwblock} และรายละเอียดการต่อสัญญาณระหว่างไมโครคอนโทรลเลอร์ โมดูลขับมอเตอร์ และเอนโคเดอร์ แสดงดังรูปที่ \ref{fig:wiring-motor} ให้ใช้สองรูปดังกล่าวประกอบการต่อสาย

ลักษณะของระบบเมื่อเปิดใช้งานครบทุกโหนดตามขั้นตอนที่ 6 แสดงดังรูปที่ \ref{fig:app-launch}

\begin{figure}[htbp]
\centering
\includegraphics[width=0.75\textwidth]{images/all_nodes_running.png}
\caption[สถานะของระบบเมื่อเปิดใช้งานครบทุกโหนด (บทที่ ก)]{สถานะของระบบเมื่อเปิดใช้งานครบทุกโหนด
ผลของคำสั่ง \texttt{ros2 node list} ขณะเปิดทั้งระบบ
ครอบคลุมฐานหุ่นยนต์ ไลดาร์ SLAM ระบบนำทาง การตรวจจับวัตถุ ชั้นผสานข้อมูล และเว็บแอปพลิเคชัน
(\texttt{transform\_listener\_impl\_*} คือโหนดช่วยที่ไลบรารี tf2 สร้างขึ้นเองภายในโหนดที่อ่านการแปลงพิกัด)}
\label{fig:app-launch}
\end{figure}

\chapter{ซอร์สโค้ดส่วนสำคัญของระบบ}

ภาคผนวกนี้แสดงตำแหน่งของซอร์สโค้ดแต่ละส่วนในที่เก็บซอร์สโค้ดของโครงงาน
พร้อมกับตัวอย่างโค้ดของสองส่วนที่เป็นแกนหลักของระบบ
รายการไฟล์ทั้งหมดแสดงดังตารางที่ \ref{tab:srcmap}

\begin{table}[htbp]
\centering
\caption[ตำแหน่งของซอร์สโค้ดส่วนสำคัญของระบบ (ภาคผนวก ข)]{ตำแหน่งของซอร์สโค้ดส่วนสำคัญของระบบ}
\label{tab:srcmap}
\small
\begin{tabularx}{\textwidth}{|X|X|}
\hline
\textbf{ส่วนของระบบ} & \textbf{ไฟล์} \\
\hline
เฟิร์มแวร์ควบคุมมอเตอร์และอ่านค่าเอนโคเดอร์
 & \texttt{esp-motor-firmware/} \newline
   \texttt{esp-motor-firmware.ino}, \texttt{pid.cpp}, \texttt{encoders.cpp},
   \texttt{motors.cpp}, \texttt{mpu6050.cpp}, \texttt{config.h} \\
\hline
Hardware Interface ของ ros2\_control
 & \texttt{my\_bot/hardware/diffdrive\_serial.cpp} \newline
   \texttt{my\_bot/hardware/serial\_port.cpp} \\
\hline
ไฟล์กำหนดค่าของคอนโทรลเลอร์และตัวกรองคาลมาน
 & \texttt{my\_bot/config/my\_controllers.yaml} \newline
   \texttt{my\_bot/config/ekf.yaml} \\
\hline
ไฟล์ URDF ที่บรรยายโครงสร้างและตำแหน่งเซนเซอร์
 & \texttt{my\_bot/description/robot.urdf.xacro} และไฟล์ย่อย
   \texttt{robot\_core}, \texttt{lidar}, \texttt{camera}, \texttt{imu},
   \texttt{ros2\_control} \\
\hline
ไฟล์กำหนดค่าของ slam\_toolbox และ Nav2
 & \texttt{my\_bot/config/} \newline
   \texttt{mapper\_params\_online\_async.yaml}, \texttt{nav2\_params.yaml} \\
\hline
โหนดตรวจจับวัตถุและการส่งออกแบบจำลอง
 & \texttt{my\_bot/scripts/yolo\_detector.py} \newline
   \texttt{yolo/export\_engine.py}, \texttt{yolo/export\_onnx.py} \\
\hline
โมดูลแปลงพิกัดวัตถุลงบนแผนที่
 & \texttt{semantic\_objects/} \newline
   \texttt{lidar\_range\_extractor.py}, \texttt{world\_point\_projector.py},
   \texttt{landmark\_store.py}, \texttt{semantic\_objects\_node.py} \\
\hline
ส่วนหลังและส่วนหน้าของเว็บแอปพลิเคชัน
 & \texttt{semantic\_bridge/} (FastAPI) \newline
   \texttt{semantic\_map\_ui/} (React) \\
\hline
\end{tabularx}
\end{table}

\section{รอบการควบคุมแบบ PID บนไมโครคอนโทรลเลอร์}

ฟังก์ชันต่อไปนี้ทำงานที่อัตรา 30 ครั้งต่อวินาที
โดยรับค่าเป้าหมายเป็นจำนวนพัลส์ต่อหนึ่งรอบการควบคุม
จุดที่ควรสังเกตคือการหยุดสะสมค่าอินทิกรัลเมื่อสัญญาณควบคุมถึงขีดจำกัดแล้ว
เพื่อป้องกันการสะสมค่าเกิน (Integral Wind-up)

{\small
\begin{verbatim}
static void doPID(SetPointInfo *p) {
  int  input  = p->Encoder - p->PrevEnc;
  long Perror = p->TargetTicksPerFrame - input;

  long output = (Kp * Perror - Kd * (input - p->PrevInput) + p->ITerm) / Ko;
  p->PrevEnc = p->Encoder;
  output += p->output;

  if (output >= MAX_PWM) {
    output = MAX_PWM;
  } else if (output <= -MAX_PWM) {
    output = -MAX_PWM;
  } else {
    p->ITerm += Ki * Perror;   /* สะสมค่าอินทิกรัลเฉพาะเมื่อยังไม่ถึงขีดจำกัด */
  }

  p->output    = output;
  p->PrevInput = input;
}
\end{verbatim}
}

\section{การหาตำแหน่งวัตถุจากจุดตัดของรังสีสายตากับวงกลมระยะทาง}

โค้ดต่อไปนี้คือการคำนวณตามสมการที่อธิบายไว้ในบทที่ 4
โดย \texttt{C} คือจุดกำเนิดของกล้อง \texttt{L} คือจุดกำเนิดของไลดาร์
และ \texttt{r} คือระยะทางที่ได้จากลำแสงไลดาร์
(ตัดส่วนการสร้างค่าที่ส่งกลับในกรณีปฏิเสธออกเพื่อความกระชับ)

{\small
\begin{verbatim}
az_base_ros = -estimate.azimuth_center + self.extrinsics.yaw
dx_ray = math.cos(az_base_ros)
dy_ray = math.sin(az_base_ros)

cx, cy = self.extrinsics.dx, self.extrinsics.dy      # จุดกำเนิดของกล้อง
lx, ly = self.extrinsics.range_origin                # จุดกำเนิดของไลดาร์
mx, my = cx - lx, cy - ly

m_dot_d = mx * dx_ray + my * dy_ray
disc    = m_dot_d * m_dot_d - (mx * mx + my * my) + r * r
if disc < 0.0:
    return WorldPoint(valid=False)                   # ปฏิเสธการตรวจจับนี้

t = -m_dot_d + math.sqrt(disc)
x_in_base = cx + t * dx_ray
y_in_base = cy + t * dy_ray
\end{verbatim}
}

\end{document}
